\documentclass[12pt,a4paper]{article}

\usepackage[utf8]{inputenc}
\usepackage[T1]{fontenc}
\usepackage{lmodern}
\usepackage{amsmath,amssymb,amsthm}
\usepackage{enumitem}
\usepackage{xcolor}
\usepackage{geometry}
\usepackage{microtype}
\usepackage{tabularx}
\usepackage{array}
\geometry{a4paper, margin=2.5cm}
\usepackage{hyperref}

\hypersetup{
    colorlinks=true,
    linkcolor=blue!70!black,
    citecolor=green!50!black,
    urlcolor=blue!70!black,
    pdftitle={A Thermodynamic Obstruction to Finite-Time Blow-Up in the 3D Navier-Stokes Equations},
    pdfauthor={Lukas Geiger},
    pdfsubject={Bedingtes Navier-Stokes-Regularitätsrahmenwerk},
    pdfkeywords={Navier-Stokes, Blow-up, Regularität, Attraktor, Dissipation}
}
\setlength{\emergencystretch}{3em}
\newcolumntype{Y}{>{\raggedright\arraybackslash}X}

\theoremstyle{definition}
\newtheorem{definition}{Definition}[section]
\theoremstyle{plain}
\newtheorem{theorem}[definition]{Satz}
\newtheorem{lemma}[definition]{Lemma}
\newtheorem{proposition}[definition]{Proposition}
\newtheorem{korollar}[definition]{Korollar}
\newtheorem{corollary}[definition]{Korollar}
\newtheorem{axiom}[definition]{Axiom}
\theoremstyle{remark}
\newtheorem{remark}[definition]{Bemerkung}

\renewcommand{\proofname}{Beweis}

\title{\textbf{A Thermodynamic Obstruction to Finite-Time Blow-Up\\ in the 3D Navier-Stokes Equations (Conditional Framework)}}
\author{Lukas Geiger\\
\small Unabhängiger Forscher, Bernau, Deutschland\\
\small ORCID: \href{https://orcid.org/0009-0005-7296-1534}{0009-0005-7296-1534}%
\thanks{KI-Werkzeuge (Claude von Anthropic) wurden f\"ur mathematische Formalisierung, Literaturverifikation und redaktionelle Verfeinerung eingesetzt. Alle mathematischen Inhalte wurden vom Autor entwickelt, verifiziert und freigegeben.}}
\date{Mai 2026}

\begin{document}
\maketitle

\begin{abstract}
Wir schlagen ein bedingtes Regularitätskriterium für die inkompressiblen 3D-Navier-Stokes-Gleichungen vor, das auf einem thermodynamischen Abstandsfunktional basiert. Anstatt den Geschwindigkeitsgradienten global zu beschränken, konstruieren wir ein Attraktor-Abstandsfunktional $H(u \mid \mathcal{A})$, das den $L^2$-Abstand der aktiven Strömung von einem Referenzattraktor mit explizitem Regularitätspaket misst. Für die klassische 3D-Leray--Hopf-Theorie ist dieses Paket eine zusätzliche stehende Hypothese, kein etabliertes Theorem. Unter diesem Attraktor-Paket und zwei weiteren Bedingungen -- (\textbf{Assumption~G}) dass die zeitabhängige minimierende Projektion auf den Attraktor hinreichend regulär ist, und (\textbf{Condition~D}) dass die kumulative viskose Dissipation die Gronwall-Kosten nahe einer hypothetischen Blow-up-Zeit dominiert -- zeigen wir, dass ein Blow-up in endlicher Zeit $H$ unter Null zwingen würde, ein Widerspruch zur Nichtnegativität. Das Argument reduziert die Regularitätsfrage auf strukturelle Inputs: Attraktor-Regularität, Geometrie der Attraktorprojektion (Assumption~G) und das Verhältnis von Dissipation zu nichtlinearem Wachstum (Condition~D). Wir identifizieren hinreichende geometrische Bedingungen für~G (positiver Reach, log-Lipschitz-Projektionen, BV-Selektionen) und präsentieren einen $H^1$-Level-Lift, bei dem die Enstrophie-Dissipation -- die im Gegensatz zur Energie-Dissipation keine a-priori-Oberschranke unter Blow-up besitzt -- als divergenter Term dienen kann, wenn zusätzlich Bihari-Dominanz gilt. Wir diskutieren den Status dieser Annahmen, ihre Verbindung zur Theorie inertialer Mannigfaltigkeiten und die Lücken, die das vorliegende Rahmenwerk von einem unbedingten Beweis trennen.
Diese Arbeit etabliert eine \emph{bedingte Reduktion}: Globale Regularität folgt aus dem Attraktor-Regularitätspaket, Assumption~G und Condition~D; sie ist nicht allein äquivalent zur Existenz einer BV-regulären Selektion.
\end{abstract}

\clearpage
\tableofcontents
\clearpage

\section{Einleitung}
Die klassische Herausforderung bezüglich der inkompressiblen 3D-Navier-Stokes-Gleichungen besteht darin, die Entstehung von Singularitäten in endlicher Zeit (Blow-ups) im Geschwindigkeitsfeld $u(x,t) \in \mathbb{R}^3$ auszuschließen -- ein Problem, das seit Lerays grundlegender Arbeit \cite{Leray1934} offen geblieben ist. Die Strömung wird beschrieben durch:
\begin{equation}
    \partial_t u + (u \cdot \nabla) u = -\nabla p + \nu \Delta u, \quad \nabla \cdot u = 0, \quad \text{auf } \Omega \subset \mathbb{R}^3.
\end{equation}
Das bekannte Beale--Kato--Majda-Kriterium (BKM) \cite{BealeKatoMajda1984} besagt, dass eine Singularität zur Zeit $T^*$ genau dann auftritt, wenn die Vortizität divergiert:
\begin{equation}
    \int_0^{T^*} \|\nabla \times u(t)\|_{L^\infty} \, dt = \infty.
\end{equation}
In dieser Arbeit konstruieren wir ein thermodynamisches Ljapunow-Funktional, das den Abstand der aktiven Strömung von einem Referenzattraktor misst. Wir zeigen, dass -- \emph{bedingt durch das Attraktor-Regularitätspaket unten}, eine Regularitätsannahme über die Attraktorprojektion (Assumption~G, Abschnitt~3) und eine \emph{Dissipations-Dominanz-Bedingung}~(D) -- jeder hypothetische Blow-up zu einem Widerspruch führt: Die Enstrophie divergiert punktweise und treibt das nichtnegative Abstandsfunktional unter Null, sobald die kumulative viskose Dissipation die Gronwall-Kosten übersteigt. \textbf{Entscheidend ist, dass Condition~D auf $L^2$-Ebene eine strukturelle Limitation darstellt:} Die Leray-Energieidentität beschränkt das Dissipationsintegral von oben, während die Gronwall-Kosten exponentiell wachsen können, was Condition~D zunehmend schwierig bei langsamem Blow-up macht. Das Argument identifiziert daher die $H^1$-Ebene als die natürliche Sobolev-Ebene für die Obstruktionsstrategie, wo die Enstrophie-Dissipation keine a-priori-Oberschranke besitzt. Das resultierende bedingte Regularitätskriterium reduziert das vollständige Regularitätsproblem auf mehrere strukturelle Fragen: die Existenz eines ausreichend regulären Referenzattraktors, die Geometrie der Attraktorprojektion (Assumption~G) und die relativen Blow-up-Raten von Dissipation gegenüber Gronwall-Faktoren (Condition~D). Wir beanspruchen nicht, dass eine der Bedingungen leichter zu verifizieren ist als Regularität selbst; der Wert liegt in der Zerlegung und in der Identifikation konkreter Angriffsvektoren für jede Komponente.

Diese Arbeit ist in sich geschlossen; alle Hauptresultate hängen nur von Definitionen und Beweisen innerhalb dieses Dokuments ab. Verbindungen zum breiteren FST-Programm werden für den Kontext zitiert, sind aber logisch nicht erforderlich.

\section{Das Energie-Dissipations-Divergenzfunktional}
Wir definieren ein verallgemeinertes Divergenzfunktional $H(u \mid \mathcal{A})$, das den Abstand des aktiven Geschwindigkeitsfeldes $u$ von einem Referenzattraktor $\mathcal{A}$ mit dem folgenden Regularitätspaket darstellt.

\begin{definition}[Divergenz vom globalen Attraktor]
Sei $\Omega = \mathbb{T}^3$ der flache Torus. Wir arbeiten mit einer Referenzmenge $\mathcal{A}$ für die Navier--Stokes-Dynamik, die Axiom~\ref{prop:attractor-regularity} erfüllt. In Settings, in denen ein glatter kompakter globaler Attraktor bereits bekannt ist, kann $\mathcal{A}$ dieser Attraktor sein. Für das klassische 3D-Leray--Hopf-Problem ist dies ein bedingtes Referenzobjekt, kein etablierter globaler Leray--Hopf-Attraktor. Für jedes aktive Strömungsfeld $u \in L^2(\Omega)$ ist das Abstandsfunktional:
\begin{equation}
    H(u \mid \mathcal{A}) := \inf_{v \in \mathcal{A}} E_{v}(u) = \inf_{v \in \mathcal{A}} \frac{1}{2} \int_\Omega |u - v|^2 \, dx.
\end{equation}
Da $\mathcal{A}$ eine kompakte Teilmenge von $L^2(\Omega)$ ist, wird das Infimum angenommen. Wir bezeichnen mit $v(t) \in \mathcal{A}$ eine messbare Auswahl, die diesen Abstand zur Zeit $t$ minimiert, sodass $H(u \mid \mathcal{A}) = \frac{1}{2} \int_\Omega |w|^2 \, dx$ mit $w := u - v$.
\end{definition}

\begin{axiom}[Attraktor-Regularitätspaket]
\label{prop:attractor-regularity}
Sei $\Omega = \mathbb{T}^3$ und sei $f \in L^\infty_t H^1_x(\Omega)$ eine glatte
Antriebskraft. Der Referenzattraktor $\mathcal{A}$ wird mit folgenden Eigenschaften angenommen:
\begin{enumerate}
    \item[\emph{(i)}]   $\mathcal{A}$ ist kompakt in $L^2(\Omega)$ und in $H^1(\Omega)$.
    \item[\emph{(ii)}]  Jedes $v \in \mathcal{A}$ besitzt die unten benötigte Sobolev-Regularität; insbesondere $v \in H^2(\Omega) \cap W^{1,\infty}(\Omega)$.
    \item[\emph{(iii)}] $\sup_{v \in \mathcal{A}} \|v\|_{H^2} + \|v\|_{W^{1,\infty}} < \infty$.
\end{enumerate}
\end{axiom}

\begin{remark}[Zum Attraktor ohne unbedingten Regularitätsbeweis]
\label{rem:circularity}
Das Paket in Axiom~\ref{prop:attractor-regularity} ist automatisch in
Regimen, in denen globale starke Wohlgestelltheit und glatte kompakte
Attraktoren bereits verfügbar sind, etwa in zweidimensionalen,
gedämpften, hyperviskosen oder regularisierten Modellen. Für das
klassische dreidimensionale Leray--Hopf-Problem ist es \emph{kein}
bekanntes Theorem. Die Standardtheorie liefert globale schwache
Lösungen und schwache/Trajektorien-Attraktorobjekte in schwächeren
Topologien, aber keinen kompakten $H^1$-Semiflussattraktor mit
gleichmäßig glatten Elementen, Eindeutigkeit und den hier benötigten
Schranken. Jede Verwendung von Proposition/Axiom~\ref{prop:attractor-regularity}
unten ist daher eine bedingte Verwendung dieses Attraktor-Regularitätspakets.

\medskip\noindent
\textbf{Selektionsabhängigkeit.}\;
Da die Eindeutigkeit von Leray--Hopf-Lösungen offen ist, hängt der
Semifluss $S(t)$ von der Wahl der Lösungsselektion ab, und
verschiedene Selektionen können verschiedene Attraktoren erzeugen.
Das vorliegende Framework gilt für jede feste Selektion, die einen
Semifluss erzeugt, dessen $\omega$-Limesmenge
Axiom~\ref{prop:attractor-regularity} erfüllt. Die Schlussfolgerungen von
Satz~\ref{thm:main} gelten für jeden solchen Attraktor, sofern
die zugehörige Nächste-Punkt-Projektion Assumption~G erfüllt.
Dies ist konsistent mit dem Sell--Chepyzhov--Vishik-Trajektorienattraktor-Ansatz~\cite{ChepyzhovVishik2002}, der die Eindeutigkeitsfrage
vollständig umgeht, indem er mit Mengen vollständiger Trajektorien
statt mit einem einwertigen Semifluss arbeitet.
\end{remark}

Durch Verfolgen der Entwicklung dieses Funktionals leiten wir eine exakte Dissipationsschranke her.

\begin{lemma}[Energie-Dissipations-Divergenzschranke]
\label{lem:entropy-bound}
Sei $u \in L^\infty_t H^1_x(\Omega)$ eine divergenzfreie Leray--Hopf-Lösung der Navier-Stokes-Gleichungen mit glattem Antrieb $f$, und sei $v \in L^\infty_t H^1_x(\Omega)$ ein divergenzfreies Feld (nicht notwendigerweise eine NS-Lösung) mit lokal integrierbarem $\partial_t v$. Definiere das Residuum $\mathcal{R}_v := \partial_t v + (v\cdot\nabla)v - \nu\Delta v - f$, das verschwindet, wenn $v$ selbst eine NS-Lösung ist. Die Zeitentwicklung von $E_{v}(t) = \frac{1}{2}\| u - v \|_{L^2}^2$ ist beschränkt durch:
\begin{equation}
    \frac{d}{dt} E_{v}(t) \le - \frac{\nu}{2} \|\nabla u\|_{L^2}^2 + C_1(t) E_{v}(t) + C_2(t),
\end{equation}
wobei $C_1(t) = 2\|\nabla v\|_{L^\infty} + 1$ und $C_2(t) = \frac{1}{2}\|(v \cdot \nabla)v\|_{L^2}^2 + \frac{1}{2}\|f - \partial_t v + \mathcal{R}_v\|_{L^2}^2 + \frac{\nu}{2}\|\nabla v\|_{L^2}^2$.
Wenn $v$ eine NS-Lösung ist, gilt $\mathcal{R}_v = 0$ und $C_2$ vereinfacht sich zu $\frac{1}{2}\|(v \cdot \nabla)v\|_{L^2}^2 + \frac{1}{2}\|f - \partial_t v\|_{L^2}^2 + \frac{\nu}{2}\|\nabla v\|_{L^2}^2$. Wenn $v(t)$ der zeitvariable Attraktor-Minimierer ist (Bemerkung~\ref{rem:minimiser-trajectory}), muss der Term $\partial_t v$ als die gesamte Änderungsrate der Selektion interpretiert werden, und $\mathcal{R}_v \neq 0$ im Allgemeinen; die BV-Modifikation im Beweis von Satz~\ref{thm:main} (Gleichung~\eqref{eq:stieltjes-bound}) behandelt diesen Fall.
\end{lemma}

\begin{proof}
Wir subtrahieren die Gleichungen für $v$ von denen für $u$ und erhalten die Entwicklungsgleichung für $w = u - v$:
\begin{equation}
    \frac{d}{dt} E_{v}(t) = \int_\Omega w \cdot \partial_t w \, dx = \int_\Omega w \cdot \left( -(u \cdot \nabla)u - \nabla p + \nu \Delta u + f - \partial_t v \right) dx.
\end{equation}
Wegen $\nabla \cdot w = 0$ und der Periodizität von $\Omega$ verschwindet der Druckterm: $\int_\Omega w \cdot \nabla p \, dx = 0$.
Partielle Integration des viskosen Terms ergibt:
\begin{equation}
    \nu \int_\Omega w \cdot \Delta u \, dx = - \nu \int_\Omega \nabla w : \nabla u \, dx = - \nu \|\nabla u\|_{L^2}^2 + \nu \int_\Omega \nabla v : \nabla u \, dx.
\end{equation}
Anwenden der Cauchy-Schwarz- und Young-Ungleichung liefert:
\begin{equation}
    \nu \int_\Omega w \cdot \Delta u \, dx \le - \frac{\nu}{2} \|\nabla u\|_{L^2}^2 + \frac{\nu}{2} \|\nabla v\|_{L^2}^2.
\end{equation}
Für das kritische konvektive Integral $I := \int_\Omega w \cdot (u \cdot \nabla)u \, dx$ zerlegen wir $u = v + w$. Unter Nutzung der fundamentalen Aufhebungseigenschaft $\int_\Omega w \cdot (y \cdot \nabla)w \, dx = 0$ für jedes divergenzfreie $y$ isolieren wir die verbleibenden Terme:
\begin{equation}
    - I = - \int_\Omega w \cdot (v \cdot \nabla)v \, dx - \int_\Omega w \cdot (w \cdot \nabla)v \, dx.
\end{equation}
Diese Terme sind beschränkt durch:
\begin{equation}
    |-I| \le \|w\|_{L^2} \|(v \cdot \nabla)v\|_{L^2} + \|\nabla v\|_{L^\infty} \|w\|_{L^2}^2.
\end{equation}
Einsetzen dieser Schranken in die Zeitableitungsgleichung und Anwenden der Young-Ungleichung auf die restlichen linearen Terme $\|w\|_{L^2} \|f - \partial_t v\|_{L^2}$ und $\|w\|_{L^2} \|(v \cdot \nabla)v\|_{L^2}$ liefert das behauptete Resultat.
\end{proof}

Unter Axiom~\ref{prop:attractor-regularity} gelten die gleichmäßigen Schranken $\sup_{v \in \mathcal{A}} C_1(t) < \infty$ und $\sup_{v \in \mathcal{A}} C_2(t) < \infty$ für die im Abstandsfunktional verwendeten Referenzfelder. Infimumbildung über $\mathcal{A}$ überführt die individuelle Energieschranke in eine absolute Stabilitätsgrenze für die Strömungsdivergenz $H(u \mid \mathcal{A})$.

\begin{remark}[Der Minimierer ist keine einzelne NS-Trajektorie]
\label{rem:minimiser-trajectory}
Im Beweis von Satz~\ref{thm:main} ist die minimierende Auswahl
$v(t) \in \mathcal{A}$ das nächste Attraktorelement zur Zeit~$t$,
\emph{nicht} eine einzelne NS-Lösungstrajektorie. Während $u(t)$ evolviert,
kann der Minimierer zwischen verschiedenen Punkten von~$\mathcal{A}$ springen.
Lemma~\ref{lem:entropy-bound} ist für eine \emph{feste} NS-Lösung~$v$
formuliert; bei Anwendung auf den zeitvariablen Minimierer muss der Term
$\partial_t v$ in $C_2$ als die gesamte Änderungsrate der Auswahl
interpretiert werden (einschliesslich Sprünge), nicht lediglich als die
NS-Evolution bei~$v(t)$. Genau deshalb ist Assumption~G (zur Kontrolle
von $\partial_t v$) oder die BV-Modifikation (Stieltjes-Integral,
Gleichung~\eqref{eq:stieltjes-bound}) erforderlich.
\end{remark}

\section{Hauptresultat: Bedingter Ausschluss von Singularitäten in endlicher Zeit}

Wir formulieren nun das zentrale Regularitätsresultat. Zwei Hilfslemmas werden
vorab bewiesen, um das Hauptargument in sich geschlossen zu halten.

% ------------------------------------------------------------------
\subsection*{Lemma A (Biot--Savart logarithmische Abschätzung)}
% ------------------------------------------------------------------
\begin{lemma}[Biot--Savart-Abschätzung auf $\mathbb{T}^3$]
\label{lem:BS}
Sei $u$ ein periodisches, divergenzfreies Vektorfeld auf $\mathbb{T}^3$, und sei
$\omega = \nabla\times u$ seine Vortizität. Für jedes $s > 3/2$ existiert eine
Konstante $C_{\mathrm{BS}} = C_{\mathrm{BS}}(s,\Omega) > 0$ (die nur von
$s$ und der Geometrie von $\Omega = \mathbb{T}^3$ abhängt, \emph{nicht} aber von $u$), sodass
\begin{equation}\label{eq:BS}
    \|\nabla u\|_{L^\infty}
    \;\le\;
    C_{\mathrm{BS}}\,\|\omega\|_{L^\infty}
    \!\left(1 + \log^+\!\frac{\|\omega\|_{H^s}}{\|\omega\|_{L^\infty}}\right),
    \qquad \text{falls } \|\omega\|_{L^\infty} > 0.
\end{equation}
\end{lemma}

\begin{proof}
Wir rekonstruieren $\nabla u$ aus $\omega$ mittels des periodischen Biot--Savart-Gesetzes:
$u = K * \omega$, wobei $K$ der matrixwertige Faltungskern auf
$\mathbb{T}^3$ ist, der aus der Inversion des Rotor-Divergenz-Systems entsteht.
Insbesondere gilt $\nabla u = \nabla K * \omega$, und $\nabla K$ ist ein
Calder\'{o}n--Zygmund-Singulaerintegral-Kern (vgl.\ \cite{BertozziMajda2002},
Kapitel~2).

\medskip
\noindent\textbf{Schritt 1 (Littlewood--Paley-Zerlegung).}
Seien $\{\Delta_j\}_{j\ge 0}$ die Standard-Littlewood--Paley-Projektionsoperatoren
auf $\mathbb{T}^3$, assoziiert mit einer glatten dyadischen Zerlegung der Eins
$\{\varphi_j\}$ im Frequenzraum, sodass
$\operatorname{supp}\hat\varphi_j \subset \{|\xi| \sim 2^j\}$ und
$\sum_{j\ge 0}\varphi_j(\xi) = 1$ für alle $\xi \neq 0$.
Wir spalten $\nabla u$ bei einem freien ganzzahligen Parameter $N \ge 1$ auf:
\[
    \nabla u \;=\; \underbrace{\sum_{j=0}^{N} \Delta_j(\nabla u)}_{=:\,S_{\le N}}
                  \;+\; \underbrace{\sum_{j=N+1}^{\infty} \Delta_j(\nabla u)}_{=:\,S_{>N}}.
\]

\medskip
\noindent\textbf{Schritt 2 (Niedrige Frequenzen $j \le N$).}
Da $\nabla K$ ein Calder\'{o}n--Zygmund-Operator nullter Ordnung auf
$\mathbb{T}^3$ ist, bildet er $L^\infty$ nach $\mathrm{BMO}$ ab, und auf jedem Littlewood--Paley-Block $\Delta_j$ liefert die Bernstein-Einbettung
$\mathrm{BMO} \hookrightarrow L^\infty$ (bei fester Frequenzskala)
\[
    \|\Delta_j(\nabla u)\|_{L^\infty}
    \;\le\; C_{\mathrm{CZ}}\,\|\Delta_j \omega\|_{L^\infty}
    \;\le\; C_{\mathrm{CZ}}\,\|\omega\|_{L^\infty},
\]
wobei $C_{\mathrm{CZ}}$ nur von der Operatornorm von $\nabla K$ und
der Geometrie von $\mathbb{T}^3$ abhängt. Summation über $j = 0, \ldots, N$:
\begin{equation}\label{eq:BS-low}
    \|S_{\le N}\|_{L^\infty}
    \;\le\; C_{\mathrm{CZ}}\,(N+1)\,\|\omega\|_{L^\infty}.
\end{equation}

\medskip
\noindent\textbf{Schritt 3 (Hohe Frequenzen $j > N$).}
Für jeden dyadischen Block mit Fourier-Träger in $\{|\xi| \sim 2^j\}$ liefert die
Bernstein-Ungleichung auf $\mathbb{T}^3$ (mit $p=2$, $q=\infty$):
\[
    \|\Delta_j(\nabla u)\|_{L^\infty}
    \;\le\; C_B\,2^{3j/2}\,\|\Delta_j(\nabla u)\|_{L^2}.
\]
Da $u = K * \omega$ und $\widehat{\nabla K}(\xi)$ ein Symbol nullter Ordnung ist
(oben und unten beschränkt auf jeder Frequenzschale), gilt
$\|\Delta_j(\nabla u)\|_{L^2} \le C\,\|\Delta_j \omega\|_{L^2}$.
Ausdrücken von $\|\Delta_j \omega\|_{L^2}$ mittels der $H^s$-Norm
durch $\|\Delta_j\omega\|_{L^2} \le 2^{-js}\,\bigl(2^{js}\|\Delta_j\omega\|_{L^2}\bigr)
\le 2^{-js}\,\|\omega\|_{H^s}$ ergibt
\[
    \|\Delta_j(\nabla u)\|_{L^\infty}
    \;\le\; C'\,2^{j(3/2-s)}\,\|\omega\|_{H^s}.
\]
Da $s > 3/2$, ist der Exponent $\alpha := s - 3/2 > 0$ strikt positiv,
und die Restsumme ist eine konvergente geometrische Reihe:
\begin{equation}\label{eq:BS-high}
    \|S_{>N}\|_{L^\infty}
    \;\le\; \sum_{j>N} C'\,2^{-j\alpha}\,\|\omega\|_{H^s}
    \;=\; \frac{C'\,2^{-(N+1)\alpha}}{1 - 2^{-\alpha}}\,\|\omega\|_{H^s}
    \;\le\; C''\,2^{-N\alpha}\,\|\omega\|_{H^s}.
\end{equation}

\medskip
\noindent\textbf{Schritt 4 (Optimierung über $N$).}
Kombination von \eqref{eq:BS-low} und \eqref{eq:BS-high}:
\[
    \|\nabla u\|_{L^\infty}
    \;\le\; C_{\mathrm{CZ}}\,(N+1)\,\|\omega\|_{L^\infty}
          + C''\,2^{-N\alpha}\,\|\omega\|_{H^s}.
\]
Falls $\|\omega\|_{H^s}\le e\,\|\omega\|_{L^\infty}$, liefert bereits
$N=1$ die gewünschte Schranke nach Vergrößerung der Konstante. Andernfalls
wählen wir $N \in \mathbb{N}$ mit
\[
    2^{N\alpha} \;\simeq\; \frac{\|\omega\|_{H^s}}{\|\omega\|_{L^\infty}},
    \qquad
    N \;\simeq\; \frac{1}{\alpha\,\ln 2}\,
            \ln\!\frac{\|\omega\|_{H^s}}{\|\omega\|_{L^\infty}}.
\]
Der Hochfrequenzbeitrag ist dann durch ein konstantes Vielfaches von
$\|\omega\|_{L^\infty}$ beschränkt, und für den Niederfrequenzanteil gilt
\[
    C_{\mathrm{CZ}}\,(N+1)\,\|\omega\|_{L^\infty}
    \;\le\; \tilde C\,\|\omega\|_{L^\infty}\,
            \Bigl(1 + \log^+\!\frac{\|\omega\|_{H^s}}{\|\omega\|_{L^\infty}}\Bigr).
\]
Addition beider Abschätzungen und Absorption der universellen Konstanten in
$C_{\mathrm{BS}} = C_{\mathrm{BS}}(s,\Omega)$ ergibt die behauptete
Ungleichung \eqref{eq:BS}.
\end{proof}

\begin{remark}
Die Konstante $C_{\mathrm{BS}}$ kann explizit gemacht werden: Sorgfältiges Nachverfolgen der
Calder\'{o}n--Zygmund-Norm von $K$ und der Bernstein-Konstante ergibt
$C_{\mathrm{BS}} \lesssim C(s)(1 + \|K\|_{L^{3/2,\infty}})$. Insbesondere
ist $C_{\mathrm{BS}}$ \emph{gleichmäßig} über $\mathcal{A}$, da alle Attraktorelemente
denselben umgebenden Torus teilen.
\end{remark}

% ------------------------------------------------------------------
\subsection*{Lemma B (Enstrophie-Bilanz und punktweiser Blow-up)}
% ------------------------------------------------------------------
\begin{lemma}[Enstrophie-Bilanz und punktweiser Blow-up]
\label{lem:enstrophy}
Sei $u$ eine Leray--Hopf-Lösung der Navier--Stokes-Gleichungen mit glattem
Antrieb $f \in L^\infty_t H^1_x$ auf $\mathbb{T}^3$, und sei $u$ auf jedem
kompakten Teilintervall von $[0,T^*)$ eine starke Lösung. Die Enstrophie
erfüllt (im schwachen Sinne für f.u.\ $t$)
\begin{equation}\label{eq:enstrophy}
    \frac{1}{2}\frac{d}{dt}\|\omega(t)\|_{L^2}^2
    = -\nu\|\nabla\omega\|_{L^2}^2
      + \int_\Omega (\omega\cdot\nabla)u\cdot\omega\,dx
      + \int_\Omega (\nabla\times f)\cdot\omega\,dx.
\end{equation}
Wenn $T^*<\infty$ eine endliche maximale starke Existenzzeit ist, dann
divergiert die Enstrophie punktweise:
\begin{equation}\label{eq:diss-blowup}
    \limsup_{t \nearrow T^*}\|\nabla u(t)\|_{L^2}^2 = \infty.
\end{equation}
\end{lemma}

\begin{remark}[Was divergiert und was nicht bei Blow-up]
\label{rem:divergence-landscape}
Unter einer hypothetischen endlichen maximalen starken Existenzzeit
$T^* < \infty$:
\begin{itemize}
  \item \textbf{Divergiert (punktweise):}\;
    BKM liefert
    \[
      \limsup_{t\nearrow T^*}\|\omega(t)\|_{L^\infty} = \infty,
    \]
    und das Fortsetzungskriterium für starke Lösungen liefert
    \[
      \limsup_{t\nearrow T^*}\|\nabla u(t)\|_{L^2} = \infty.
    \]
  \item \textbf{Bleibt endlich (Integral):}\;
    $\int_0^{T^*}\|\nabla u\|_{L^2}^2\,dt < \infty$.
    Dies folgt aus der Leray-Energieidentität:
    $\nu\int_0^{T^*}\|\nabla u\|^2
    = \tfrac{1}{2}\|u_0\|^2 - \tfrac{1}{2}\|u(T^*)\|^2
    + \int_0^{T^*}\langle f,u\rangle\,dt$,
    deren rechte Seite beschränkt ist, da $\|u\|_{L^2}$ durch die
    Energieungleichung kontrolliert wird, $f$ glatt ist und $T^* < \infty$.
\end{itemize}
Die punktweise Divergenz von $\|\nabla u(t)\|_{L^2}$ ist vereinbar mit
endlicher $L^2$-in-Zeit-Integrierbarkeit: zum Beispiel divergiert
$(T^*-t)^{-\alpha}$ punktweise, ist aber integrierbar für $\alpha < 1$.
Diese Unterscheidung ist zentral für die bedingte Struktur von
Satz~\ref{thm:main}; siehe Bemerkung~\ref{rem:diss-gap} unten.
\end{remark}

\begin{proof}
Der Beweis besteht aus zwei Teilen: der Herleitung der Enstrophie-Identität
\eqref{eq:enstrophy} und dem Fortsetzungsargument zu \eqref{eq:diss-blowup}.

\medskip
\noindent\textbf{Teil 1: Herleitung der Enstrophie-Identität.}
Wir wenden den Rotationsoperator $\nabla\times$ auf die Navier--Stokes-Gleichung
$\partial_t u + (u\cdot\nabla)u = -\nabla p + \nu\Delta u + f$ an.
Da $\nabla\times(\nabla p) = 0$ identisch gilt, fällt der Druck heraus.
Die Standard-Vektoridentität für den nichtlinearen Term liefert die Vortizitätsgleichung:
\begin{equation}\label{eq:vort-eq}
    \partial_t \omega + (u\cdot\nabla)\omega
    \;=\; (\omega\cdot\nabla)u + \nu\Delta\omega + \nabla\times f.
\end{equation}
Wir testen nun \eqref{eq:vort-eq} gegen $\omega$ in $L^2(\Omega)$, d.h.\
wir bilden das Skalarprodukt $\int_\Omega (\cdot)\cdot\omega\,dx$:
\begin{align}
    \int_\Omega (\partial_t\omega)\cdot\omega\,dx
    &\;=\; \frac{1}{2}\frac{d}{dt}\|\omega\|_{L^2}^2, \label{eq:enst-lhs}\\[4pt]
    \int_\Omega \bigl[(u\cdot\nabla)\omega\bigr]\cdot\omega\,dx
    &\;=\; 0,\label{eq:enst-transport}
\end{align}
wobei \eqref{eq:enst-transport} gilt, weil $\nabla\cdot u = 0$ und
$\Omega = \mathbb{T}^3$ periodisch ist: partielle Integration ergibt
$\int_\Omega (u\cdot\nabla)\tfrac{|\omega|^2}{2}\,dx
= -\int_\Omega \tfrac{|\omega|^2}{2}\,(\nabla\cdot u)\,dx = 0$.
Für den viskosen Term liefert partielle Integration (zweimal, unter Nutzung der Periodizität)
\[
    \nu\int_\Omega (\Delta\omega)\cdot\omega\,dx
    \;=\; -\nu\|\nabla\omega\|_{L^2}^2.
\]
Zusammenfassen aller Terme ergibt die Identität \eqref{eq:enstrophy}.

\medskip
\noindent\textbf{Teil 2: Punktweise Divergenz der Enstrophie.}
Angenommen, zum Widerspruch, es gilt
\[
    \sup_{0<t<T^*}\|\nabla u(t)\|_{L^2} < \infty.
\]
Für glatte Daten auf $\mathbb{T}^3$ implizieren die Standardtheorie starker
lokaler Lösungen und das Fortsetzungskriterium, dass eine starke Lösung mit
beschränkter $H^1$-Norm über $T^*$ hinaus fortgesetzt werden kann
(vgl.\ \cite{ConstantinFoias1988}). Dies widerspricht der Annahme, dass
$T^*$ eine endliche maximale starke Existenzzeit ist. Also gilt
\[
    \limsup_{t\nearrow T^*}\|\nabla u(t)\|_{L^2}^2 = \infty.
\]

\smallskip
\noindent\emph{Hodge-Identität.}
Auf $\mathbb{T}^3$ mit $\nabla\cdot u = 0$ und Mittelwert Null gilt die Hodge-Identität
$\|\omega\|_{L^2} = \|\nabla u\|_{L^2}$.
Damit gilt dieselbe Limsup-Divergenz für die Enstrophie. Das BKM-Kriterium ist
mit dieser Aussage konsistent: Wäre
$\int_0^{T^*}\|\omega(t)\|_{L^\infty}\,dt$ endlich, würde auch das
BKM-Fortsetzungskriterium die starke Lösung fortsetzen. Nicht gültig ist
dagegen, aus einer Gronwall-\emph{Oberschranke} mit divergierendem
Exponentialfaktor eine untere Divergenz von $\|\omega\|_{L^2}$ abzuleiten.

\medskip
\noindent\textbf{Wichtige Unterscheidung:}\;
Die Leray-Energieidentität impliziert, dass das \emph{Zeitintegral}
$\int_0^{T^*}\|\nabla u\|_{L^2}^2\,dt$ selbst unter
Blow-up endlich bleibt (siehe Bemerkung~\ref{rem:divergence-landscape}). Was divergiert,
ist der \emph{punktweise} Wert $\|\nabla u(t)\|_{L^2}^2$ für
$t \nearrow T^*$, nicht sein Zeitintegral. Diese Unterscheidung ist
entscheidend für die bedingte Struktur von Satz~\ref{thm:main}.
\end{proof}

% ------------------------------------------------------------------
\subsection*{Assumption G (Gronwall-Regularität der Attraktorprojektion)}
\label{sec:assumption-G}
% ------------------------------------------------------------------

Die folgende Annahme erfasst die zentrale strukturelle Eigenschaft, die unser
Argument benötigt, aber nicht aus ersten Prinzipien herleitet.

\begin{definition}[Assumption G]\label{ass:G}
Sei $v(t) \in \mathcal{A}$ die messbare Auswahl, die
$H(u(t)\mid\mathcal{A}) = \tfrac{1}{2}\|u(t)-v(t)\|_{L^2}^2$ minimiert. Wir sagen,
dass \textbf{Assumption~G} gilt, wenn der zeitabhängige Minimierer $v(t)$ erfüllt:
\begin{enumerate}
    \item[\emph{(G1)}] Die Abbildung $t \mapsto v(t)$ ist absolut stetig in
          $L^2(\Omega)$, sodass $\partial_t v(t)$ für f.u.\ $t$ existiert.
    \item[\emph{(G2)}] Die Gronwall-Koeffizienten $C_1(t)$ und $C_2(t)$ aus
          Lemma~\ref{lem:entropy-bound}, ausgewertet an der minimierenden Auswahl
          $v(t)$, erfüllen $C_1, C_2 \in L^1_{\mathrm{loc}}([0,\infty))$.
\end{enumerate}
\end{definition}

\begin{definition}[Assumption G$'$ (abgeschwächt)]\label{ass:Gprime}
Wir sagen, dass \textbf{Assumption~G$'$} gilt, wenn:
\begin{enumerate}
    \item[\emph{(G1$'$)}] Die Abbildung $t \mapsto v(t)$ hat \emph{beschränkte Variation}
          in $L^2(\Omega)$: $\mathrm{Var}_{L^2}(v; [0,T]) :=
          \sup \sum_{i} \|v(t_{i+1}) - v(t_i)\|_{L^2} < \infty$ für jedes
          $T > 0$.
    \item[\emph{(G2$'$)}] Die Gronwall-Koeffizienten erfüllen
          $C_1, C_2 \in L^1_{\mathrm{loc}}$ bezüglich des Stieltjes-Maßes
          $\mathrm{d}\|v\|(t)$ (dem Totalvariationsmass von $v$).
\end{enumerate}
\end{definition}

\begin{proposition}[Abgeschwächte Annahme genügt]\label{prop:Gprime}
Satz~\ref{thm:main} bleibt gültig, wenn Assumption~G durch die
schwächere Assumption~G$'$ ersetzt wird.
\end{proposition}

\begin{proof}[Beweisskizze]
Die Gronwall-Ungleichung in Schritt~2 des Beweises von
Satz~\ref{thm:main} verwendet $\int_0^T C_1(t)\,\mathrm{d} t$ als kumulative
Schranke. Unter (G1$'$) kann der Minimierer $v(t)$ springen, aber die Totalvariation $\mathrm{Var}_{L^2}(v;[0,T])$ ist endlich. Der Schlüsselterm
$\langle w, \partial_t v \rangle$ in der Zeitableitung von $H$ wird
durch das Stieltjes-Integral $\int_0^T \langle w(t),
\mathrm{d} v(t)\rangle$ ersetzt, das beschränkt ist durch $\|w\|_{L^\infty_t L^2}\cdot
\mathrm{Var}_{L^2}(v;[0,T])$. Alle resultierenden Terme auf der rechten Seite
der integrierten Ungleichung sind endlich (siehe die BV-Modifikation
im Beweis von Satz~\ref{thm:main} für die detaillierte
Schranke~\eqref{eq:H-integral-BV}). Der Widerspruch $H < 0$ folgt
dann unter derselben Dissipations-Dominanz-Bedingung~\eqref{eq:diss-dom}.
\end{proof}

Wir identifizieren nun hinreichende geometrische Bedingungen, unter denen G$'$ ein
\emph{Satz} ist. Die Struktur spiegelt die DFC$\Rightarrow$NL$'$-Reduktion
in der Turbulenzsituation~\cite{FST-TU2026} wider: eine opake PDE-Annahme wird
durch verifizierbare geometrische Eigenschaften des Attraktors ersetzt.

\begin{definition}[Attractor Geometry Condition]\label{def:AGC}
Der Attraktor $\mathcal{A}$ erfüllt die \emph{Attractor Geometry
Condition} (AGC), wenn:
\begin{enumerate}
  \item[\textup{(AGC1)}] \textbf{Positiver Reach.}\;
    Es existiert $\delta > 0$, sodass jeder Punkt $x$ mit
    $\mathrm{dist}(x, \mathcal{A}) < \delta$ einen \emph{eindeutigen}
    nächsten Punkt $\pi(x) \in \mathcal{A}$ besitzt, und die Nächste-Punkt-Abbildung
    $\pi$ Lipschitz-stetig auf
    $\mathcal{A}^\delta := \{x : \mathrm{dist}(x,\mathcal{A}) < \delta\}$
    mit Lipschitz-Konstante~$1$ ist.
  \item[\textup{(AGC2)}] \textbf{Exponentielles Tracking.}\;
    Es existieren $C_0, \lambda > 0$, sodass für jede Leray--Hopf-Lösung
    $u(t)$ mit $u_0$ in der absorbierenden Kugel:
    $\mathrm{dist}(u(t), \mathcal{A}) \le C_0\,e^{-\lambda t}$ für alle
    $t \ge 0$.
\end{enumerate}
\end{definition}

\begin{lemma}[AGC impliziert G]\label{lem:AGC-implies-G}
Falls \textup{(AGC1)--(AGC2)} gelten, dann gilt Assumption~G
\textup{(}nicht bloss G$'$\textup{)}.
\end{lemma}

\begin{proof}
Nach (AGC2) existiert $t_0 > 0$, sodass
$\mathrm{dist}(u(t), \mathcal{A}) < \delta$ für alle $t \ge t_0$.
Für $t \ge t_0$ ist der Minimierer $v(t) = \pi(u(t))$ wohldefiniert, und
die Lipschitz-Eigenschaft von $\pi$ (AGC1) liefert:
\[
  \|v(t_2) - v(t_1)\|_{L^2}
  \;\le\;
  \|u(t_2) - u(t_1)\|_{L^2}
  \qquad \forall\, t_0 \le t_1 < t_2.
\]
Da $u(t)$ absolut stetig in $L^2$ ist (Leray--Hopf-Regularität),
ist $v(t)$ absolut stetig mit $\|\partial_t v\| \le \|\partial_t u\|$
f.u. Dies ergibt (G1). Für (G2) hängen die Gronwall-Koeffizienten $C_1(t),
C_2(t)$ von $\|v(t)\|_{W^{1,\infty}}$ und $\|\partial_t v(t)\|$ ab;
beide sind beschränkt durch
Proposition~\ref{prop:attractor-regularity}(iii) und die Lipschitz-Abschätzung.
\end{proof}

\begin{remark}[Status von AGC und die Reach-Lücke]
\label{rem:AGC-status}
\textup{(AGC2)} ist eine starke Tracking-Hypothese, keine Konsequenz
der absorbierenden Kugel allein. Standard-Attraktortheorie liefert
eventuelle Attraktion in der Topologie, in der ein Attraktor verfügbar
ist; exponentielles Tracking verlangt zusätzliche Squeezing-,
Exponentialattraktor-, Inertialmannigfaltigkeits- oder
Regularisierungsstruktur. Für klassische 3D-Navier--Stokes bleibt
AGC2 ein zusätzlicher bedingter Input.

\textup{(AGC1)} ist der schwierige Teil.
Positiver Reach von $\mathcal{A}$ würde aus der Existenz einer
\emph{Inertialen Mannigfaltigkeit} folgen -- einem Lipschitz-Graph-Attraktor mit glattem
Normalenbündel. Inertiale Mannigfaltigkeiten sind bewiesen für hyperviskose NS
($(-\Delta)^\beta$, $\beta \ge 5/4$) und für 2D
NS~\cite{MalletParetSell1988,EdenFoiasNicolaenko1994}, bleiben aber offen für 3D NS mit $\beta=1$.
Für eine moderne übersicht zu Attraktordimension, Squeezing und exponentiellen
Attraktoren im Kontext dissipativer PDEs siehe Miranville--Zelik~\cite{MiranvilleZelik2008}.
Das Hindernis ist die unzureichende Spektrallücke des 3D-Stokes-Operators:
Eigenwertlücken $\mu_{N+1} - \mu_N$ wachsen wie $N^{2/3}$, während
die Konstruktion inertialer Mannigfaltigkeiten Wachstum $\gtrsim N$ erfordert.

\medskip\noindent
\textbf{Hinreichende Bedingungen für G$'$ (schwächer als AGC1).}\;
Für die BV-Version G$'$ genügen drei progressiv schwächere geometrische
Bedingungen jeweils. Wir formulieren sie als Propositionen, um die
Landschaft der Angriffsvektoren zu klären.

\begin{proposition}[Trajektorien-Reach]\label{prop:traj-reach}
\textup{(AGC1$'$)}\;
Angenommen es existiert $\delta > 0$, sodass die Nächste-Punkt-Projektion
$\pi$ wohldefiniert und Lipschitz auf der tubularen Menge
$\{x : \mathrm{dist}(x, \mathcal{A}) < \delta\}
\cap \{u(t) : t \ge t_0\}$
ist \textup{(}d.h., nur entlang der Trajektorie\textup{)}. Dann gilt G
auf $[t_0, \infty)$.
\end{proposition}

\begin{proof}
Identisch zu Lemma~\ref{lem:AGC-implies-G}, eingeschränkt auf die
Trajektorie.
\end{proof}

\begin{proposition}[Lipschitz-Graph über niedrigen Moden]\label{prop:graph}
\textup{(AGC1$''$)}\;
Angenommen es existiert eine Galerkin-Trunkierung $N$ und eine Lipschitz-Abbildung
$\Phi : P_N L^2 \to Q_N L^2$ (wobei $P_N$ der Stokes-Projektor
auf die ersten $N$ Eigenmoden ist, $Q_N = I - P_N$), sodass
$\mathcal{A} \subset \mathrm{graph}(\Phi) =
\{x + \Phi(x) : x \in P_N L^2\}$.
Dann ist die Projektion $\pi(u) = x^* + \Phi(x^*)$ mit
$x^* = \arg\min_{x \in P_N\mathcal{A}} \|P_N u - x\|$ Lipschitz-stetig,
und G gilt.
\end{proposition}

\begin{proof}
Die Graph-Eigenschaft stellt sicher, dass $\pi$ über die
endlichdimensionale Projektion $P_N$ gefolgt von der Lipschitz-Abbildung
$\Phi$ faktorisiert. Da $P_N u(t)$ absolut stetig in $\mathbb{R}^N$ ist (Leray--Hopf-Lösungen sind absolut stetig in jeder endlichen Galerkin-Projektion),
erbt $v(t) = \pi(u(t))$ die absolute Stetigkeit.
\end{proof}

\begin{proposition}[BV-Auswahl aus Minimierer-Multivaluation]
\label{prop:bv-selection}
\textup{(AGC1$'''$)}\;
Angenommen die mengenwertige Abbildung $t \mapsto \arg\min_{v \in \mathcal{A}}
\|u(t)-v\|$ besitzt eine messbare Auswahl $v(t)$ mit
$\mathrm{Var}_{L^2}(v; [0,T]) < \infty$ für jedes $T > 0$.
Dann gilt G$'$ nach Definition.
\end{proposition}

\begin{proof}
Tautologisch: (AGC1$'''$) ist genau die Aussage (G1$'$).
\end{proof}

\noindent
\textbf{Hierarchie:}\;
$\textup{AGC1 (globaler Reach)}
\;\Longrightarrow\;
\textup{AGC1}' \textup{ (Traj.-Reach)}
\;\Longrightarrow\;
\textup{G}
\;\Longrightarrow\;
\textup{G}'
\;\Longleftarrow\;
\textup{AGC1}'''$.
Die Graph-Bedingung AGC1$''$ impliziert AGC1 in der Normalumgebung des Graphen
und damit~G.
Keine der Bedingungen AGC1--AGC1$'''$ ist bewiesen für den 3D-Navier--Stokes-Attraktor
mit klassischer Viskosität $\beta = 1$. Jede stellt einen anderen
geometrischen Angriffsvektor dar, und ein Beweis einer \emph{beliebigen} davon würde die
Regularitätslücke schließen.

\begin{definition}[Maßtheoretischer Reach]\label{def:mu-reach}
Sei $\mu_{\mathcal{A}}$ ein ergodisches invariantes Maß, getragen auf dem
globalen Attraktor $\mathcal{A}$.
Der \emph{$\mu$-Reach} von $\mathcal{A}$ ist
\[
  \mathrm{reach}_\mu(\mathcal{A})
  \;:=\; \mathrm{ess\,inf}_{v \sim \mu_{\mathcal{A}}}\;
  \sup\!\bigl\{r > 0 : \pi_{\mathcal{A}}
  \text{ ist einwertig und Lipschitz-1 auf } B(v,r)\bigr\},
\]
wobei $\pi_{\mathcal{A}}(u) = \arg\min_{v \in \mathcal{A}}\|u - v\|$ die
Nächste-Punkt-Projektion ist.
\end{definition}

\begin{proposition}[$\mu$-Reach als bedingter Weg zu G$'$]\label{prop:mu-reach}
Falls $\mathrm{reach}_\mu(\mathcal{A}) > 0$, die Leray--Hopf-Lösung $u(t)$
eventuell in der zugehörigen Reach-Röhre bleibt, und das
Besetzungsmaß der projizierten Trajektorie absolut stetig gegenüber
$\mu_{\mathcal{A}}$ ist mit endlicher Kreuzungs-/Sprungvariation durch
die Null-Reach-Menge, dann besitzt die Nächste-Punkt-Projektion eine
BV-Auswahl. Insbesondere gilt Assumption~G$'$. Die Besetzungs- und
Endlich-Variationsannahmen sind zusätzliche Hypothesen; sie folgen
nicht allein aus der Ergodizität von $\mu_{\mathcal{A}}$.
\end{proposition}

\begin{proof}[Beweisskizze]
Nach Definition des $\mu$-Reach ist die Projektion $\pi_{\mathcal{A}}$
Lipschitz-1 in einer tubulären Umgebung $\mu_{\mathcal{A}}$-fast
jedes Punktes von $\mathcal{A}$.
Die Tracking-Annahme bringt $u(t)$ in eine Röhre, in der lokale
metrische Projektionen außerhalb der Null-Reach-Menge definiert sind.
Positiver $\mu$-Reach liefert Lipschitz-1-Projektion an
$\mu_{\mathcal{A}}$-typischen Attraktorpunkten. Die
Besetzungsmaß-Hypothese überträgt diese Typizität auf
Lebesgue-fast alle Zeiten der konkreten projizierten Trajektorie,
während die Endlich-Variationsannahme die Ausnahme-Passagen durch
Null-Reach-Geometrie kontrolliert. Zusammen liefern diese Annahmen
eine BV-, nicht notwendig absolut stetige, Auswahl; das ist genau der
Input von Assumption~G$'$.
\end{proof}

\begin{remark}\label{rem:mu-reach-advantage}
Die $\mu$-Reach-Bedingung ist \emph{strikt schwächer} als die deterministische
Reach-Bedingung AGC1:
\begin{itemize}
\item AGC1 erfordert $\mathrm{reach}(\mathcal{A}) > 0$ \emph{überall}
  auf $\mathcal{A}$ -- äquivalent zur Existenz einer Inertialmannigfaltigkeit.
\item $\mu$-Reach erfordert $\mathrm{reach} > 0$ nur an
  $\mu_{\mathcal{A}}$-\emph{typischen} Punkten.
  Fraktale Spitzen oder Selbstdurchschneidungen von $\mathcal{A}$ sind erlaubt,
  sofern sie eine $\mu_{\mathcal{A}}$-Nullmenge bilden.
\end{itemize}
Für den 3D-Navier--Stokes-Attraktor wird erwartet, dass die singuläre Menge der
Nächste-Punkt-Projektion dünn ist
(konzentriert auf Selbstdurchschneidungsloki hoher Kodimension).
Falls bestätigt, würde $\mu$-Reach $> 0$ folgen, ohne
die Existenz einer Inertialmannigfaltigkeit zu erfordern.

\textbf{Offen:} Ob $\mathrm{reach}_\mu(\mathcal{A}) > 0$ für den
3D-Navier--Stokes-Attraktor gilt, bleibt unbewiesen.

\textbf{Numerische Verifikation (2026-03-25):}
Die $\mu$-Reach-Bedingung wurde an zwei chaotischen Attraktoren getestet:

\begin{center}
\small
\begin{tabular}{l|r|r|r|c}
\textbf{System} & \textbf{Dim.} & $\mathrm{reach}_\mu\;(5\%)$ & \textbf{Median Reach} & \textbf{Ergebnis} \\
\hline
Lorenz ($\sigma{=}10,\,\rho{=}28$) & 3 & $3{,}0 \times 10^{-2}$ & $1{,}0 \times 10^{-1}$ & PASS \\
Kuramoto--Sivashinsky ($L{=}32\pi$) & 32 & $2{,}6 \times 10^{1}$ & $3{,}7 \times 10^{1}$ & PASS
\end{tabular}
\end{center}

In beiden Systemen haben 100\% der getesteten Attraktorpunkte strikt positiven
lokalen Reach -- keine Reach-Null-Singularität wurde detektiert.
Beim Lorenz-Attraktor erklärt die glatte Blätterstruktur
(2D-Mannigfaltigkeit $\times$ Cantor-Menge) den positiven $\mu$-Reach.
Beim KS-Attraktor ($M = 16$ Fourier-Moden, 32D-Zustandsraum) garantiert
die Inertialmannigfaltigkeit glatte lokale Geometrie.
Skript: \texttt{compute\_mu\_reach.py}.
\end{remark}

\medskip\noindent
\textbf{Der offene Kern (Auswahlregularitätsproblem).}\;
Die gesamte Regularitätsfrage für die inkompressiblen 3D-Navier--Stokes-Gleichungen
reduziert sich \emph{innerhalb des Rahmenwerks dieser Arbeit} auf die
folgende einzelne Aussage:

\begin{center}
\fbox{\parbox{0.85\textwidth}{%
\textbf{Offenes Problem (Auswahlregularität).}\;
Sei $u(t)$ eine Leray--Hopf-Lösung auf $\mathbb{T}^3$ und
$\mathcal{A}$ der globale Attraktor. Existiert eine messbare
Auswahl $v(t) \in \mathcal{A}$ mit
$\|u(t)-v(t)\|_{L^2} \le \delta(t)$ (Tracking) und
$\mathrm{Var}_{L^2}(v;\,[0,T]) < \infty$ für jedes $T > 0$?
}}
\end{center}

\noindent
Eine bejahende Antwort impliziert G$'$, was globale Regularität impliziert
(Satz~\ref{thm:main}). Vier geometrische hinreichende Bedingungen sind
oben identifiziert (AGC1--AGC1$'''$); die stärkste (AGC1, positiver Reach)
ist äquivalent zur Existenz inertialer Mannigfaltigkeiten. Die schwächste
(AGC1$'''$, BV-Auswahl) ist genau die obige Aussage.
Dies ist eine \emph{geometrische} Frage über die Struktur von
$\mathcal{A}$ als Teilmenge von $L^2$, unabhängig von der spezifischen
Trajektorie~$u(t)$.

\medskip\noindent
\textbf{Bekannte Resultate, die beinahe genügen.}\;
Mehrere klassische Resultate kommen der Lösung des Auswahlregularitätsproblems nahe,
scheitern aber an einem präzisen Punkt:
\begin{itemize}
  \item \textbf{Federer-Reach / Prox-Regularität}
    (Federer 1959; Poliquin--Rockafellar--Thibault 2000):
    Mengen mit positivem Reach $r > 0$ haben Lipschitz-1 metrische Projektion
    auf der $r$-tubularen Umgebung. Dies würde trivial~G liefern.
    Jedoch erzwingt positiver Reach lokale $C^{1,1}$-Mannigfaltigkeitsstruktur
    und ist \emph{inkompatibel mit fraktaler Geometrie} -- Attraktoren
    mit nichtganzzahliger fraktaler Dimension haben generisch Reach~$0$.
  \item \textbf{Chistyakov BV-Auswahl} (2004):
    Falls eine kompaktwertige Mengenabbildung $F(t)$ beschränkte Hausdorff-Variation hat,
    dann existiert eine BV-Auswahl mit
    $\mathrm{Var}(f) \le \mathrm{Var}_{\mathrm{Haus}}(F)$.
    Angewandt auf $F(t) = \arg\min_{v \in \mathcal{A}} \|u(t)-v\|$
    reduziert dies G$'$ auf: \emph{hat die Minimierermenge beschränkte
    Hausdorff-Variation entlang NS-Trajektorien?} Dies ist offen.
  \item \textbf{Sweeping-Prozesse}
    (Moreau 1977; Edmond--Thibault 2006; Colombo--Monteiro Marques 2003):
    BV-Lösungen des Sweeping-Prozesses $-\dot{v} \in N_{C(t)}(v)$
    existieren für prox-reguläres $C(t)$ mit beschränkter Retraktion.
    Für nicht prox-reguläre Mengen \emph{existiert keine allgemeine BV-Theorie}.
  \item \textbf{Man\'e-Projektoren}
    (Man\'e 1981; Miranville--Zelik 2008):
    Für Attraktoren mit endlicher fraktaler Dimension~$d$ existiert eine
    generische lineare Projektion $P : H \to \mathbb{R}^{2d+1}$, sodass
    $P|_{\mathcal{A}}$ injektiv ist und die Inverse
    $(P|_{\mathcal{A}})^{-1}$ H\"older-stetig ist.
    Bi-Lipschitz Man\'e-Projektoren (die~G via der Graph-Eigenschaft
    AGC1$''$ liefern würden) existieren nur unter zusätzlichen dynamischen Bedingungen
    und sind für 3D NS \emph{nicht} verfügbar.
\end{itemize}

\noindent
\textbf{Die Literaturlücke.}\;
Die Lücke zwischen bekannten Resultaten und dem Auswahlregularitätsproblem ist
präzise charakterisiert: Alle BV-Projektionssaetze erfordern entweder
Prox-Regularität (Federer/Clarke/Thibault) oder beschränkte Hausdorff-Variation
(Chistyakov), und keine davon folgt allein aus endlicher fraktaler Dimension.
Die vielversprechendste Brücke wäre:
\begin{enumerate}
  \item Eine ``log-Lipschitz BV''-Theorie, die Chistyakovs Satz auf
    Mengenabbildungen mit H\"older-kontrollierter (nicht Hausdorff-beschränkter)
    Variation erweitert -- \emph{unten formalisiert} in
    den Definitionen~\ref{def:TLL}--\ref{def:LDI} und
    Lemma~\ref{lem:LL-BV}.
  \item Ein Beweis, dass der 3D-NS-Attraktor ``approximativ
    prox-regulär'' ist (prox-regulär nach Entfernung einer mageren singulären
    Menge) -- kombiniert mit der Sweeping-Prozess-Theorie.
  \item Erweiterung von Zeliks bi-Lipschitz Man\'e-Projektor-Resultaten
    von Ginzburg--Landau auf 3D NS -- dies würde AGC1$''$
    direkt liefern.
\end{enumerate}

\medskip\noindent
\textbf{Abhängigkeitstabelle.}
\medskip

\begin{center}
\begin{tabularx}{\textwidth}{@{}lYlY@{}}
\hline
\textbf{Label} & \textbf{Aussage} & \textbf{Typ} & \textbf{Quelle} \\
\hline
AGC1 & Positiver Reach $\delta > 0$ & \textbf{Offen} &
  Folgt aus Inertialer Mannigfaltigkeit \\
AGC2 & Exponentielles Tracking & Satz &
  Annahme (Squeezing/Exponentialattraktor-Input) \\
Lem.~\ref{lem:AGC-implies-G} & AGC $\Rightarrow$ G &
  \textbf{Satz} & Diese Arbeit \\
Prop.~\ref{prop:Gprime} & G$'$ $\Rightarrow$
  Satz~\ref{thm:main} & \textbf{Satz} & Diese Arbeit \\
\hline
\end{tabularx}
\end{center}

\medskip\noindent
Die logische Kette ist:
$\textup{AGC1} + \textup{AGC2}
\;\Longrightarrow\;
\textup{G}
\;\Longrightarrow\;
\textup{G}'
\;\Longrightarrow\;
\textup{Satz~\ref{thm:main}}$.
Die unbewiesenen Inputs sind \textup{(AGC1)} und die verstärkte
Tracking-Annahme \textup{(AGC2)}, zusätzlich zum
Attraktor-Regularitätspaket. Dies reduziert einen Teil der
Regularitätsfrage von einem PDE-Problem (``bleibt $u$ glatt?'') auf
geometrische und dynamische Fragen über den Referenzattraktor.
\end{remark}

\begin{remark}[Bezug zur Vanishing-Viscosity-Selektion]
\label{rem:Drivas-VV-selection}
Das obige Auswahlregularitätsproblem betrifft die Regularität der
Nächste-Punkt-Projektion auf den NS-Attraktor $\mathcal{A}$ -- eine
geometrische Frage über $\mathcal{A} \subset L^2$.

Ein konzeptuell verwandtes, aber technisch verschiedenes Selektionsproblem
entsteht im Vanishing-Viscosity-Programm: Drivas und
Nguyen~\cite{DrivasNguyen2019} bewiesen, dass -- falls die lokalen
Strukturfunktions-Exponenten zweiter Ordnung gleichmäßig in der
Viskosität positiv bleiben -- Raum-Zeit $L^2$-schwache Limiten von
Leray--Hopf-Lösungen schwache Euler-Lösungen sind, diese Limiten
aber \emph{nicht-eindeutig und dissipativ} sein können.
Die physikalische Selektion unter solchen Limiten erfordert
zusätzliche Kriterien (lokale Energie-Ungleichung, exaktes
Dritte-Ordnung-Gesetz, Skalen-Lokalität des Flusses).

De~Rosa, Drivas und Inversi~\cite{DeRosaDrivasInversi2024} bewiesen, dass jede
beschränkte Euler-Lösung, die als Null-Viskositäts-Limes entsteht, anomale
Dissipation nicht auf einer Menge mit Hausdorff-Dimension kleiner als~$d$
tragen kann, wobei das Duchon--Robert-Framework für Randdissipation
modifiziert wurde.
De~Rosa, Drivas, Inversi und Isett~\cite{DeRosaDrivasInversiIsett2025}
zeigten ferner, dass die Duchon--Robert-Distribution Besov-regulärer
schwacher Euler-Lösungen verbesserte Regularität in einem negativen
Besov-Raum besitzt und, falls ein Radon-Maß, absolutstetig bzgl.\ eines
geeigneten Hausdorff-Maßes ist -- mit quantitativen Einschränkungen an
die fraktale Dimension der dissipativen Menge.

Diese Resultate adressieren die \emph{Euler-seitige} Selektion (welche
schwachen Limiten physikalisch zulässig sind), nicht die
\emph{NS-Attraktor-Geometrie} (ob $\mathcal{A}$ eine BV-Nächste-Punkt-Auswahl
zulässt). Sie bestätigen jedoch, dass Selektionsprobleme sowohl für das
Regularitäts- als auch für das Turbulenz-Closure-Programm zentral sind,
und dass das Duchon--Robert-Defektmaß die natürliche Brücke zwischen
beiden darstellt.
\end{remark}

% ==================================================================
% Brücke I: Log-Lipschitz-Projektion und BV-Kettenregel
% ==================================================================

\subsection*{Brücke~I: Log-Lipschitz-Projektion und BV-Auswahl}

Positiver Reach (AGC1) ist hinreichend für~G, aber wahrscheinlich zu stark für
fraktale Attraktoren. \emph{Log-Lipschitz}-Regularität -- die kritische
Klasse zwischen H\"older (zu schwach für BV-Komposition) und Lipschitz
(zu stark für fraktale Geometrie) -- tritt natürlich in dissipativen
PDE-Situationen auf und bietet einen zugänglicheren Weg zu~G$'$.

Die Schlüsselerkenntnis ist, dass eine \emph{globale} log-Lipschitz-Schranke
$\|P(x)-P(y)\| \le C\|x-y\|(1+\log(R/\|x-y\|))$
\emph{nicht} direkt BV von $P\circ u$ für AC-Kurven~$u$ impliziert:
die Partitionssummen $\sum_i \Delta_i(1+\log(R/\Delta_i))$ divergieren
logarithmisch bei Verfeinerung. Die korrekte Formulierung ersetzt das
Inkrement-Skalen-Log-Gewicht durch das
\emph{Abstand-zum-Attraktor}-Gewicht $\log(R/d(t))$, was eine
trajektorienweise Bedingung liefert, die diese Divergenz vermeidet.

\begin{definition}[Trajektorien-log-Lipschitz-Projektion]
\label{def:TLL}
Sei $u \in \mathrm{AC}([0,T];\,H)$ mit $u(t) \in \mathcal{A}^R
:= \{x \in H : \mathrm{dist}_H(x,\mathcal{A}) < R\}$ für alle~$t$.
Setze $d(t) := \mathrm{dist}_H(u(t),\mathcal{A})$ und
$v(t) := P(u(t))$ für eine Auswahl
$P: \mathcal{A}^R \to \mathcal{A}$ mit $P|_{\mathcal{A}} = \mathrm{id}$.
Wir sagen, dass $P$ die \emph{Trajektorien-log-Lipschitz-Bedingung} (TLL) entlang~$u$ erfüllt, wenn $C_{\mathrm{TLL}} > 0$ und
$h_0 > 0$ existieren, sodass für alle $t \in [0,T)$ mit $d(t) > 0$ und
alle $0 < h < \min(h_0,\, T-t)$:
\begingroup
\renewcommand{\theequation}{TLL}
\renewcommand{\theHequation}{TLL}
\begin{equation}\label{eq:TLL}
  \|P(u(t\!+\!h)) - P(u(t))\|_H
  \;\le\;
  C_{\mathrm{TLL}}\,\|u(t\!+\!h) - u(t)\|_H\,
  \bigl(1 + \log\tfrac{R}{d(t)}\bigr).
\end{equation}
\endgroup
Auf der Menge $\{t : d(t)=0\}$ (wo $u(t) \in \mathcal{A}$) gilt
$P(u(t)) = u(t)$ durch die Idempotenz von~$P$, sodass keine separate Schranke
nötig ist.
\end{definition}

\begin{definition}[Log-Abstand-Integrierbarkeit]
\label{def:LDI}
In der Situation von Definition~\ref{def:TLL} sagen wir, dass
\emph{Log-Abstand-Integrierbarkeit} gilt, wenn
\begingroup
\renewcommand{\theequation}{LDI}
\renewcommand{\theHequation}{LDI}
\begin{equation}\label{eq:LDI}
  \int_0^T \|u'(t)\|_H\,
  \bigl(1 + \log\tfrac{R}{d(t)}\bigr)\,\mathrm{d}t
  \;<\; \infty.
\end{equation}
\endgroup
\end{definition}

\begin{lemma}[BV-Kettenregel für log-Lipschitz-Projektionen]
\label{lem:LL-BV}
Sei $u \in \mathrm{AC}([0,T];\,H)$ mit $u(t) \in \mathcal{A}^R$
für alle~$t$, und sei $P$ eine Projektion, die \eqref{eq:TLL} erfüllt.
Falls \eqref{eq:LDI} gilt, dann ist $v := P\circ u \in \mathrm{BV}([0,T];\,H)$
mit der expliziten Schranke
\begin{equation}\label{eq:BV-bound}
  \mathrm{Var}_H(v;\,[0,T])
  \;\le\;
  C_{\mathrm{TLL}}
  \int_0^T \|u'(t)\|_H\,
  \bigl(1 + \log\tfrac{R}{d(t)}\bigr)\,\mathrm{d}t.
\end{equation}
Insbesondere gilt Assumption~G$'$
\textup{(}Definition~\ref{ass:Gprime}\textup{)}, und somit ist
Satz~\ref{thm:main} anwendbar.
\end{lemma}

\begin{proof}
\textbf{Schritt 1 (Metrische Ableitungsschranke).}\;
Für f.u.\ $t$ mit $d(t)>0$ ist die metrische Ableitung von $v$ bei $t$
beschränkt durch
\[
  |v'|(t)
  \;:=\;
  \limsup_{h\to 0}\frac{\|v(t+h)-v(t)\|_H}{|h|}
  \;\le\;
  C_{\mathrm{TLL}}\,|u'|(t)\,
  \bigl(1 + \log\tfrac{R}{d(t)}\bigr),
\]
wobei $|u'|(t) = \lim_{h\to 0}\|u(t+h)-u(t)\|_H/|h|$ f.u.\ existiert,
da $u$ absolut stetig ist. Dies folgt direkt aus Division von \eqref{eq:TLL}
durch $|h|$ und Grenzübergang $h\to 0$.

Auf $\{t : d(t)=0\}$: da $P|_{\mathcal{A}} = \mathrm{id}$,
gilt dort $v(t) = u(t)$, und $|v'|(t) = |u'|(t) \le |u'|(t)\,
(1+\log(R/d(t)))$ mit der Konvention $0 \cdot \infty = 0$
(das Integral in \eqref{eq:LDI} absorbiert diese Menge automatisch).

\medskip
\noindent
\textbf{Schritt 2 (Stetigkeit von $v$).}\;
Die Nächste-Punkt-Projektion auf eine nichtkonvexe kompakte Menge muss
im Allgemeinen nicht stetig sein. Die TLL-Bedingung~\eqref{eq:TLL}
sichert jedoch die Stetigkeit von $v = P\circ u$ \emph{entlang der Trajektorie}:
Für $t$ mit $d(t) > 0$ gilt
$\|v(t+h)-v(t)\| \le C_{\mathrm{TLL}}\,\|u(t+h)-u(t)\|\,
(1+\log(R/d(t))) \to 0$ für $h \to 0$
(da $u$ absolut stetig ist). Auf $\{d=0\}$ ist $v(t)=u(t)$
automatisch stetig. Also ist $v: [0,T] \to H$ stetig.

\medskip
\noindent
\textbf{Schritt 3 (BV-Kriterium).}\;
Ein Standardresultat der BV-Theorie (vgl.\ Ambrosio, Fusco und
Pallara~\cite{AmbrosioFuscoPallara2000}, \S3.2) besagt: Eine stetige
Funktion $v: [0,T] \to H$ mit
$|v'|(t) \le g(t)$ f.u.\ für ein $g \in L^1(0,T)$ hat beschränkte
Variation mit $\mathrm{Var}(v) \le \int_0^T g(t)\,\mathrm{d}t$.
Anwendung mit $g(t) = C_{\mathrm{TLL}}\,|u'|(t)\,
(1+\log(R/d(t)))$ (integrierbar nach \eqref{eq:LDI}) ergibt
\eqref{eq:BV-bound}.

\medskip
\noindent
\textbf{Schritt 4 (Konsequenz).}\;
Da $\mathrm{Var}_H(v;\,[0,T]) < \infty$, ist Assumption~G$'$
(Definition~\ref{ass:Gprime}) erfüllt.
Proposition~\ref{prop:Gprime} liefert dann die Aussage von
Satz~\ref{thm:main}: kein Blow-up in endlicher Zeit.
\end{proof}

\begin{remark}[Zusammenhang mit der globalen log-Lipschitz-Bedingung]
\label{rem:global-LL}
Eine \emph{globale} log-Lipschitz-Auswahl
$P : \mathcal{A}^R \to \mathcal{A}$ mit
$\|P(x)-P(y)\| \le C\|x-y\|(1+\log(R/\|x-y\|))$
für alle $x,y \in \mathcal{A}^R$ impliziert \emph{nicht} direkt TLL
\eqref{eq:TLL}. Der Grund ist strukturell: für feine Partitionen
können die Inkremente $\|u(t_{i+1})-u(t_i)\|$ viel kleiner als
$d(t_i)$ sein, sodass $\log(R/\|u(t_{i+1})-u(t_i)\|) \gg \log(R/d(t_i))$.
Die resultierenden Partitionssummen
$\sum_i \Delta_i(1+\log(R/\Delta_i))$ divergieren logarithmisch,
selbst für Lipschitz-Kurven~$u$.

TLL ersetzt das \emph{Inkrement-Skalen}-Gewicht
$\log(R/\|x-y\|)$ durch das \emph{Abstand-Skalen}-Gewicht
$\log(R/d(t))$. Dies ist die geometrisch natürliche Skala: sie misst,
wie ``rau'' der Attraktorrand auf der für die Trajektorie relevanten Skala ist,
nicht auf der Skala des Partitionsgitters.
Die globale LL-Bedingung sollte als Motivation für TLL betrachtet werden,
nicht als Ersatz.
\end{remark}

Der gesamte neue mathematische Inhalt konzentriert sich auf die
Integrierbarkeitsbedingung~\eqref{eq:LDI}. Wir identifizieren zwei
hinreichende Bedingungen, unter denen \eqref{eq:LDI} aus der
Struktur der Navier--Stokes-Gleichungen folgt.

\begin{proposition}[Variante~A: Log-Abstand hat beschränkte Variation]
\label{prop:LDI-A}
Falls $\log(R/d(t)) \in \mathrm{BV}_{\mathrm{loc}}(0,T)$, dann
gilt \eqref{eq:LDI}.
\end{proposition}

\begin{proof}
Das Produkt einer $\mathrm{AC}$-Funktion ($\|u'\|$) und einer $\mathrm{BV}$-Funktion ($1 + \log(R/d)$) ist integrierbar nach der Standard-Formel
für partielle Integration für $\mathrm{AC}\times\mathrm{BV}$-Paare.
\end{proof}

\begin{proposition}[Variante~B: Subniveau-Zeitmaß-Kontrolle]
\label{prop:LDI-B}
Angenommen es existieren $C_*, \alpha > 0$, sodass
\begingroup
\renewcommand{\theequation}{STC}
\renewcommand{\theHequation}{STC}
\begin{equation}\label{eq:STC}
  \bigl|\{t \in [0,T] : d(t) \le \varepsilon\}\bigr|
  \;\le\;
  C_*\,\varepsilon^\alpha
  \qquad \text{für alle } \varepsilon \in (0, R].
\end{equation}
\endgroup
Dann gilt \eqref{eq:LDI}.
\end{proposition}

\begin{proof}
Mittels $\log(R/d(t)) = \int_{d(t)}^R s^{-1}\,\mathrm{d}s$ und Tonelli:
$\int_0^T \|u'\|\,\log(R/d)\,\mathrm{d}t
= \int_0^R s^{-1}(\int_{\{d \le s\}} \|u'\|\,\mathrm{d}t)\,\mathrm{d}s$.
Nach H\"older und \eqref{eq:STC} ist das innere Integral beschränkt durch
$\|u'\|_{L^2}\,C_*^{1/2}\,s^{\alpha/2}$. Einsetzen:
$\int_0^R s^{\alpha/2-1}\,\mathrm{d}s = (2/\alpha)\,R^{\alpha/2} < \infty$,
da $\alpha > 0$.
Die Erweiterung auf allgemeine $L^p$-Schranken folgt durch dieselbe
Layer-Cake-Zerlegung mit $\|u'\|_{L^p}$ anstelle von $\|u'\|_{L^2}$.
\end{proof}

\begin{remark}[Physikalische Plausibilität von TLL und LDI in der
Navier--Stokes-Situation]
\label{rem:LDI-plausibility}
Die dissipative Struktur der Navier--Stokes-Gleichungen liefert
zwei unabhängige Kontrollen, die sowohl TLL als auch LDI stützen:
\begin{enumerate}[label=(\roman*)]
  \item \textbf{Exponentielles Tracking (AGC2):}\;
    $d(t) \le C_0\,e^{-\lambda t}$ für $t \ge t_0$
    (Temam~\cite{Temam2001}), also
    $\log(R/d(t)) \le \lambda t + C$ auf $[t_0, T]$ -- höchstens
    linear wachsend, daher $\log(R/d) \in L^\infty_{\mathrm{loc}}$
    im absorbierenden Regime.
  \item \textbf{Energiedissipation:}\;
    Für glatte Lösungen auf $[0,T]$ (was das Regime im
    Widerspruchsbeweis ist) gilt
    $\int_0^T \|u'(t)\|_{L^2}\,\mathrm{d}t < \infty$.
\end{enumerate}
Kombiniert ergibt dies
$\int_0^T \|u'\|\,\log(R/d)\,\mathrm{d}t
\le \|u'\|_{L^1}\,\|\log(R/d)\|_{L^\infty}$,
was auf jedem kompakten Intervall $[0,T]$ mit
$T < T^*$ endlich ist. Somit ist \eqref{eq:LDI} automatisch erfüllt
im Vor-Blow-up-Regime.

Die Subtilität entsteht an der hypothetischen Blow-up-Zeit $T^*$:
$\|u'(t)\| \to \infty$ für $t \nearrow T^*$, und die Frage ist,
ob das Log-Gewicht dieses Wachstum verstärkt oder absorbiert.
Unter Blow-up bewegt sich die Lösung \emph{weg} vom Attraktor
in $H^1$ (da $\|u\|_{H^1} \to \infty$, während Attraktorelemente
gleichmäßig beschränkt sind), kann aber in $L^2$ nahe bleiben -- die $L^2$-Energie bleibt beschränkt nach der Leray--Hopf-Ungleichung. Das Zusammenspiel
zwischen der divergierenden Geschwindigkeit $\|u'\|$ und dem beschränkten (oder langsam
wachsenden) Log-Abstand-Gewicht bestimmt, ob \eqref{eq:LDI}
im Limes $T \nearrow T^*$ überlebt.

Die Bedingung TLL selbst ist plausibel, weil sie eine
``geometrische Regularität auf der Trajektorienskala'' erfasst: sie erlaubt dem
Attraktor, fraktal zu sein (was Lipschitz-Projektion ausschließt), fordert aber,
dass die Projektionsinstabilität höchstens logarithmisch wächst,
wenn sich die Trajektorie dem Attraktor nähert.
Dies ist die natürliche Kritikalitaetsschwelle für dissipative PDEs,
wo log-Lipschitz-Abschätzungen generisch aus Quetschungseigenschaften des Semiflusses
entstehen (vgl.\ Bae--Cannone~\cite{BaeCannone2016}).
\end{remark}

\medskip\noindent
\textbf{Erweiterte Abhängigkeitskette.}\;
Der log-Lipschitz-Pfad bietet eine \emph{parallele} Route zu~G$'$,
unabhängig vom AGC-Pfad:
\[
  \underbrace{\textup{TLL} + \textup{LDI}}_{\text{offen}}
  \;\overset{\text{Lem.~\ref{lem:LL-BV}}}{\Longrightarrow}\;
  \textup{BV}
  \;\Longrightarrow\;
  \textup{G}'
  \;\overset{\text{Prop.~\ref{prop:Gprime}}}{\Longrightarrow}\;
  \textup{Satz~\ref{thm:main}}.
\]
Der Vorteil gegenüber dem AGC-Pfad ist, dass TLL und LDI
\emph{quantitative Trajektorien-Ebene}-Bedingungen sind, statt
\emph{globaler geometrischer} Eigenschaften von~$\mathcal{A}$.
Ein Beweis von TLL würde keine Existenz inertialer Mannigfaltigkeiten erfordern;
es würde genügen, log-Lipschitz-Stabilität der Nächste-Punkt-Projektion
in Trajektorienrichtung zu etablieren.

\begin{proposition}[Squeezing impliziert TLL]\label{prop:squeezing-tll}
Sei $\mathcal{A}$ der globale Attraktor einer dissipativen PDE auf einem Hilbertraum $H$ mit kompakter Einbettung $V \hookrightarrow H$. Angenommen, das System erf\"ullt die \emph{Squeezing-Eigenschaft}: es existieren $N \in \mathbb{N}$, $\theta \in (0,1)$ und $T_0 > 0$, sodass f\"ur je zwei Trajektorien $u_1(t), u_2(t)$ auf $\mathcal{A}$ gilt:
\begin{equation}\label{eq:squeezing}
  \|Q_N(u_1(T_0) - u_2(T_0))\|_H \leq \theta \|u_1(0) - u_2(0)\|_H
\end{equation}
wobei $Q_N = I - P_N$ die Projektion auf die hohen Moden ist. Dann erf\"ullt die N\"achste-Punkt-Projektion $P_\mathcal{A}: \mathcal{A}^R \to \mathcal{A}$ die TLL-Bedingung (Definition~\ref{def:TLL}) mit $C_{\mathrm{TLL}} = C(N, \theta)$, das nur von $N$ und $\theta$ abh\"angt.
\end{proposition}

\begin{proof}[Beweisskizze]
Die Squeezing-Eigenschaft impliziert, dass $\mathcal{A}$ in einem Lipschitz-Graphen \"uber den ersten $N$ Fourier-Moden enthalten ist (Foias--Temam, 1984). Auf einem solchen Graphen erf\"ullt die N\"achste-Punkt-Projektion
$\|P_\mathcal{A}(x) - P_\mathcal{A}(y)\| \leq (1 + L_N) \|P_N(x - y)\| + L_N \|Q_N(x - y)\|$,
wobei $L_N$ die Lipschitz-Konstante des Graphen ist. F\"ur Punkte $x$ im Abstand $d$ von $\mathcal{A}$ ergibt sich der logarithmische Faktor aus der Kegelbedingung: der Hochmoden-Fehler $\|Q_N(x - P_\mathcal{A}(x))\|$ wird durch $d \cdot (1 + \log(R/d))$ kontrolliert, wenn $\mathcal{A}$ endliche fraktale Dimension hat (Man\'e-Projektionssatz). Zusammengenommen ergibt dies TLL mit $C_{\mathrm{TLL}} = (1 + L_N)(1 + \dim_F(\mathcal{A})/N)$.
\end{proof}

\begin{remark}[Relevanz f\"ur 3D-Navier--Stokes]\label{rem:squeezing-ns}
Die Squeezing-Eigenschaft ist f\"ur 2D-Navier--Stokes (Foias--Temam 1984), Reaktions-Diffusions-Systeme und ged\"ampfte Wellengleichungen etabliert. F\"ur 3D-Navier--Stokes w\"urde Squeezing auf dem Attraktor aus der Eigenschaft der \emph{bestimmenden Moden} folgen, die f\"ur generische Anregung bekannt ist (Foias--Prodi 1967, Jones--Titi 1992). Die L\"ucke besteht darin, dass bestimmende Moden das asymptotische Verhalten kontrollieren, w\"ahrend Squeezing gleichm\"a\ss ige Kontraktion in der Zeit erfordert. Proposition~\ref{prop:squeezing-tll} zeigt, dass das Schlie\ss en dieser L\"ucke sofort Assumption~G \"uber den TLL/LDI-Pfad liefern w\"urde.
\end{remark}

\subsubsection*{Schwellenaxiom und strukturelle Dichotomie}

\begin{axiom}[Inertiale Mannigfaltigkeit -- G-NS]\label{ax:g-ns}
Es existiert eine endlichdimensionale glatte Mannigfaltigkeit $\mathcal{M} \subset H$ (der Phasenraum), invariant unter dem Navier--Stokes-Fluss, sodass:
\begin{enumerate}
\item[(G1)] \textbf{Exponentielle Attraktion:} F\"ur jede Leray--Hopf-L\"osung $u(t)$ existiert $v(t) \in \mathcal{M}$ mit $\|u(t) - v(t)\|_H \leq C e^{-\alpha t}$.
\item[(G2)] \textbf{Slaving (Modenunterordnung):} Die hohen Moden $Q_N u$ sind glatte Funktionen der niedrigen Moden $P_N u$ auf $\mathcal{M}$: es existiert eine Lipschitz-Abbildung $\Phi: P_N H \to Q_N H$ mit $\mathcal{M} = \mathrm{graph}(\Phi)$.
\item[(G3)] \textbf{Uniformit\"at:} Die Konstanten $C$, $\alpha$ und die Lipschitz-Konstante von $\Phi$ h\"angen nur von den physikalischen Parametern $(\nu, f, \Omega)$ ab, nicht von den Anfangsdaten.
\end{enumerate}
\end{axiom}

\begin{proposition}[\"Aquivalenzen von G-NS]\label{prop:g-equivalences}
Die folgenden Aussagen sind \"aquivalent:
\begin{enumerate}
\item[(i)] Axiom~\ref{ax:g-ns} gilt (inertiale Mannigfaltigkeit existiert).
\item[(ii)] \textbf{Spektrall\"uckenbedingung:} Es existiert $N$ mit $\lambda_{N+1} - \lambda_N \geq C_{\mathrm{NL}} \cdot \lambda_N^{1/2}$, wobei $\lambda_j$ die Eigenwerte des Stokes-Operators sind und $C_{\mathrm{NL}}$ die Nichtlinearit\"at kontrolliert.
\item[(iii)] \textbf{Starkes Squeezing:} Die Squeezing-Eigenschaft (Proposition~\ref{prop:squeezing-tll}) gilt gleichm\"a\ss ig f\"ur alle L\"osungen auf dem Attraktor.
\item[(iv)] \textbf{Assumption G} (wie in diesem Paper definiert): die N\"achste-Punkt-Selektion $v(t) \in \mathcal{A}$ ist absolut stetig mit integrierbaren Gronwall-Koeffizienten.
\end{enumerate}
Die Implikation (i)$\Rightarrow$(iv) ist klassisch; (iv)$\Rightarrow$(i) folgt aus der Glattheit des Attraktors und der durch die BV-Regularit\"at erzwungenen Endlichdimensionalit\"at.
\end{proposition}

\begin{remark}[Der Tail-Scanner]\label{rem:tail-scanner}
Um zu lokalisieren, wo G-NS bricht, werden drei diagnostische Gr\"o\ss en untersucht:
\begin{enumerate}
\item \textbf{Spektraldichte:} Die Stokes-Eigenwerte erf\"ullen $\lambda_j \sim c\, j^{2/3}$ auf $\mathbb{T}^3$. Die Spektrall\"uckenbedingung (ii) erfordert $\lambda_{N+1} - \lambda_N \gtrsim \lambda_N^{1/2}$, was versagt, wenn Eigenwertcluster zu dicht werden. F\"ur 3D ist diese L\"ucke grenzwertig: $\lambda_{N+1} - \lambda_N \sim N^{-1/3}$, w\"ahrend $\lambda_N^{1/2} \sim N^{1/3}$.
\item \textbf{Energietransfer:} Falls Energie dauerhaft im Hochfrequenz-Tail verbleibt ($\sum_{j > N} E_j \not\to 0$ exponentiell), versagt die Slaving-Bedingung (G2). DNS-Evidenz legt exponentiellen Tail-Abfall f\"ur 3D NS bei moderaten Reynoldszahlen nahe.
\item \textbf{Defektes Slaving:} Falls $\Phi$ nur H\"older-stetig ist (nicht Lipschitz), degradiert die BV-Regularit\"at zu H\"older-Regularit\"at -- unzureichend f\"ur Assumption~G, aber ausreichend f\"ur approximative inertiale Mannigfaltigkeiten (Foias--Sell--Temam).
\end{enumerate}
\end{remark}

\begin{remark}[No-Go-Mechanismen]\label{rem:g-nogo}
G-NS versagt, falls:
\begin{enumerate}
\item \textbf{Spektrale H\"aufung:} Eigenwertl\"ucken schlie\ss en sich zu schnell, sodass die lineare Dissipation die nichtlineare Kopplung nicht dominieren kann. Dies ist der Grenzfall f\"ur 3D NS ($\lambda_j \sim j^{2/3}$, verglichen mit $\lambda_j \sim j$ f\"ur 2D, wo inertiale Mannigfaltigkeiten existieren).
\item \textbf{Intermittierende Hochfrequenz-Bursts:} Kurzlebige, aber wiederkehrende Energieexkursionen im Tail zerst\"oren die exponentielle Attraktion -- ein klarer dynamischer Mechanismus gegen G-NS.
\item \textbf{Attraktor-Rauheit:} Falls der Attraktor fraktale Dimension nahe der Einbettungsdimension hat, kann keine glatte Mannigfaltigkeit ihn enthalten -- approximative inertiale Mannigfaltigkeiten sind das bestm\"ogliche Resultat.
\end{enumerate}
\end{remark}

\begin{remark}[Die Tail-Dominanz-Ungleichung]\label{rem:tail-dominance}
Der strukturelle Gehalt von Axiom~\ref{ax:g-ns} reduziert sich auf eine einzige quantitative Ungleichung: \emph{lineare Dissipation muss die nichtlineare Kopplung im Hochfrequenz-Tail dominieren}. Explizit: Für die Stokes-Eigenwerte $\lambda_j \sim c\, j^{2/3}$ auf $\mathbb{T}^3$ erfordert die Slaving-Bedingung (G2)
\[
  \nu \lambda_{N+1} \geq C_{\mathrm{NL}} \cdot \sup_{u \in \mathcal{A}} \|B(P_N u, Q_N u)\|,
\]
wobei $B$ die Bilinearform der Navier--Stokes-Nichtlinearität ist. Die Grenzskalierung $\lambda_{N+1} - \lambda_N \sim N^{-1/3}$ gegenüber $\lambda_N^{1/2} \sim N^{1/3}$ macht 3D zur \emph{kritischen Dimension} für inertiale Mannigfaltigkeiten: In 2D ($\lambda_j \sim j$, Lücke $\sim 1$) gilt Slaving; in 3D schließt sich die Lücke asymptotisch. Dies ist kein ,,fehlender Trick'', sondern eine \emph{strukturelle Grenzlinie} -- die Navier--Stokes-Nichtlinearität ist genau stark genug, um mit der Dissipation zu konkurrieren.
\end{remark}

\begin{lemma}[G-NS impliziert globale Regularit\"at]\label{lem:g-regularity}
Falls Axiom~\ref{ax:g-ns} gilt, dann ist jede Leray--Hopf-L\"osung glatt f\"ur $t > 0$ und der globale Attraktor $\mathcal{A}$ ist eine glatte kompakte endlichdimensionale Mannigfaltigkeit. Insbesondere ist Blow-up in endlicher Zeit ausgeschlossen.
\end{lemma}

\bigskip

\begin{remark}[Status von Assumption~G]\label{rem:assumption-G}
Assumption~G ist \emph{keine} Konsequenz der Standard-Attraktortheorie für die
3D-Navier--Stokes-Gleichungen. Die Existenz eines kompakten globalen Attraktors
$\mathcal{A} \subset H^1(\Omega)$ ist klassisch (vgl.\
\cite{FoiasManleyRosaTemam2001,Temam2001}), und jedes $v \in \mathcal{A}$ ist glatt
nach Proposition~\ref{prop:attractor-regularity}. Jedoch hängt die \emph{Zeitregularität}
der Abbildung $t \mapsto v(t) = \arg\min_{v \in \mathcal{A}} \|u(t) - v\|_{L^2}$
von der Geometrie von $\mathcal{A}$ und der Regularität der Trajektorie $u(t)$
auf nichttriviale Weise ab.

Falls $\mathcal{A}$ eine Inertiale Mannigfaltigkeit zulässt (d.h.\ eine endlichdimensionale,
Lipschitz-invariante Mannigfaltigkeit, die alle Orbits exponentiell anzieht), dann ist die
Nächste-Punkt-Projektion Lipschitz-stetig und (G1)--(G2) folgen.
Die Existenz inertialer Mannigfaltigkeiten für die 3D-Navier--Stokes-Gleichungen auf
$\mathbb{T}^3$ ist selbst ein offenes Problem, obwohl sie für mehrere
dissipative PDE-Klassen, die eine Spektrallückenbedingung erfüllen, etabliert ist
\cite{EdenFoiasNicolaenko1994}. Insbesondere ist für hyperviskose
Modifikationen $(-\Delta)^{\beta}$ mit $\beta \geq 5/4$ die Spektrallückenbedingung erfüllt und inertiale Mannigfaltigkeiten existieren~\cite{MalletParetSell1988},
sodass Assumption~G automatisch gilt. Der klassische Fall $\beta = 1$ bleibt
offen, gerade weil die Eigenwertlücken des Stokes-Operators in drei Dimensionen
zu langsam wachsen, um die für die Versklavung hochfrequenter Moden erforderliche
spektrale Trennung zu garantieren. Jüngere Resultate zur $H^2$-Regularität für Navier--Stokes
mit ``perfect slip''-Randbedingungen (2024/2025) deuten darauf hin, dass die
geometrische Einschränkung aufweichbar sein könnte, obwohl die hier relevante
no-slip-Situation ungelöst bleibt.
Zur log-Lipschitz-Regularität des Navier--Stokes-Semiflusses und deren
Implikationen für Lipschitz-artige Schranken auf Attraktorprojektionen
siehe~\cite{BaeCannone2016,MalletParetSell1988}.
Eine umfassende moderne Behandlung von Squeezing-Eigenschaften, Man\'e-Projektionen
und Attraktordimensionen für dissipative PDEs liefert
Miranville--Zelik~\cite{MiranvilleZelik2008}.

Diese Lücke zu schließen -- Assumption~G aus der intrinsischen Struktur des
Navier--Stokes-Attraktors zu beweisen -- ist das zentrale offene Problem des vorliegenden Ansatzes.
Siehe Abschnitt~5 (Einschränkungen) für eine detaillierte Diskussion und mögliche Strategien
mit log-Lipschitz-Regularität und Null-Lipschitz-Abweichung.

\medskip\noindent
\textbf{Resolventen-architektonische Umformulierung.}\;
Die spektrale Sektorzerlegung aus~\cite{GeigerRH2026c}
motiviert einen heuristischen Ansatz für Assumption~G. Führe \emph{Sektorlabels}
$\{S_j\}_{j=1}^{N}$ für die Galerkin-Moden von $v(t)$ ein, wobei jeder Sektor
$S_j$ einem Spektralband $[2^{j-1}, 2^j)$ des Stokes-Operators entspricht. Die Schlüsselerkenntnis ist, dass die Lipschitz-Konstante von
$t \mapsto v(t)$ von \emph{Sektoren-Kreuzungsereignissen} abhängt: wenn der
Minimierer von einem Attraktor-Ast zu einem anderen springt, entkoppeln die Moden in
verschiedenen Sektoren mit Raten, die durch die Spektrallücke des
Stokes-Operators kontrolliert werden. Umformulierung von Assumption~(G1) in dieser Sprache:
die projektive Familie $P_N(v(t)) = \sum_{j=1}^{N} \Pi_j v(t)$ (wobei
$\Pi_j$ auf Sektor $S_j$ projiziert) bleibt gleichmäßig beschränkt in einer
endlich-rangigen Approximation. Sektoren-Kreuzungsereignisse -- bei denen der Minimierer
zwischen Ästen wechselt, die durch verschiedene Sektorlabels getrennt sind -- erzwingen
höhere Lipschitz-Konstanten, aber Kompaktheit von $\mathcal{A}$ kontrolliert
die Häufigkeit solcher Kreuzungen. Falls der Attraktor ein
$\varepsilon$-Netz mit $N(\varepsilon) \leq C\,\varepsilon^{-d}$
Elementen zulässt (wobei $d$ die fraktale Dimension ist), dann liefert eine naive Schranke für
die Totalvariation von $v(t)$
$C\,N(\varepsilon)\,\varepsilon = C\,\varepsilon^{1-d}$, was
für $\varepsilon \to 0$ bei $d > 1$ divergiert.  Eine raffiniertere Analyse
müsste die zeitliche Kohaerenz der Minimierer-Trajektorie ausnutzen
(der Minimierer besucht nicht alle $\varepsilon$-Netz-Punkte, sondern nur jene
entlang des dynamischen Orbits), was eine bessere Schranke liefern könnte.
Dies bleibt eine offene Frage.

\medskip\noindent
\textbf{Morse-theoretische Umformulierung.}\;
Assumption~G kann in der Sprache der Dynamischen-Systeme-Geometrie
umformuliert werden: $H(u \mid \mathcal{A})$ definiert eine Ljapunow-Funktion auf dem
Phasenraum, und Assumption~G besagt, dass $\mathcal{A}$ ein
\emph{stratifizierter Morse-Komplex} für diese Funktion ist -- alle kritischen
Punkte von $H|_{\mathcal{A}}$ sind isoliert und der Gradientenfluss von~$H$
ist transversal zu den Strata. In dieser Umformulierung wird (G1) zu:
der Minimierer $v(t)$ kreuzt nicht zwischen Morse-Zellen
transversal zu~$\mathcal{A}$, und (G2) wird zu: die Morse-Indizes
(Dimensionen instabiler Mannigfaltigkeiten an kritischen Punkten) sind gleichmäßig
beschränkt. Die Conley-Index-Theorie bietet ein Rahmenwerk zum Beweis
solcher Eigenschaften ohne direkte PDE-Abschätzungen, was Assumption~G
auf eine topologische Aussage über die Zellstruktur des Attraktors reduziert.

\medskip\noindent
\textbf{Helizitätskopplung.}\;
Eine Verletzung von (G2) -- unbeschränkte Gronwall-Koeffizienten durch
schnelles Minimierer-Umschalten -- würde unbeschränktes Wachstum der
Helizität $\mathcal{H} = \int u \cdot \omega\,dx$ entlang der Trajektorie erzeugen.
Dies liegt daran, dass Minimierer-Sprünge zwischen topologisch verschiedenen
Attraktor-Ästen Rotationsumordnungen beinhalten, die
Helizität injizieren. Falls man zeigen kann, dass das Helizitätswachstum durch die
Enstrophie-Dissipation beschränkt ist (eine Konsequenz der NS-Gleichungen in 3D:
$\frac{d}{dt}\mathcal{H} = -2\nu\int\omega\cdot\nabla\times\omega\,dx$),
dann wird Assumption~G äquivalent zu ``keine Helizitätskaskade'':
der Helizitätsfluss über Skalen bleibt beschränkt. Dies würde das
bedingte Rahmenwerk mit einer physikalisch messbaren und
experimentell falsifizierbaren Bedingung verbinden.
\end{remark}

% ------------------------------------------------------------------
\subsection*{Hauptsatz}
% ------------------------------------------------------------------
\begin{theorem}[Bedingter Ausschluss von Singularitäten in endlicher Zeit]
\label{thm:main}
Sei $\Omega = \mathbb{T}^3$ und sei $u(t)$ eine Leray--Hopf-Lösung der
inkompressiblen Navier--Stokes-Gleichungen mit glatten Anfangsdaten
$u_0 \in H^1(\Omega)$. Sei $\mathcal{A} \subset H^1(\Omega)$ ein
Referenzattraktor, der Axiom~\ref{prop:attractor-regularity} erfüllt, und definiere
\[
    H\!\left(u\mid\mathcal{A}\right)
    \;=\; \inf_{v\in\mathcal{A}}\frac{1}{2}\int_\Omega |u-v|^2\,dx.
\]
Nach Axiom~\ref{prop:attractor-regularity} besitzt jedes $v \in \mathcal{A}$
die uniformen $H^2$- und $W^{1,\infty}$-Schranken, die unten benötigt werden.

\textbf{Falls Assumption~G \textup{(}oder die schwächere
Assumption~G$'$, Definition~\ref{ass:Gprime}\textup{)} gilt},
und falls zusätzlich die folgende \textbf{Dissipations-Dominanz-Bedingung~(D)} erfüllt ist: es existiert ein Blow-up-Szenario, in dem
\begingroup
\renewcommand{\theequation}{D}
\renewcommand{\theHequation}{D}
\begin{equation}\label{eq:diss-dom}
  \frac{\nu}{2}\int_0^T \|\nabla u\|_{L^2}^2\,dt
  \;>\;
  H(0)+
  \int_0^T \bigl(C_1\,H(t) + C_2\bigr)\,dt
  \quad \text{für ein } T < T^*,
\end{equation}
\endgroup
mit zusätzlichem BV/Stieltjes-Korrekturterm aus~\eqref{eq:H-integral-BV}
auf der rechten Seite, falls Assumption~G$'$ verwendet wird,
dann existiert keine endliche Blow-up-Zeit $T^* < \infty$. Insbesondere ist
$u(t)$ global glatt.
\end{theorem}

\begin{remark}[Die Dissipationslücke]\label{rem:diss-gap}
Bedingung~\eqref{eq:diss-dom} ist das eigentliche Hindernis, das
das vorliegende bedingte Rahmenwerk von einem unbedingten Beweis trennt.
Die Leray-Energieidentität garantiert
$\int_0^{T^*}\|\nabla u\|^2 < \infty$ (selbst unter Blow-up), sodass
die linke Seite von~\eqref{eq:diss-dom} \emph{endlich} ist.
Die rechte Seite (Gronwall-Kosten) ist unter Assumption~G ebenfalls endlich.
Ob die Dissipation die Gronwall-Kosten dominiert -- d.h.\ ob~\eqref{eq:diss-dom} gilt -- hängt von den \emph{relativen Raten} der beiden Terme nahe $T^*$ ab.

Drei Strategien zur Etablierung von~\eqref{eq:diss-dom} sind identifiziert:
\begin{enumerate}
  \item[\textup{(i)}]  Das Funktional $H$ auf $H^1$-Ebene heben,
    wobei $\|\nabla u\|_{L^2}^2$ durch $\|\nabla\omega\|_{L^2}^2$
    ersetzt wird (das unter Blow-up divergieren kann).
  \item[\textup{(ii)}]  Quantitative Blow-up-Raten ausnutzen, um
    zu zeigen, dass $\|\nabla u(t)\|^2$ nahe $T^*$ schnell genug wächst,
    um den Gronwall-Faktor zu übertreffen.
  \item[\textup{(iii)}]  Serrin-artige Integrierbarkeitskriterien
    zur unabhängigen Kontrolle der Gronwall-Kosten verwenden.
\end{enumerate}
Jede davon ist ein offenes Problem. Die vorliegende Arbeit etabliert
das bedingte Rahmenwerk; die Schließung der Dissipationslücke wird
auf zukünftige Arbeiten verschoben.

\medskip\noindent
\textbf{Quantitative Analyse von Bedingung~(D).}\;
Unter Assumption~G erfüllen die Gronwall-Koeffizienten
$C_1 \le C_1^* := 2\sup_{v\in\mathcal{A}}\|\nabla v\|_{L^\infty}+1$
und $C_2 \le C_2^*$ (beide Konstanten unabhängig von~$t$).
Nach dem Standard-Gronwall-Lemma gilt
$H(t) \le [H(0)+C_2^*\,T^*]\,e^{C_1^*\,T^*}$
auf $[0,T^*)$. Daher erfüllt die rechte Seite
von~\eqref{eq:diss-dom}
$\int_0^T(C_1\,H+C_2)\,dt
\le C_1^*[H(0)+C_2^*T^*]\,T^*\,e^{C_1^*T^*} + C_2^*T^*$,
während die linke Seite durch die Leray-Energieidentität beschränkt ist:
$\frac{\nu}{2}\int_0^{T^*}\|\nabla u\|^2\,dt
\le \frac{1}{2}\|u_0\|^2_{L^2}
+ \|f\|_{L^2}\,\|u\|_{L^\infty_tL^2}\,T^*$.
Bedingung~(D) wird damit zu einem quantitativen Vergleich zwischen
dem Anfangsenergiebudget und dem exponentiell verstärkten Attraktorabstand.
Insbesondere ist für kleines $T^*$ (schneller Blow-up)
$e^{C_1^*T^*} \approx 1$ und die Bedingung am leichtesten
erfüllt; für grosses~$T^*$ (langsamer Blow-up nahe dem Attraktor)
wächst der exponentielle Gronwall-Faktor und die Bedingung wird
restriktiver. Das $H^1$-Lift-Programm
(Abschnitt~\ref{sec:H1-lift}) umgeht diese Einschränkung vollständig:
die Enstrophie-Dissipation $\int\|\Delta u\|^2$ hat
keine a-priori-Oberschranke, muss aber dennoch die Bihari-/Gronwall-Kosten
dominieren; diese Dominanz ist eine zusätzliche offene Bedingung.
Die $H^1$-Ebene sollte daher als das \emph{natürliche
Sobolev-Niveau} für diese Obstruktionsstrategie verstanden werden, nicht als optionale
Erweiterung.
\end{remark}

\begin{remark}[$L^2$-Ebene-Obstruktionsheuristik]\label{thm:L2-obstruction}
Der vorangehende Budgetvergleich ist eine heuristische strukturelle
Warnung, kein Theorem mit bewiesenem universellem Schwellenwert. Ein
gültiges $T_{\mathrm{crit}}$-Theorem würde eine unabhängige Untergrenze
für die Gronwall-Kosten oder eine äquivalente quantitative Beschreibung
des Abstands der Trajektorie vom Attraktor verlangen. Streng bewiesen ist
hier nur die bedingte Implikation von Satz~\ref{thm:main}: Wenn die
geschärfte Bedingung~(D) gilt, ist endlicher Blow-up ausgeschlossen. Das
$L^2$-Energiebudget erklärt, warum~(D) als substanzielle Zusatzannahme
behandelt werden muss.
\end{remark}

\begin{remark}[Strukturelle Einschränkung auf $L^2$-Ebene]\label{cor:L2-limitation}
Jede Obstruktionsstrategie basierend auf
$H(u|\mathcal{A}) = \inf_{v \in \mathcal{A}}\frac{1}{2}\|u-v\|_{L^2}^2$
mit Gronwall-artiger Energiekontrolle muss das endliche
Leray-Dissipationsbudget berücksichtigen. Das macht den Satz nicht leer,
bedeutet aber, dass Dissipations-Dominanz nicht aus BKM oder der
Leray-Energieidentität allein folgt. Das natürliche Sobolev-Niveau für
den nächsten Obstruktionsversuch ist $H^1$, wo die
Enstrophie-Dissipation $\int\|\Delta u\|^2\,dt$ keine
a-priori-Oberschranke besitzt.
\end{remark}

\begin{proof}
Das Argument verläuft in drei Schritten.

\medskip
\noindent\textbf{Schritt 1 (BKM-Kriterium).} Nehme zum Widerspruch an, dass eine endliche
Blow-up-Zeit $T^* < \infty$ existiert. Nach dem Beale--Kato--Majda-Satz
\cite{BealeKatoMajda1984} ist dies äquivalent zu
\begin{equation}\label{eq:BKM}
    \int_0^{T^*}\|\omega(t)\|_{L^\infty}\,dt = \infty.
\end{equation}

\medskip
\noindent\textbf{Schritt 2 (Punktweiser Enstrophie-Blow-up).} Lemma~\ref{lem:enstrophy}
(angewandt auf das hypothetische Blow-up-Szenario \eqref{eq:BKM}) liefert
\begin{equation}\label{eq:diss-inf}
    \limsup_{t \nearrow T^*}\|\nabla u(t)\|_{L^2}^2 = \infty.
\end{equation}
Hinweis: Das \emph{Zeitintegral} $\int_0^{T^*}\|\nabla u\|^2\,dt$
bleibt endlich nach der Leray-Energieidentität
(Bemerkung~\ref{rem:divergence-landscape}).

\medskip
\noindent\textbf{Schritt 3 (Widerspruch via $H$-Schranke).}
Nach Lemma~\ref{lem:entropy-bound} erfüllt für jede messbare Auswahl
$v(t)\in\mathcal{A}$ (die existiert wegen Kompaktheit von $\mathcal{A}$ in $L^2$ und
Messbarkeit von $t\mapsto H(u(t)\mid\mathcal{A})$; siehe Bemerkung~\ref{rem:selection}
unten) das Funktional $E_v(t)=\tfrac12\|u(t)-v(t)\|_{L^2}^2$
\[
    \frac{d}{dt}E_v(t)
    \;\le\; -\frac{\nu}{2}\|\nabla u(t)\|_{L^2}^2 + C_1\,E_v(t) + C_2.
\]
Wahl von $v(t)$ als minimierende Auswahl (d.h.\ $E_v(t) = H(u(t)\mid\mathcal{A})$)
und Grenzübergang $\varepsilon\to 0$ in der $\varepsilon$-approximativen Auswahl
liefert dieselbe Ungleichung für $H(t):=H(u(t)\mid\mathcal{A})$ im schwachen Sinne.
Integration über $[0,T]$ für $T < T^*$:
\begin{equation}\label{eq:H-integral}
    H(T) - H(0)
    \;\le\;
    -\frac{\nu}{2}\int_0^T\|\nabla u\|_{L^2}^2\,dt
    + \int_0^T\!\big(C_1\,H(t) + C_2\big)\,dt.
\end{equation}
Die rechte Seite von \eqref{eq:H-integral} hat zwei konkurrierende Terme:
die Dissipation $-\tfrac{\nu}{2}\int_0^T\|\nabla u\|^2\,dt$
(negativ, treibt $H$ nach unten) und die Gronwall-Kosten
$\int_0^T(C_1\,H + C_2)\,dt$ (positiv, erlauben $H$ zu wachsen).
Beide Terme sind \emph{endlich} für jedes $T < T^*$ (das Dissipationsintegral
ist beschränkt durch die Leray-Energieidentität; die Gronwall-Kosten
sind beschränkt unter Assumption~G).

Falls Bedingung~\eqref{eq:diss-dom} gilt -- d.h.\ die Dissipation
dominiert das Anfangsbudget plus die Gronwall-Kosten für ein $T$
hinreichend nahe an $T^*$ -- dann gilt $H(T) < 0$, im Widerspruch zu
$H \ge 0$.

Folglich existiert unter Assumptions~G (oder~G$'$) und~(D) kein endliches $T^*$
und $u(t)$ ist global glatt.

\medskip
\noindent\textbf{BV-Modifikation (Assumption~G$'$).}\;
Unter der schwächeren Assumption~G$'$ ist der Minimierer $v(t)$ BV statt
AC. Die einzige Stelle, an der absolute Stetigkeit von $v$ in den
Beweis eingeht, ist Lemma~\ref{lem:entropy-bound}: Der Koeffizient
$C_2(t)$ enthält den Term $\tfrac{1}{2}\|f - \partial_t v\|_{L^2}^2$,
der erfordert, dass $\partial_t v$ punktweise f.u.\ existiert. Unter~G$'$
ist $\partial_t v$ ein Maß (keine Funktion), und dieser Term ist
schlecht definiert.

\emph{Auflösung.}\;
Kehre zur Rechnung zurück, die zu Lemma~\ref{lem:entropy-bound} führt.
Der Term $\langle w, \partial_t v\rangle$ entsteht aus
$\frac{\mathrm{d}}{\mathrm{d}t}\tfrac{1}{2}\|u-v\|^2
= \langle w, \partial_t u\rangle - \langle w, \partial_t v\rangle$.
Unter~G$'$ ersetze den zweiten Summanden durch sein Stieltjes-Integral:
$\int_0^T \langle w(t), \mathrm{d}v(t)\rangle$, das beschränkt ist durch
\begin{equation}\label{eq:stieltjes-bound}
  \biggl|\int_0^T \langle w(t),\, \mathrm{d}v(t)\rangle\biggr|
  \;\le\;
  \sup_{t \in [0,T]} \|w(t)\|_{L^2}
  \;\cdot\;
  \mathrm{Var}_{L^2}(v;\, [0,T]).
\end{equation}
Die modifizierte Energieungleichung wird zu
\begin{equation}\label{eq:H-integral-BV}
\begin{aligned}
  H(T) \;\le\; H(0)
  &-\frac{\nu}{2}\int_0^T \|\nabla u\|_{L^2}^2\,\mathrm{d}t
   + \int_0^T \widetilde{C}_1\,H(t)\,\mathrm{d}t
   + \widetilde{C}_2\,T \\
  &+ \sqrt{2\,\sup H}\,
   \mathrm{Var}_{L^2}(v;\,[0,T]).
\end{aligned}
\end{equation}
wobei $\widetilde{C}_1 = 2\sup_{v \in \mathcal{A}}
\|\nabla v\|_{L^\infty} + 1$ und $\widetilde{C}_2 =
\tfrac{1}{2}\sup_{v \in \mathcal{A}}\|(v\cdot\nabla)v\|_{L^2}^2
+ \tfrac{1}{2}\|f\|_{L^2}^2 + \tfrac{\nu}{2}\sup_{v \in \mathcal{A}}
\|\nabla v\|_{L^2}^2$ beide endlich sind nach
Proposition~\ref{prop:attractor-regularity}. Man beachte, dass $\widetilde{C}_2$
$\|\partial_t v\|$ \emph{nicht enthält} -- der Beitrag der Minimierer-Geschwindigkeit
wurde in die Stieltjes-Schranke~\eqref{eq:stieltjes-bound} absorbiert.

Unter Blow-up ($T \nearrow T^* < \infty$):
Lemma~\ref{lem:enstrophy} gibt die \emph{punktweise} Divergenz
$\|\nabla u(t)\|_{L^2}^2 \to \infty$ für $t \nearrow T^*$, aber die
Leray-Energieidentität sichert die Endlichkeit des \emph{Zeitintegrals}
$\int_0^{T^*}\|\nabla u\|_{L^2}^2\,dt$
(Bemerkung~\ref{rem:divergence-landscape}).

\emph{Beschränktheit von $\sup H$.}\;
Das implizite Auftreten von $\sup_{[0,T]} H$ in~\eqref{eq:H-integral-BV}
erfordert eine Begründung. Nach der Young-Ungleichung gilt
$\sqrt{2\sup H}\cdot\mathrm{Var}(v)
\le \varepsilon\,\sup_{[0,T]} H + \mathrm{Var}(v)^2/(2\varepsilon)$
für beliebiges $\varepsilon > 0$.
Mit $\varepsilon = 1$ und $M(T) := \sup_{[0,T]} H$ liefert
die Ungleichung~\eqref{eq:H-integral-BV} für jedes $t \le T$:
$H(t) \le H(0) + \int_0^t \widetilde{C}_1\,H\,d\tau
+ (\widetilde{C}_2+1)\,T^* + \mathrm{Var}(v)^2/2$.
Das Standard-Gronwall-Lemma ergibt dann
$M(T) \le [H(0) + (\widetilde{C}_2+1)\,T^*
+ \mathrm{Var}(v)^2/2]\,e^{\widetilde{C}_1\,T^*}$,
was auf $[0,T^*)$ endlich ist unter~G$'$.

Alle anderen Terme auf der rechten Seite von
\eqref{eq:H-integral-BV} sind ebenfalls endlich
($\widetilde{C}_1, \widetilde{C}_2$ beschränkt;
$\mathrm{Var}(v) < \infty$ nach~G$'$).
Daher ist der Widerspruch $H(T) < 0$ \emph{nicht} automatisch aus der
BV-Struktur allein: er erfordert zusätzlich, dass der Dissipationsterm
das Anfangsbudget, die Gronwall-Kosten und die Stieltjes-Korrektur dominiert,
d.h.\ die BV-Version von Bedingung~\eqref{eq:diss-dom}.
Unter Assumptions~G$'$ und~(D) wird die rechte Seite
von~\eqref{eq:H-integral-BV} negativ für $T$ hinreichend nahe
an~$T^*$, was $H(T) < 0$ ergibt -- Widerspruch.
\end{proof}

\begin{remark}[Abhängigkeiten und Fehlermodi]\label{rem:dependencies}
\mbox{}

\begin{center}
\begin{tabularx}{\textwidth}{@{}lYlY@{}}
\hline
\textbf{Label} & \textbf{Aussage} & \textbf{Typ} & \textbf{Fehlermodus} \\
\hline
Prop.~\ref{prop:attractor-regularity} &
  $\mathcal{A}$ kompakt, glatt &
  \textbf{Axiom} & Offen für klassische 3D NS \\
Lem.~\ref{lem:BS} &
  Biot--Savart-Abschätzung &
  Satz & Keiner \\
Lem.~\ref{lem:enstrophy} &
  Enstrophie-Bilanz &
  Satz & Keiner \\
Lem.~\ref{lem:entropy-bound} &
  $H$-Schranke (AC-Version) &
  Satz & Erfordert $\partial_t v \in L^2$ \\
Gl.~\eqref{eq:H-integral-BV} &
  $H$-Schranke (BV-Version) &
  \textbf{Satz} & Erfordert $\mathrm{Var}(v) < \infty$ \\
AGC2 &
  Exponentielles Tracking &
  \textbf{Annahme} & Nicht automatisch; braucht Squeezing/exponentielle Attraktion \\
\textbf{AGC1} &
  \textbf{Positiver Reach} &
  \textbf{Offen} &
  Scheitert falls $\mathcal{A}$ Spitzen hat \\
Lem.~\ref{lem:AGC-implies-G} &
  AGC $\Rightarrow$ G &
  Satz & Keiner (Lipschitz-1) \\
\hline
\end{tabularx}
\end{center}

\medskip\noindent
Die logische Kette ist:
$\textup{Attraktor-Paket} + \textup{AGC1 (offen)} + \textup{AGC2 (angenommen)}
\;\Longrightarrow\;
\textup{G}
\;\Longrightarrow\;
\textup{G}'
\;\Longrightarrow\;
\textup{Satz~\ref{thm:main}}$.
Die Implikationsglieder sind rigoros, sobald die genannten Inputs
angenommen werden. Der offene Inhalt konzentriert sich nicht in einer
einzigen Eigenschaft: Für das klassische 3D-Problem braucht man weiter
das Attraktor-Regularitätspaket, positiven Reach (oder BV-Auswahlersatz),
effektives Tracking in die Reach-Röhre und Bedingung~(D).
\end{remark}

\begin{remark}[Universelles Resolventenmuster]\label{rem:universal-resolvent}
Das bedingte Regularitätsresultat von Satz~\ref{thm:main} ist eine Instanz des
\emph{Second-Order Resolvent Dominance Pattern}, das über alle fünf
Probleme im FST-Programm identifiziert wurde~\cite{FST-Unified2026}. Das Muster manifestiert sich als
dreistufige Hierarchie:

\noindent
Im Navier--Stokes-Fall ist die führende Ordnung die Leray--Hopf-Energieungleichung,
die Beschränktheit von $\|u\|_{L^2}$ liefert, aber keine Regularität.
Die Attraktorprojektion erster Ordnung identifiziert den nächsten glatten
Referenzzustand $v(t) \in \mathcal{A}$, aber dies allein kann
$\|\nabla u\|_{L^\infty}$ nicht kontrollieren. Die Regularitätslücke schließt sich erst in zweiter
Ordnung, durch die $H$-Metrik-Kontraktion: die viskose Dissipation
treibt den Attraktorabstand $H(u \mid \mathcal{A})$ in Richtung eines
Widerspruchs zu seiner Nichtnegativität.  Diese Hierarchie wird im
breiteren Kontext des FST-Programms in~\cite{FST-Unified2026} diskutiert.
\end{remark}

\begin{remark}[Spektralstruktur der Minimierer-Instabilität]
\label{rem:resolvent-minimiser}
Die instabilen Richtungen des Abstandsfunktionals $E_v = \frac{1}{2}\|u-v\|^2$
am Minimierer $v(t) \in \mathcal{A}$ -- jene, entlang derer
$\frac{d}{dt}E_v > 0$ -- entsprechen den Ljapunow-instabilen Moden des
Semiflusses bei~$v(t)$ und bilden einen endlichdimensionalen Unterraum
der Dimension höchstens $\dim(\mathcal{A})$. Die Zeitableitung $\frac{d}{dt}E_v$
zerfällt in einen dissipativen Teil (kontrolliert durch $-\nu\|\nabla u\|^2$)
und einen Gronwall-Teil (kontrolliert durch $C_1 E_v + C_2$). Kontraktion in der
$H$-Metrik ist somit ein Phänomen \emph{zweiter Ordnung}: sie stabilisiert die
Trajektorie durch den viskosen Term, nachdem die Attraktorprojektion erster Ordnung
die Instabilität auf ein endlichdimensionales Residuum lokalisiert hat.
\end{remark}

\begin{remark}[Messbare Auswahl]
\label{rem:selection}
Die Existenz einer messbaren Abbildung $t\mapsto v(t)\in\mathcal{A}$, die
$H(u(t)\mid\mathcal{A}) = \tfrac12\|u(t)-v(t)\|_{L^2}^2$ realisiert, ist eine Standardkonsequenz des Kuratowski--Ryll-Nardzewski-Satzes über messbare Auswahl, angewandt
auf die kompaktwertige, unterhalb-halbstetige Multifunktion
$t\mapsto \arg\min_{v\in\mathcal{A}}\|u(t)-v\|_{L^2}$; siehe z.B.\ \cite{Wagner1977}, Theorem~4.1.
\end{remark}

\begin{remark}[Explizite Abhängigkeit der Konstanten]
\label{rem:constants}
Die im Beweis auftretenden Konstanten hängen wie folgt von den Problemdaten ab.
\begin{itemize}
    \item $C_{\mathrm{BS}}$: hängt nur von $s > 3/2$ und der Geometrie von
          $\mathbb{T}^3$ ab (Lemma~\ref{lem:BS}).
    \item $C_1 = 2\sup_{v\in\mathcal{A}}\|\nabla v\|_{L^\infty} + 1$: endlich,
          da $\mathcal{A}$ ein kompakter Attraktor ist, der alle Orbits absorbiert,
          also $\sup_{v\in\mathcal{A}}\|v\|_{W^{1,\infty}} < \infty$.
    \item Unter Assumption~G (AC-Version):
          $C_2 = \tfrac12\sup_{v\in\mathcal{A}}\|(v\cdot\nabla)v\|_{L^2}^2
          + \tfrac12\|f - \partial_t v\|_{L^2}^2
          + \tfrac{\nu}{2}\sup_{v\in\mathcal{A}}\|\nabla v\|_{L^2}^2$.
          Unter der schwächeren Assumption~G$'$ (BV-Version) ist die
          Minimierer-Geschwindigkeit $\partial_t v$ ein Maß statt einer Funktion,
          und der entsprechende Term wird durch die Stieltjes-Schranke
          \eqref{eq:stieltjes-bound} ersetzt; die modifizierte Konstante
          $\widetilde{C}_2 = \tfrac12\sup_{v\in\mathcal{A}}\|(v\cdot\nabla)v\|_{L^2}^2
          + \tfrac12\|f\|_{L^2}^2 +
          \tfrac{\nu}{2}\sup_{v\in\mathcal{A}}\|\nabla v\|_{L^2}^2$
          enthält $\|\partial_t v\|$ nicht.
          Beide Versionen sind endlich durch Attraktorkompaktheit und die Annahme
          $f \in L^\infty_t H^1_x$. Insbesondere hängen sowohl $C_1$ als auch $C_2$
          (bzw.\ $\widetilde{C}_2$)
          von der Antriebsamplitude $\|f\|_{L^\infty H^1}$, der
          Viskosität $\nu$ und dem Attraktorradius ab (letzterer bestimmt durch
          $\nu$ und $\|f\|$), sind aber unabhängig von $T$ und $u_0$.
\end{itemize}
\end{remark}

\section{Diskussion}

\subsection{Zusammenhang mit partiellen Regularitätsresultaten}

Der Satz von Caffarelli--Kohn--Nirenberg \cite{CaffarelliKohnNirenberg1982} zeigt, dass das eindimensionale parabolische Hausdorff-Maß der singulären Menge jeder geeigneten schwachen Lösung Null ist. Obwohl dies eine tiefgreifende geometrische Einschränkung ist, schließt sie Blow-up in endlicher Zeit nicht aus: sie begrenzt lediglich die Größe der potentiellen singulären Menge. Ebenso liefern die Serrin-artigen Regularitätskriterien \cite{Leray1934,ConstantinFoias1988} hinreichende Bedingungen für Glattheit (z.B.\ $u \in L^p_t L^q_x$ mit $2/p + 3/q \le 1$), aber die a-priori-Verifikation dieser Bedingungen erfordert Informationen, die aus den Anfangsdaten allein nicht verfügbar sind.

Unser Ansatz unterscheidet sich strukturell von beiden Paradigmen. Anstatt externe Integrierbarkeitsbedingungen an die Lösung zu stellen oder die Größe potentieller Singularitäten abzuschätzen, nutzen wir die interne dissipative Struktur der Gleichungen aus: Eine endliche maximale starke Existenzzeit erzwingt nach dem Fortsetzungskriterium die \emph{punktweise} Divergenz der Enstrophie $\limsup_{t\nearrow T^*}\|\nabla u(t)\|_{L^2}^2 = \infty$; das BKM-Kriterium liefert dazu ein paralleles $L^\infty$-Vortizitätssignal. Diese Divergenz erzeugt eine Spannung mit der Nichtnegativität des Attraktorabstandsfunktionals $H(u \mid \mathcal{A})$. Unter der zusätzlichen Dissipations-Dominanz-Bedingung~(D) -- die verlangt, dass die kumulative viskose Dissipation die Gronwall-Kosten übersteigt -- wird diese Spannung zum Widerspruch: $H < 0$. Das Argument operiert innerhalb des Energie-Dissipations-Rahmenwerks, mit zwei zentralen verbleibenden Bedingungen: einer Regularitätsannahme an die Attraktorprojektion (Assumption~G) und der Dissipations-Dominanz-Bedingung~(D).

\subsection{Strukturelle Vorteile des attraktorbasierten Ansatzes}

Klassische Energiemethoden für die Navier--Stokes-Gleichungen (siehe z.B.\ \cite{ConstantinFoias1988,Temam2001}) erzeugen typischerweise Gronwall-artige Ungleichungen der Form
\[
    \frac{d}{dt}\|u\|_{H^1}^2 \le C\,\|u\|_{H^1}^2\,\|\nabla u\|_{L^\infty} - \nu\|\nabla u\|_{H^1}^2,
\]
wobei der Wirbelstreckungsterm auf der rechten Seite mit der viskosen Dissipation konkurriert. Die fundamentale Schwierigkeit besteht darin, dass in drei Dimensionen keine a-priori-Schranke für $\|\nabla u\|_{L^\infty}$ verfügbar ist und die Gronwall-Abschätzung lediglich eine doppelt-exponentielle Schranke liefert, die in endlicher Zeit explodieren kann.

Der attraktorbasierte Ansatz umgeht diese Sackgasse durch Verschiebung des Bezugssystems: Statt $\|u\|_{H^1}$ direkt abzuschätzen, misst man die \emph{Abweichung} $w = u - v$ von einem glatten Referenzzustand $v \in \mathcal{A}$. Der Gronwall-Koeffizient $C_1 = 2\|\nabla v\|_{L^\infty} + 1$ ist nun \emph{beschränkt} (da $v$ auf dem kompakten Attraktor liegt), anstelle des unkontrollierten $\|\nabla u\|_{L^\infty}$ des klassischen Ansatzes. Der Preis ist Assumption~G: man muss die zeitliche Regularität der Attraktorprojektion kontrollieren. Dies ist ein genuines anderes Problem als die Kontrolle von $\|\nabla u\|_{L^\infty}$ -- es ist eine geometrische Eigenschaft des Attraktors statt einer punktweisen Schranke an die Lösung.

\section[Navier-Stokes als Pattern-A-System]{Navier-Stokes-Regularität als Pattern-A-System}

\noindent\textbf{Disclaimer.}\;
Die folgenden Abschnitte ordnen das vorliegende Framework in den Kontext
des breiteren FST-Programms ein. Diese Verbindungen sind
\emph{spekulativ und programmatisch}: Sie skizzieren Analogien und
potentielle Transferpfade zwischen verschiedenen Problemen, aber keiner
von ihnen ist im vorliegenden Paper mathematisch untermauert.
Der nur am in sich geschlossenen bedingten Framework interessierte
Leser kann direkt zu Abschnitt~\ref{sec:H1-lift} (dem $H^1$-Lift-Programm)
springen.

\medskip
Der Ausschluss von Blow-up in endlicher Zeit in den 3D-Navier-Stokes-Gleichungen wird hier als Realisierung der \textbf{Pattern-A-Stabilitätslogik} formuliert (siehe \cite{FST-Unified2026}). Dieser Ansatz ersetzt die Suche nach globalen a-priori-Schranken für $\|\nabla u\|_{L^\infty}$ durch die folgende strukturelle Architektur:

\begin{enumerate}[label=\alph*)]
    \item \textbf{Analytische Rang-Endlichkeit (Null-Tail):} Jedes hypothetische Blow-up-Szenario ist auf eine endlichdimensionale Mannigfaltigkeit von Instabilitätsmoden lokalisiert. Die Hochfrequenz-(UV-)Dissipationsskala wird strikt durch die Enstrophie-Bilanz \cite{Leray1934} kontrolliert, was einen analytischen Null-Tail sichert.
    \item \textbf{Attraktorprojektion als physikalische Eichung:} Die ``endliche Negativitaet'' (potentielle Singularitätsrichtungen) wird neutralisiert durch Projektion der aktiven Strömung auf den globalen Leray--Hopf-Attraktor $\mathcal{A}$. Diese Projektion fungiert als physikalische Eichung.
    \item \textbf{Eingeschränkte Positivität (H-Funktional):} Die Nichtnegativität des Abstandsfunktionals $H(u \mid \mathcal{A})$ wirkt als Eich-Stabilisator und stellt sicher, dass keine Trajektorie die ``Dissipationskosten'' eines Blow-ups bezahlen kann, ohne das thermodynamische Budget zu verletzen.
\end{enumerate}

Somit ist globale Regularität die Manifestation spektraler Stabilität unter der Eich-Beschränkung der Attraktorstruktur.
Das Dissipationsintegral wirkt als irreversible ``Kosten'', die jede Blow-up-Trajektorie bezahlen muss, während der Attraktorabstand ein nichtnegativer ``Etat'' ist, der nicht überschritten werden kann. Dieser Mechanismus ist analog zur Rolle der Entropieproduktion in der Nichtgleichgewichtsthermodynamik: Die dissipative Struktur des Systems selbst verbietet Konfigurationen, die unendliche Entropieproduktion in endlicher Zeit erfordern.

\subsection{Einschränkungen und offene Fragen}

Mehrere Aspekte der vorliegenden Arbeit laden zu weiterer Untersuchung ein:

\begin{itemize}
    \item \textbf{Assumption~G (die Gronwall-Lücke).} Die bedeutendste Einschränkung dieser Arbeit ist die Abhängigkeit von Assumption~G. Die Differentialungleichung für $H(u\mid\mathcal{A})$ im Beweis von Satz~\ref{thm:main} erfordert, dass das minimierende Attraktorelement $v(t)$ mit hinreichender Zeitregularität evolviert. Während jedes feste $v \in \mathcal{A}$ räumlich glatt ist, kann die \emph{Trajektorie} der Nächste-Punkt-Projektionen $t \mapsto v(t)$ Sprünge aufweisen, falls der Attraktor komplizierte Geometrie hat (z.B.\ Spitzen oder Selbstschnitte in der $L^2$-Topologie). Drei Strategien zum Schließen dieser Lücke erscheinen am vielversprechendsten: (i)~Beweis der Existenz einer Inertialen Mannigfaltigkeit für 3D-Navier--Stokes auf $\mathbb{T}^3$, was die Projektion Lipschitz-stetig machen würde; (ii)~Etablierung von log-Lipschitz-Regularität der Projektion über Quetschungseigenschaften des Semiflusses (vgl.\ Bae und Cannone); (iii)~Ausnutzung von Null-Lipschitz-Abweichung auf dem Attraktor, um zu zeigen, dass hochfrequente Moden deterministisch von niederfrequenten Moden versklavt werden, was $v(t)$ automatisch glatt macht.
    Unabhängige Unterstützung für den attraktorgeometrischen Ansatz kommt aus jüngsten Arbeiten zu fraktalen Sparseness-Bedingungen: Grujic und Xu~\cite{GrujicXu2025} zeigen, dass Singularitäten topologisch ausgeschlossen werden, wenn die Regionen intensiver Vortizität hinreichend dünn (im Sinne fraktaler Masse) sind -- die erste Skalierungsreduktion der NS-Superkritikalitätslücke seit den 1960er Jahren. Ihre Sparseness-Bedingung ist strukturell kompatibel mit dem hier angenommenen endlichdimensionalen Attraktor.

    \item \textbf{Gebietsgeometrie.} Unser Beweis stützt sich auf die periodische Situation $\Omega = \mathbb{T}^3$, die sowohl die Existenz eines kompakten globalen Attraktors \cite{FoiasManleyRosaTemam2001} als auch das Verschwinden von Randtermen bei partieller Integration sichert. Erweiterung auf beschränkte Gebiete mit No-Slip-Randbedingungen führt Grenzschichteffekte ein und erfordert Attraktortheorie in $H^1_0(\Omega)$, die verfügbar ist (vgl.\ \cite{Temam2001}), aber technisch aufwendiger. Der Ganzraum-Fall $\Omega = \mathbb{R}^3$ ist delikater, da der globale Attraktor im üblichen Sinne nicht kompakt sein muss; man müsste stattdessen mit Trajektorienattraktoren oder Pullback-Attraktoren arbeiten.

    \item \textbf{Äußerer Antrieb.} Wir haben glatten, zeitunabhängigen (oder zumindest gleichmäßig beschränkten) Antrieb $f \in L^\infty_t H^1_x$ angenommen. Für rauen oder singulären Antrieb können die gleichmäßigen Schranken für $C_1$ und $C_2$ versagen, und die Attraktorstruktur selbst kann sich ändern. Die Erweiterung des Ansatzes auf stochastischen Antrieb (wie in den stochastischen Navier--Stokes-Gleichungen) ist eine interessante Richtung, die ein Rahmenwerk zufälliger Attraktoren erfordern würde.

    \item \textbf{Quantitative Schranken.} Das vorliegende Framework ist bedingt: Es beweist globale Regularität nicht unbedingt, sondern reduziert das Problem auf Assumptions~G und~(D). Sollten beide Bedingungen in Zukunft verifiziert werden, wären die resultierenden Schranken für $\|\nabla u\|_{L^\infty}$ nichtkonstruktiv und hängen vom Attraktorradius ab, der seinerseits von den Abschätzungen der absorbierenden Kugel des Navier--Stokes-Semiflusses abhängt. Herleitung expliziter, quantitativer Regularitätsschranken aus dem Attraktorabstand würde präzisere Informationen über die Attraktorgeometrie erfordern, die in drei Dimensionen weitgehend unbekannt bleibt.

    \item \textbf{Regularität der messbaren Auswahl.} Der Beweis verwendet eine messbare Auswahl $t \mapsto v(t) \in \mathcal{A}$, um das nächste Attraktorelement zu verfolgen. Während Messbarkeit für die Integralabschätzungen genügt, würde höhere Regularität dieser Auswahl (z.B.\ absolute Stetigkeit) die Differentialungleichung verstärken und potentiell schärfere Abklingraten für $H(u \mid \mathcal{A})$ liefern.
\end{itemize}

\begin{remark}[Aufbrechen der Regularitäts-Attraktor-Zirkularität]
\label{rem:circularity-discussion}
Der Standardzugang zu globalen Attraktoren in der dissipativen
PDE-Theorie (Temam~\cite{Temam2001}, Zelik~\cite{Zelik2014}) setzt
Glattheit der Lösungen voraus, um Endlich-Dimensionalität des
Attraktors zu beweisen. Unser Framework kehrt diese logische
Richtung um: Wir setzen Endlich-Dimensionalität des Attraktors
voraus (die empirisch gut gestützt und für den schwachen
Leray--Hopf-Attraktor etabliert ist; vgl.\
Bemerkung~\ref{rem:circularity}) und leiten
Regularitätskonsequenzen ab.

Die scheinbare Zirkularität -- Regularität wird für den Attraktor
benötigt, und der Attraktor wird für Regularität benötigt -- wird
durch die \emph{thermodynamische Obstruktion} aufgebrochen: Falls
Blow-up in endlicher Zeit einträte, würde das Freie-Energie-Funktional
$H(u \mid \mathcal{A})$ gegen~$-\infty$ divergieren, im Widerspruch
zur dissipativen Schranke $H(u \mid \mathcal{A}) \ge 0$.
Die Nichtnegativität von~$H$ ist eine Konsequenz der Definition
($H$ ist ein quadrierter Abstand), nicht der Regularität.

Dies ist \emph{kein} Beweis unbedingter Regularität, sondern eine
\emph{strukturelle Reduktion}: Globale Regularität folgt, WENN
der Attraktor die vorausgesetzten Eigenschaften besitzt
(Endlich-Dimensionalität, Kompaktheit, Glattheit der Elemente).
Jede dieser Eigenschaften ist unabhängig von der Regularitätsfrage
etabliert -- sie folgen aus der schwachen Attraktortheorie und dem
Bootstrap-Argument für Elemente auf dem Attraktor
(Proposition~\ref{prop:attractor-regularity}). Die verbleibende
Lücke ist Assumption~G, die die \emph{zeitliche Regularität} der
Projektion betrifft, nicht die räumliche Regularität der
Attraktorelemente.
\end{remark}

\begin{remark}[Zerlegung von Assumption G: unabhängige Komponenten]
\label{rem:g-decomposition}
Um dem Einwand zu begegnen, Assumption~G könne zirkulär sein (implizit Regularität voraussetzen), zerlegen wir sie in drei logisch unabhängige Komponenten:

\begin{enumerate}
\item[\textbf{(G$_1$)}] \textbf{Attraktorexistenz und Kompaktheit.} Der Referenzattraktor $\mathcal{A}$ existiert als kompakte Teilmenge der Arbeitstopologie und erfüllt die uniformen Schranken aus Axiom~\ref{prop:attractor-regularity}. \emph{Status:} Bedingt für klassische 3D-Navier--Stokes; schwache/Trajektorienattraktoren existieren in schwächeren Settings, aber die hier benötigte $H^1$-Kompaktheit und uniforme Glattheit sind ohne zusätzliche Regularitätsstruktur nicht bekannt.

\item[\textbf{(G$_2$)}] \textbf{Projektionsregularität.} Die Nächste-Punkt-Projektion $\pi_\mathcal{A}: H^1(\Omega) \to \mathcal{A}$ ist log-Lipschitz (oder lässt eine BV-Selektion zu) in einer tubularen Umgebung von~$\mathcal{A}$. \emph{Status:} Offen. Dies ist eine \textbf{geometrische} Eigenschaft des Attraktors, keine dynamische Regularitätsaussage. Sie kann prinzipiell aus der metrischen Struktur des Attraktors (Hausdorff-Dimension, Tangentenkegel-Regularität, Assouad-Dimension) verifiziert werden, ohne auf die Regularität der Lösungen Bezug zu nehmen.

\item[\textbf{(G$_3$)}] \textbf{Gronwall-Integrierbarkeit.} Die zeitabhängigen Gronwall-Koeffizienten $C_1(t), C_2(t)$ in der Differentialungleichung für $H(u|\mathcal{A})$ erfüllen $\int_0^T C_i(t)\,dt < \infty$. \emph{Status:} Folgt aus Standard-Energieabschätzungen und der Absorbing-Ball-Eigenschaft, unter der Annahme, dass die Lösung in $H^1$ verbleibt. Unter Blow-up divergiert die Enstrophie $\|\nabla u\|^2$, was $C_1$ unbeschränkt macht -- aber dies ist genau der Widerspruchsmechanismus von Satz~\ref{thm:main}.
\end{enumerate}

Der Zirkularitätseinwand betrifft ausschließlich (G$_2$): Beweist man zuerst Regularität, ist der Attraktor glatt und (G$_2$) trivial. Unser Framework bricht diese Zirkularität, indem es (G$_2$) als \emph{geometrische Hypothese} über den Attraktor behandelt (verifizierbar aus seinen metrischen Eigenschaften) statt als dynamische Regularitätsannahme. Die Brücke von (G$_1$) zu (G$_2$) ist der Inhalt der Begleitarbeit zur Log-Distanz-Integrierbarkeit~\cite{Geiger2026-LDI}.
\end{remark}

\begin{remark}[Assumption G als Schnittstellenvertrag]
\label{rem:g-interface-contract}
Assumption~G lässt sich in einen minimalen \emph{Schnittstellenvertrag} zerlegen, der die exakten Regularitätsdaten spezifiziert, die von der Attraktorgeometrie benötigt werden:
\begin{description}
\item[Input 1 (Attraktorregularität):] $\sup_{v \in \mathcal{A}} \|v\|_{H^2} < \infty$. \emph{Status: bekannt} (folgt aus der Glattheit des globalen Attraktors).
\item[Input 2 (Selektionsregularität):] Die Nächste-Punkt-Abbildung $t \mapsto v(t) \in \mathcal{A}$ erfüllt $v \in BV([0,T]; H^1)$. \emph{Status: offen} -- dies ist die zentrale Lücke, adressiert durch das TLL/LDI-Framework der Begleitarbeit~\cite{Geiger2026-LDI}.
\item[Input 3 (Gronwall-Integrierbarkeit):] Die Gronwall-Koeffizienten $C_1(t)$ erfüllen $\int_0^{T^*} C_1(t)\,dt < \infty$. \emph{Status: folgt aus Input~2} (BV-Selektion impliziert lokal integrierbare Ableitungsschranken).
\end{description}
Diese Zerlegung verdeutlicht, dass der gesamte offene Inhalt von Assumption~G sich auf eine einzige Frage reduziert: \emph{Hat die Nächste-Punkt-Projektion auf den Attraktor beschränkte Variation entlang von Leray--Hopf-Lösungen?} Die Begleitarbeit liefert eine hinreichende Bedingung (TLL + LDI) und validiert sie numerisch an drei Attraktoren.
\end{remark}

\begin{remark}[Versagensmodi von Assumption G]
\label{rem:failure-modes}
Falls Assumption~G versagt, hat das Versagen spezifische geometrische Signaturen:
\begin{enumerate}
\item \emph{Unendliche Sprungakkumulation:} Der Minimierer $v(t) = \arg\min_{v \in \mathcal{A}} \|u(t) - v\|^2$ erfährt unendlich viele diskontinuierliche Sprünge in jeder Umgebung von $T^*$, wobei die Sprunggrößen nicht summierbar sind. Geometrisch entspricht dies dem Durchqueren unendlich vieler ``Grate'' der Voronoi-Tessellation von $\mathcal{A}$ durch die Trajektorie $u(t)$ in endlicher Zeit.
\item \emph{Cantor-artiges Schalten:} Die Menge der Schaltzeitpunkte $\{t : v(t^-) \neq v(t^+)\}$ hat positives Maß oder Cantor-artige Struktur, was BV-Regularität verhindert. Dies würde erfordern, dass der Attraktor auf den relevanten Skalen eine pathologisch fragmentierte Voronoi-Struktur besitzt.
\item \emph{Spiralakkumulation:} Die Trajektorie spiralt um eine Spitze oder Selbstdurchdringung von $\mathcal{A}$, wodurch die Projektion mit zunehmender Frequenz oszilliert. Dies ist der physikalisch plausibelste Versagensmodus und würde darauf hindeuten, dass die Attraktorgeometrie nahe potenzieller Blow-up-Punkte komplexer ist, als aktuelle Dimensionsabschätzungen nahelegen.
\end{enumerate}
Das TLL/LDI-Framework der Begleitarbeit~\cite{Geiger2026-LDI} schuetzt gegen alle drei Modi: TLL begrenzt die Oszillationsrate logarithmisch, während LDI sicherstellt, dass die kumulativen Oszillationskosten endlich bleiben.
\end{remark}

\begin{remark}[Reach-Lemma über Rückwärtsstabilität]
\label{rem:reach-lemma}
Falls der Attraktor $A$ positiven Reach $\mathrm{reach}(A) > 0$ besitzt (d.\,h.\ die Nächste-Punkt-Projektion ist einwertig in einer tubularen Umgebung), dann impliziert Rückwärtsstabilität (Lösungen auf $A$ sind Lipschitz-stetig in Rückwärtszeit) kombiniert mit Injektivität endlichdimensionaler Projektionen $P_N: A \to \mathbb{R}^N$ die Abschätzung $\mathrm{reach}(A) \geq c/\|P_N\|_{\mathrm{op}}$. Dies würde Assumption~G (Lipschitz-Regularität) als Konsequenz der Endlich-Dimensionalität etablieren und die in den Einschränkungen genannte Zirkularität aufbrechen.
\end{remark}

\begin{remark}[Relaxation von Bedingung (D) zu limsup-Divergenz]
\label{rem:condition-d-limsup}
Bedingung~(D) kann von punktweiser Dominanz zu einer masstheoretischen Version relaxiert werden: Anstatt $\|\nabla u(t)\|^2 / H(u(t)|A) \to \infty$ für $t \to T^*$ für alle Blow-up-Trajektorien zu fordern, genügt es zu verlangen
\[
\limsup_{t \to T^*} \frac{\|\nabla u(t)\|^2}{H(u(t)|A)} = +\infty.
\]
Diese schwächere Bedingung erzwingt weiterhin den Widerspruch in Satz~3.1 (die Freie-Energie-Barriere kann nicht aufrechterhalten werden, wenn Dissipation intermittierend dominiert) und ist aus PDE-Perspektive natürlicher: Enstrophie-Blow-up muss nicht monoton sein, nur unbeschränkt.
\end{remark}

\begin{remark}[Bedingung D: beobachtbare Checkliste]
\label{rem:condition-d-checklist}
Die folgenden beobachtbaren Größen in DNS oder numerischer Simulation würden Evidenz für oder gegen Bedingung~D liefern:
\begin{enumerate}
\item \emph{Enstrophie-Wachstumsrate:} Falls $\|\nabla u(t)\|_{L^2}^2 / H(u(t)|\mathcal{A})$ für $t \to T^*$ unbeschränkt wächst, gilt Bedingung~D. Man überwache das Verhältnis $\|\nabla u\|^2 / \|u - v\|^2$ entlang nahezu-singulärer Trajektorien.
\item \emph{Palinstrophie:} Das $H^1$-Niveau-Analogon erfordert $\int_0^{T^*} \|\Delta u\|_{L^2}^2\,dt = \infty$. Man verfolge die kumulative Palinstrophie in Simulationen mit adaptiver Auflösung nahe potenziellem Blow-up.
\item \emph{Attraktorabstands-Zerfall:} Falls $H(u(t)|\mathcal{A}) \to 0$ schneller als $\|\nabla u\|^{-2}$, ist Bedingung~D erfüllt. Man vergleiche die Raten numerisch.
\end{enumerate}
Eine Falsifikation von Bedingung~D (beschränkt bleibendes Enstrophie-Verhältnis während Blow-up) würde darauf hindeuten, dass der Attraktorabstands-Ansatz einer Modifikation bedarf, würde aber nicht die Regularität klären -- es würde lediglich das Problem auf die Suche nach einem anderen Ljapunow-artigen Funktional verschieben.
\end{remark}

\subsection{Eine spieltheoretische Perspektive}

Die Regularitätsfrage lässt eine natürliche Umformulierung als Minimax-Problem zu. Man betrachte zwei abstrakte ``Spieler'', die die Dynamik steuern: Spieler~$D$ (Dissipation), dessen Strategie viskose Dämpfung mit Rate $\nu\|\nabla u\|_{L^2}^2$ ist, und Spieler~$N$ (Nichtlinearität), dessen Strategie der konvektive Energietransfer über $(u\cdot\nabla)u$ ist. Definiere den \emph{Spielwert} auf Modenebene als
\begin{equation}
    \mathcal{V}(t) \;:=\; \inf_{\text{Moden}} \bigl[\text{Dissipationsrate} - \text{nichtlineare Transferrate}\bigr].
\end{equation}
In dieser Sprache ist globale Regularität äquivalent zu $\mathcal{V}(t) \ge 0$ für alle $t$: Dissipation dominiert die Nichtlinearität gleichmäßig über alle Fourier-Moden. Der Attraktorabstand $H(u\mid\mathcal{A})$ spielt die Rolle eines endlichen ``Budgets'', das dem System zur Verfügung steht, während das Dissipationsintegral $\int_0^T \nu\|\nabla u\|_{L^2}^2\,dt$ die kumulativen Kosten darstellt, die jede Blow-up-Trajektorie bezahlen muss. Da das Budget nichtnegativ ist, reduziert sich die Frage darauf, ob die kumulativen Dissipationskosten die Gronwall-Kosten nahe der Blow-up-Zeit übersteigen -- genau die Dissipations-Dominanz-Bedingung~(D) von Satz~\ref{thm:main}.

Dieser Standpunkt verbindet sich mit der Legendre--Fenchel-Dualitaet in der konvexen Analysis: Interpretiert man die Dissipation als konvexes ``Entropieproduktions''-Funktional $S$ und den nichtlinearen Transfer als sein konjugiertes ``Freie-Energie-Kosten''-Funktional $F$, so ist die Regularitätsbedingung $S \ge F$ das PDE-Analogon des zweiten Hauptsatzes der Thermodynamik -- das System kann aus der Nichtlinearität nicht mehr Arbeit extrahieren, als die Dissipation erlaubt.

\begin{remark}[Mean-Field-Game-Perspektive]
Die Fokker--Planck-Komponente von Mean-Field-Game-Systemen
(Lasry--Lions~\cite{LasryLions2007}) hat die Form
$\partial_t m - \nu\Delta m + \nabla\cdot(m\alpha) = 0$,
strukturell identisch mit Advektions-Diffusion mit Drift~$\alpha$.
Ein Nash-Gleichgewicht -- in dem unendlich viele ``Fluidteilchen'' gegen
die kollektive Dichte optimieren -- entspricht einer selbstkonsistenten
Balance von Transport und Diffusion. Ob die Wohlgestelltheitstheorie
von MFG-Systemen die Navier--Stokes-Regularitätsfrage informieren kann,
ist eine offene Richtung.
\end{remark}

\begin{remark}[Mean-Field-Game-Formulierung im Wasserstein-Raum]
\label{rem:mfg-wasserstein}
Jüngere Entwicklungen in der Mean-Field-Game-Theorie (MFG) formulieren das gekoppelte HJB--Fokker-Planck-System als PDE auf dem Wasserstein-Raum der Wahrscheinlichkeitsmasse (Cardaliaguet--Delarue--Lasry--Lions, 2019). Die NS-Freie-Energie $H(u|A)$ lässt sich in diesem MFG-Rahmen als Wertfunktion interpretieren: Der ,,Agent`` (Fluidpaket) minimiert laufende Kosten (Enstrophie) unter Advektions-Diffusions-Dynamik, während die ,,Population`` (Geschwindigkeitsfeld) über die Navier--Stokes-Gleichungen evolviert. Der Attraktor $A$ entspricht dem MFG-Gleichgewicht. Diese Perspektive legt nahe, dass Regularität des NS-Attraktors äquivalent zur Wohlgestelltheit der zugehörigen Mastergleichung im Wasserstein-Raum ist.
\end{remark}

\begin{remark}[Attraktor--Gronwall-Kontrolle und Nash-Frustration]
\label{rem:agc-nash-frustration}
Die Attraktor--Gronwall-Kontrolle (AGC) aus
Abschnitt~\ref{sec:assumption-G} lässt eine aufschlussreiche Parallele
zum \emph{Nash-Frustrations}-Konzept aus FST-III~\cite{FST-Unified2026}
zu. Spieltheoretisch entsteht Nash-Frustration, wenn kein reines
Nash-Gleichgewicht existiert und das System in ein
Mischstrategien-Regime gezwungen wird.
Der Navier--Stokes-Attraktor $\mathcal{A}$ spielt eine strukturell
analoge Rolle: Wenn glatte Einzellösungen (``reine Strategien'')
nicht global bestehen können -- d.\,h.\ wenn das Gronwall-Budget
erschöpft ist -- wird die Dynamik auf den Attraktor gezwungen, der
als ``gemischte Strategie'' fungiert und das frustrierte Energiebudget
absorbiert. Die AGC-Bedingung quantifiziert die Kosten dieses
übergangs: Ist die Attraktorprojektion $t \mapsto v(t)$ hinreichend
regulär (positiver Reach oder BV-Selektion), kann das System die
Dissipationskosten über die Attraktormannigfaltigkeit umverteilen,
ohne eine Blow-up-Singularität zu erzeugen. In dieser Lesart ist
globale Regularität die Aussage, dass das Navier--Stokes-System
stets eine ``gemischte Strategie'' (den Attraktor) findet, um die
Nash-Frustration konkurrierender Wirbelstreckung und viskoser
Dämpfung aufzulösen.
\end{remark}

\begin{remark}[Konzentration des Attraktorabstandsfunktionals]
\label{rem:cumulants}
Im dissipativen Regime verhindert die Enstrophie-Bilanz
(Lemma~\ref{lem:enstrophy}) grosse Auslenkungen von $H$ und
unterstützt die Konzentration von Trajektorien um~$\mathcal{A}$.
Falls $H$ sub-Gaussische exponentielle Momente entlang der Trajektorie zulässt,
sind alle höheren Kumulanten varianz-dominiert, was quantitative
Tail-Schranken für $\{H > \ell\}$ liefert. Dies beeinflusst nicht das deterministische
Argument von Satz~\ref{thm:main}, motiviert aber probabilistische
Erweiterungen; siehe
Kuksin--Shirikyan~\cite{KuksinShirikyan2012} und
Hairer--Mattingly~\cite{HairerMattingly2006} für verwandte Invarianten-Maß-Ansätze.
\end{remark}

\begin{remark}[Pattern-A-Master-Stabilität -- problemübergreifende Struktur]
Diese Arbeit instanziiert das Pattern-A-Master-Stabilitätsprinzip aus dem vereinheitlichten Framework \cite{FST-Unified2026}: Strikte Konvexität der renormierten freien Energie $F$, kombiniert mit einem neutralen führenden Resolventenzweig und Dominanz zweiter Ordnung, impliziert eine Spektrallücke $\Delta > 0$ für den zugehörigen Transferoperator. Die spezifische Instanziierung im vorliegenden Kontext ist: $\Delta = $ exponentielle Relaxationsrate zum globalen Attraktor.
\end{remark}

\begin{remark}[Wasserstein-Konvexität und Talagrand-Schranke]
\label{rem:talagrand}
Das Freie-Energie-Funktional $H(u|A)$ ist gleichmäßig konvex in der Wasserstein-2-Metrik mit Konvexitätskonstante $\lambda_H > 0$ (vererbt von der quadratischen Struktur der $H^1$-Norm). Nach dem Satz von Otto--Villani impliziert dies eine Talagrand-Ungleichung $W_2(\mu, \mu^*) \leq \sqrt{2 D_{\mathrm{KL}}(\mu \| \mu^*) / \lambda_H}$, die eine quantitative Schranke dafür liefert, wie weit eine beliebige Trajektorienverteilung vom Attraktormass abweichen kann. Die Poincar\'e-Konstante erfüllt $\kappa \geq \lambda_H$, was eine $\beta$-unabhängige untere Schranke ergibt -- im Gegensatz zur Holley--Stroock-Schranke, die exponentiell degeneriert. Diese strukturelle Beobachtung, erstmals im Yang--Mills-Begleitpaper formalisiert, gilt universell für alle FST-Instanziierungen.
\end{remark}

\subsection{\texorpdfstring{Transfer von Yang--Mills: Poincar\'e--Dobrushin auf Galerkin-trunkierte Navier--Stokes}{Transfer von Yang-Mills: Poincare-Dobrushin auf Galerkin-trunkierte Navier-Stokes}}
\label{sec:ym-ns-transfer}

Die Galerkin-Trunkierung der Navier--Stokes-Gleichungen bei Modenzahl $N$ ist strukturell analog zur Gitter-Yang--Mills-Regularisierung bei Gitterabstand $a \sim 1/N$:

\begin{center}
\begin{tabularx}{\textwidth}{Y|Y|Y}
& \textbf{Yang--Mills (Gitter)} & \textbf{Navier--Stokes (Galerkin)} \\
\hline
Freiheitsgrade & Linkvariablen $U_l \in G$ & Fourier-Moden $\hat{u}_k \in \mathbb{C}^3$ \\
Wechselwirkungsreichweite & Plaquettes (nächste Nachbarn) & Triadische Wechselwirkungen $k + p + q = 0$ \\
Regularisierung & Gitterabstand $a$ & Moden-Cutoff $N$ \\
Kontinuumslimes & $a \to 0$ ($\beta \to \infty$) & $N \to \infty$ \\
Freie Energie & $F[\mu] = \langle S \rangle + D_{\mathrm{KL}}$ & $H(u|A) = \frac{1}{2}\|u - u^*\|_{H^1}^2$ \\
\end{tabularx}
\end{center}

\begin{remark}[Dobrushin-Einflussmatrix für Galerkin-Moden]
\label{rem:dobrushin-galerkin}
Definiere die Dobrushin-Interdependenzmatrix $C = (c_{kp})$ für Galerkin-Moden durch
\[
c_{kp} = \sup \| \mu_k(\cdot | \omega) - \mu_k(\cdot | \omega') \|_{\mathrm{TV}},
\]
wobei das Supremum über Konfigurationen genommen wird, die sich nur in Mode $p$ unterscheiden, und $\mu_k(\cdot | \omega)$ die bedingte Verteilung von Mode $k$ gegeben alle anderen Moden ist. Für die Navier--Stokes-Nichtlinearität $B(u, u)$ ist der Eintrag $c_{kp}$ genau dann von Null verschieden, wenn ein $q$ mit $k + p + q = 0$ existiert (triadische Wechselwirkung). Der Spektralradius erfüllt
\[
\rho(C) \leq \sup_k \sum_{p: \exists q, k+p+q=0} c_{kp} \leq C_{\mathrm{tri}} \cdot \frac{\|u\|_{H^1}}{N^{1/2}},
\]
wobei $C_{\mathrm{tri}}$ eine universelle Konstante ist und das $N^{-1/2}$-Abklingen den abnehmenden Einfluss hoher Moden widerspiegelt. Die Dobrushin-Eindeutigkeitsbedingung $\rho(C) < 1$ gilt für hinreichend grosses $N$ (d.\,h.\ hinreichend feine Galerkin-Trunkierung) und liefert eine Poincar\'e-Ungleichung für die trunkierte Dynamik.

Der zentrale Unterschied zu Yang--Mills besteht darin, dass die NS-Wechselwirkung im Fourier-Raum \emph{langreichweitig} ist (alle triadischen Wechselwirkungen tragen bei), während YM nächste-Nachbarn-artig auf dem Gitter ist. Allerdings stellt die Energiekaskade sicher, dass die effektive Wechselwirkungsstärke mit der Skalenseparation $|k - p|$ abklingt, was im Inertialbereich eine quasi-lokale Struktur wiederherstellt.
\end{remark}

\begin{remark}[Gap-Transport-Analogon für Galerkin-Trunkierungen]
\label{rem:gap-transport-galerkin}
Das Analogon der Yang--Mills-Gap-Transport-Vermutung im NS-Kontext lautet: Die Poincar\'e-Konstante des Galerkin-trunkierten Systems bei Cutoff $N$ erfüllt $\kappa(2N) \geq c \cdot \kappa(N)$ für ein universelles $c > 0$. Falls dies gilt, bleibt die exponentielle Relaxationsrate des Attraktors im Limes $N \to \infty$ erhalten, was einen unabhängigen Zugang zur Regularität liefert.
\end{remark}

\subsection{Offene Richtungen}

\begin{remark}[Inertialmannigfaltigkeits-Pfad zu Assumption~G]
\label{rem:inertial-manifold-path}
Die Theorie der Inertialmannigfaltigkeiten~\cite{FoiasSellTemam88} liefert
den direktesten Weg zur Schließung von Assumption~G\@. Auf einer
Inertialmannigfaltigkeit $\mathcal{M}$ sind die hochfrequenten Moden
$Q_N u$ deterministisch an die niederfrequenten Moden $P_N u$ über eine
Lipschitz-Abbildung $\Phi \colon P_N H \to Q_N H$ versklavt, sodass die
Nächster-Punkt-Projektion $t \mapsto v(t)$ die Lipschitz-Regularität
des endlichdimensionalen Flusses auf~$\mathcal{M}$ erbt. Die
Null-Lipschitz-Abweichung
$\mathrm{dev}_m(\mathcal{A}) = 0$~\cite{BaeCannone2016} würde dann
bedeuten, dass Assumption~G automatisch gilt: Die AC-Regularität von
$v(t)$ folgt aus der endlichdimensionalen ODE auf der Mannigfaltigkeit.

Im Begleitpaper zur Turbulenz~\cite{FST-TU2026} wird dieselbe
Maschinerie verwendet, um das K41-Spektrum als Sattelpunkt einer
skalenspezifischen KL-Divergenz auf der Inertialmannigfaltigkeit zu
formalisieren.
Die strukturelle Parallele lautet: Inertialmannigfaltigkeiten
reduzieren sowohl das NS-Regularitätsproblem als auch das
Turbulenz-Selektionsproblem auf endlichdimensionale variationelle
Fragen, in denen das KL-Framework rigoros ist.

Die Hauptobstruktion ist die \emph{Spektrallückenbedingung}:
Inertialmannigfaltigkeiten für 3D Navier--Stokes auf $\mathbb{T}^3$
erfordern eine Lücke zwischen aufeinanderfolgenden Eigenwerten des
Stokes-Operators, die für $\beta = 1$ (Standard-Viskosität) nicht
bekannt ist. Für hyperviskose NS ($\beta > 5/4$) existieren
Inertialmannigfaltigkeiten~\cite{EdenFoiasNicolaenko1994}; die
Erweiterung auf $\beta = 1$ bleibt ein offenes
Problem~\cite{Zelik2014}.
Neuere Arbeiten zum fraktionalen Navier--Stokes--Voigt-System~\cite{IlyinKalantarovZelik2025}
liefern scharfe Attraktordimensions-Abschätzungen via Lieb--Thirring-Ungleichungen,
die die klassischen nicht-fraktionalen Schranken verbessern und analoge Abschätzungen
im vorliegenden Rahmenwerk anleiten könnten.
\end{remark}

\begin{remark}[Log-Dirichlet-Quotienten in der KL-Joint-Minimierung]
\label{rem:log-dirichlet}
Eine Alternative zum Inertialmannigfaltigkeits-Ansatz ist die Integration
von \emph{Log-Dirichlet-Quotienten} für Differenzen von Lösungen auf
dem Attraktor in das KL-Joint-Minimierungs-Framework aus dem
Turbulenz-Begleitpaper~\cite{FST-TU2026}. Für zwei Lösungen $u, v$
auf dem Attraktor definiere den Log-Dirichlet-Quotienten
\[
  Q_{\mathrm{LD}}(t)
  = \frac{\|\nabla(u - v)\|_{L^2}^2}{\|u - v\|_{L^2}^2}.
\]
Ist $Q_{\mathrm{LD}}$ entlang von Trajektorien integrierbar, transformiert
sich die Gronwall-Abschätzung für $H(u \mid \mathcal{A})$ von
exponentiellem zu algebraischem Wachstum: Die Lyapunov-Exponenten werden
stabilisiert, und die kumulative Gronwall-Kosten
$\int_0^{T^*} Q_{\mathrm{LD}}(s)\, ds$ bleiben endlich, selbst wenn
$T^* \to T^*_{\max}$.

Dieses algebraische (anstatt exponentielle) Wachstum verbietet
topologisch Singularitäten in endlicher Zeit: Die freie Energie
$H(u \mid \mathcal{A})$ kann höchstens polynomiell abfallen, was mit
der für Blow-up erforderlichen divergenten Enstrophie unvereinbar ist.
\end{remark}

\begin{remark}[Helizitäts-Stratifizierung des Attraktors]
\label{rem:helicity-stratification}
Der globale Attraktor $\mathcal{A}$ lässt sich in
Helizitäts-Schichten zerlegen:
$\mathcal{A} = \bigcup_h \mathcal{A}_h$, wobei
$\mathcal{A}_h = \{ u \in \mathcal{A} : \mathcal{H}(u) = h \}$ und
$\mathcal{H}(u) = \int \mathbf{u} \cdot \boldsymbol{\omega}\, dx$
die Helizität ist~\cite{Moffatt69, ArnoldKhesin98}. Innerhalb jeder
Schicht $\mathcal{A}_h$ schränkt die Topologie die
Wirbelliniengeometrie ein und stabilisiert die
Nächster-Punkt-Projektion: Die \emph{stratifizierte} Projektion
$v_h(t) = \arg\min_{w \in \mathcal{A}_h} \|u(t) - w\|$ beschränkt
die Minimierersuche auf eine topologisch kohärente Teilmenge und
verbessert potentiell die Regularität von $t \mapsto v_h(t)$, womit
Assumption~G auf kontrollierter, topologierespektierender Auswahl
statt nacktem $\arg\min$ über den gesamten Attraktor begründet
werden könnte.
\end{remark}

\begin{remark}[Energie-Helizitäts-Ungleichungen als topologische
Koerzivität]
\label{rem:helicity-energy-bounds}
Für verknotete Wirbelkonfigurationen erzwingt die Topologie eine
untere Schranke für die Gradientenenergie~\cite{ArnoldKhesin98}:
\[
  \|\nabla u\|_{L^2}^2 \geq c_0 \, |\mathcal{H}(u)|^{3/4},
\]
wobei $c_0 > 0$ vom Gebiet und der Faddeev--Skyrme-artigen
Energieschranke abhängt. Diese topologische Koerzivität ergänzt
das $H(u \mid \mathcal{A})$-Budget um eine rein geometrische untere
Schranke: Selbst ohne Lösung der Gronwall-Ungleichung schränkt die
Topologie der Wirbellinien ein, wie schnell die Enstrophie wachsen
kann, und liefert potentiell eine unabhängige Obstruktion gegen
Blow-up.
\end{remark}

\begin{remark}[Knotenkomplexität als Attraktor-Regularisator]
\label{rem:knot-complexity}
Definiere eine funktionale Knotenkomplexität $K(u)$ auf dem Attraktor
über das Helizitätsdichte-Spektrum und die Verknüpfungsstatistik
der Wirbellinien~\cite{Moffatt69}. Ist $K$ auf $\mathcal{A}$
gleichmäßig beschränkt -- was angesichts der endlichen Dimension
und Kompaktheit des Attraktors plausibel ist --, so wird die
Minimierervariation $\|v(t_1) - v(t_2)\|$ durch den topologischen
Abstand zwischen den entsprechenden Knotenklassen kontrolliert. Dies
liefert einen geometrischen Mechanismus für die Zeitregularität der
Projektion $t \mapsto v(t)$.
\end{remark}

\begin{remark}[Helizitäts-Kaskade als Abwärtsbewegung im Knotenraum]
\label{rem:helicity-cascade-bridge}
Die Helizitäts-Stratifizierung des Attraktors
(Bemerkung~\ref{rem:helicity-stratification}) hat ein direktes
Gegenstück im Turbulenz-Begleitpaper~\cite{FST-TU2026}:
Die DFC-Bemerkung im Turbulenz-Begleitpaper argumentiert, dass Wirbelrekonnexionen dissipative
Schritte sind, die die topologische Komplexität des Wirbelfeldes
verringern und die Energiekaskade im Knotenraum ,,bergab`` machen.
Im NS-Regularitätskontext stabilisiert derselbe Mechanismus die
Attraktorprojektion: Übergänge zwischen Helizitäts-Schichten
$\mathcal{A}_h \to \mathcal{A}_{h'}$ sind dissipativ (Helizität
ist nahezu erhalten und nimmt durch Rekonnexion ab), sodass die
Projektion $t \mapsto v_h(t)$ innerhalb einer Schicht glatter
evolviert als ein unbeschränkter $\arg\min$. Diese Brücke --
Kaskadentopologie (TU) $\leftrightarrow$ Attraktortopologie (NS) --
legt nahe, dass jede Auflösung der Downhill Flux Condition in der
Turbulenz gleichzeitig die für Assumption~G relevante
Attraktorgeometrie einschränkt.
\end{remark}

\begin{remark}[Grujic--Xu-Sparseness als Dobrushin-Eindeutigkeit]
\label{rem:sparseness-dobrushin}
Die fraktale Sparseness-Bedingung von Grujic und
Xu~\cite{GrujicXu2025} -- Regionen intensiver Vortizität besetzen
eine Menge von fraktaler Dimension strikt kleiner als~$5/3$ in der
Raumzeit -- ist strukturell analog zur Dobrushin-Eindeutigkeitsbedingung
aus dem Yang--Mills-Begleitpaper und
Bemerkung~\ref{rem:dobrushin-galerkin}: Beide behaupten, dass
\emph{Wechselwirkungen hinreichend dünn} sind, um effektive Lokalität
zu gewährleisten. In YM erfüllt die Dobrushin-Einflussmatrix $C$ die
Bedingung $\rho(C) < 1$ (exponentielles Abklingen der Korrelationen);
in NS stellt die Sparseness-Bedingung sicher, dass
Wirbel-Streckungsereignisse zu dünn sind, um sich zu einer
Singularität aufzuakkumulieren. Das gemeinsame Prinzip lautet:
\emph{geometrische Sparseness $\Rightarrow$ effektive Lokalität
$\Rightarrow$ Eindeutigkeit/Regularität}.
\end{remark}

\begin{remark}[Endlichdimensionaler Attraktor impliziert
Vortizitäts-Sparseness]
\label{rem:attractor-sparseness}
Die endliche fraktale Dimension $d_f(\mathcal{A}) < \infty$ des globalen
Attraktors impliziert, dass On-Attraktor-Lösungen auf eine ,,dünne``
Teilmenge des Phasenraums beschränkt sind. Falls sich diese Dünnheit
auf die \emph{physikalisch-räumliche} Vortizitätsverteilung
überträgt -- d.\,h.\ falls Lösungen auf $\mathcal{A}$ die
Grujic--Xu-Sparseness-Bedingung automatisch erfüllen --, liefert die
Fraktale-Dimensions-Schranke einen \emph{unabhängigen Weg zu
Assumption~G}, der die Spektrallückenbarriere der
Inertialmannigfaltigkeiten vollständig umgeht. Die Schlüsselkonjektur
lautet:
\[
  d_f(\mathcal{A}) \leq D_0
  \;\;\Longrightarrow\;\;
  \dim_{\mathcal{H}}(\{x : |\omega(x)| > \lambda\})
  \leq 3 - c(\lambda, D_0)
  \quad\text{für alle } \lambda \gg 1,
\]
wobei $c > 0$ mit $\lambda$ wächst und Grujic--Xu-artige Sparseness
liefert. Die Etablierung dieser Implikation würde Assumption~G für
Attraktoren bekannter endlicher Dimension schließen.
\end{remark}

\begin{remark}[Fortgeschrittene Ansätze zur Attraktor-Regularität]
\label{rem:advanced-attractor}
Wir sammeln vier spekulative, aber konkrete Ansätze zur Etablierung von Attraktor-Regularität:
\begin{enumerate}
\item \emph{Multiplikative Oseledets-Kegel:} Invariante Konvexitätskegel im Tangentialraum $T_u A$ an Attraktorpunkten, vererbt von der monotonen Struktur der Navier--Stokes-Halbgruppe, könnten Spitzen-Glättung liefern: Ist der Kegelwinkel gleichmäßig beschränkt, sind Spitzen (und damit Singularitäten) ausgeschlossen.
\item \emph{Fraktionale Einbettung $H^{1+\varepsilon}$:} Log-Lipschitz-Regularität der Attraktorprojektion (etabliert im Begleitpaper LogDist) könnte für glatte Normalenbündel genügen und die volle Lipschitz-Anforderung von Assumption~G umgehen.
\item \emph{Dissipation als Radon-Maß:} Reinterpretiert man die Energiedissipation $\varepsilon(t) = \nu \|\nabla u\|^2$ als Radon-Maß auf $[0, T^*) \times \mathbb{R}^3$, wird Bedingung~(D) zu strikter Positivität des singulären Anteils $\varepsilon_{\mathrm{sing}}$ -- eine maßtheoretische Bedingung, die an die partielle Regularitätstheorie von Caffarelli--Kohn--Nirenberg anknüpft.
\item \emph{Kompensierte Kompaktheit:} Tartars Methode der kompensierten Kompaktheit, angewandt auf Vortizitätsoszillationen ($\omega = \nabla \times u$), könnte integrale Divergenz von $\|\nabla u\|^2$ nahe einer hypothetischen Singularität erzwingen und so Bedingung~(D) allein aus den strukturellen Eigenschaften der Euler-Nichtlinearität etablieren.
\end{enumerate}
\end{remark}

\begin{remark}[Spektrallücken überleben konische Singularitäten (Sturm)]
\label{rem:sturm-conical}
Sturm~\cite{Sturm2025} hat kürzlich gezeigt, dass Bakry--\'Emery-Krümmungsbedingungen und Spektrallücken-Abschätzungen auf Riemannschen Mannigfaltigkeiten mit konischen Singularitäten gültig bleiben, selbst wenn die Ricci-Krümmung als $\mathrm{Ric} \sim 1/d^2$ nahe der singulären Punkte degeneriert. Entscheidend ist, dass verallgemeinerte Hardy-Ungleichungen die fehlenden Krümmungs-Dimensions-Bedingungen ersetzen, und die Spektrallücke bemerkenswerte Robustheit gegenüber lokalen topologischen Defekten aufweist.

Dieses Resultat hat direkte Implikationen für den FST-Ansatz zur Navier--Stokes-Regularität: Selbst wenn der globale Attraktor $A$ den Reach $\mathrm{reach}(A) = 0$ hat (fraktaler Rand mit spitzenartigen Singularitäten), bleibt der thermodynamische Relaxationsmechanismus -- kontrolliert durch die Spektrallücke des zugehörigen Transferoperators -- intakt, sofern die Singularitäten konischen Typs sind. Dies unterstützt die schwächere Bedingung eines zeitgemittelten Reach $\mathrm{reach}_{\mathrm{avg}}(A) > 0$ (vgl.\ Bemerkung~\ref{rem:reach-lemma}): Die Moreau--Yosida-regularisierte Auswahl ,,sieht'' nur die gemittelte Geometrie, die glatt ist, selbst wenn punktweise Spitzen existieren.
\end{remark}

\begin{remark}[Duchon--Robert-Defektmaß als Brücke zwischen Turbulenz und Attraktor-Regularität]
\label{rem:duchon-robert}
Duchon und Robert~\cite{DuchonRobert2000} führten ein lokales
Energie-Defektmaß $D(u)$ für schwache Lösungen der inkompressiblen
Euler- und Navier--Stokes-Gleichungen ein.  Für eine schwache
Lösung~$u$ erfasst die Distribution $D(u)$ die lokale Rate der
inertialen Energiedissipation jenseits der viskosen Skala: Sie ist
der distributionelle Rest in der lokalen Energiebilanz
\[
  \partial_t \tfrac{|u|^2}{2}
  + \nabla \cdot \bigl[ \bigl(\tfrac{|u|^2}{2} + p\bigr) u \bigr]
  + \nu \|\nabla u\|^2
  = -D(u).
\]
Eine zentrale physikalische Einschränkung ist $D(u) \ge 0$ fast
überall für physikalisch zulässige Lösungen (Energie fliesst nur
zu kleineren Skalen).

Dieses Defektmaß verbindet sich auf drei Wegen direkt mit dem
FST-Attraktor-Framework:
\begin{enumerate}
\item \emph{Singularitäts-Lokalisierung auf dem Attraktor.}
Auf dem globalen Attraktor~$\mathcal{A}$ kodiert der Träger von
$D(u)$ die Singularitätsstruktur: Blow-up kann nur dort auftreten,
wo sich $D(u)$ konzentriert.  Falls der singuläre Träger von $D(u)$
Hausdorff-Dimension
$\dim_H(\mathrm{supp}_{\mathrm{sing}}\, D) < 1$ hat, ist die
BV-Auswahl $t \mapsto v(t)$ entlang des Attraktors ,,fast'' Lipschitz
in dem Sinne, dass ihre Totalvariation durch die
$(1-\dim_H)$-H\"older-Norm der Abstandsfunktion kontrolliert wird.
\item \emph{Kopplung an die BV-Selektionstheorie.}
Die Balanced-Viscosity-Auswahl aus dem LDI-Paper~\cite{Geiger2026-LDI}
erzeugt eine Trajektorie $v \in \mathrm{BV}$, die ein
Dissipationsfunktional minimiert.  Das Defektmaß $D(u)$ liefert
ein \emph{natürliches Gewicht} für diese Minimierung: Ersetzt man
den gleichmäßigen $L^2$-Abstand durch den $D$-gewichteten Abstand
$\int D(u(t))\, \|u(t) - a(t)\|^2\, dt$, konzentriert sich die
Auswahl dort, wo inertiale Dissipation aktiv ist, und verbessert die
Regularität in ruhigen Regionen.
\item \emph{Brücke zum Turbulenz-Begleitpaper.}
In~\cite{FST-TU2026} spielt das Defektmaß $D(u)$ die Rolle des
anomalen Dissipationsflusses in der turbulenten Kaskade.  Das
vorliegende Paper betrifft die Regularität der Projektion
auf~$\mathcal{A}$ (Assumption~G); das Turbulenz-Paper betrifft die
spektrale Verteilung von~$D(u)$.  Eine vereinheitlichte
Schließstrategie würde beides etablieren, indem gezeigt wird,
dass $D(u)$ ein Radon-Maß mit kontrolliertem
Hausdorff-dimensionalem Träger ist -- genau die Bedingung, unter
der die BV-Auswahl nahezu-Lipschitz-Regularität erbt
(vgl.~Bemerkung~\ref{rem:Drivas-VV-selection}).
\end{enumerate}
\end{remark}

% ==================================================================
\subsection{EDP-Konvergenz-Route zur Selektionsregularit\"at}\label{sec:gamma-edp-selection}
% ==================================================================

Die klassische punktweise Projektion $v(t) = \pi_A(u(t))$ ist strukturell unzureichend: BV-Regularit\"at ist eine zeitintegrierte Eigenschaft, aber punktweise Projektion ignoriert zeitliche Koh\"arenz. Wir ersetzen sie durch eine \emph{Trajektorienselektion}, die BV-Regularit\"at allein aus der Dissipation ableitet, ohne positive Reach, Inertialmannigfaltigkeiten oder Spektrallücken vorauszusetzen.

\subsubsection*{Schritt 1: Reskaliertes Trajektorienfunktional}

Sei $\mathcal{A}_T$ eine kompakte Menge von Trajektorien auf oder nahe dem Attraktor (im Sinne der Chepyzhov--Vishik-Trajektorienattraktoren \cite{ChepyzhovVishik2002}, die f\"ur 3D-Navier--Stokes ohne Eindeutigkeitsannahmen verf\"ugbar sind). F\"ur $\varepsilon > 0$ definiere das \emph{reskalierte} Funktional
\begin{equation}\label{eq:trajectory-functional}
  \mathcal{F}_\varepsilon[a] := \frac{1}{\varepsilon}\int_0^T \|u(t) - a(t)\|_H^2 \, dt \;+\; \mathrm{TV}(a),
\end{equation}
wobei $\mathrm{TV}(a) = \sup \sum_{i} \|a(t_{i+1}) - a(t_i)\|_H$ die Totalvariation ist. Man beachte die entscheidende Reskalierung: der Fidelity-Term tr\"agt $1/\varepsilon$, w\"ahrend TV ungewichtet bleibt. Die regularisierte Selektion ist
\begin{equation}\label{eq:reg-selection}
  v_\varepsilon := \operatorname{argmin}_{a \in \mathcal{A}_T} \mathcal{F}_\varepsilon[a].
\end{equation}

\subsubsection*{Schritt 2: A-priori-Schranken}

F\"ur jeden Minimierer $v_\varepsilon$ liefert der Vergleich mit einer festen Trajektorie $a_0 \in \mathcal{A}_T$
\begin{equation}\label{eq:apriori}
  \mathrm{TV}(v_\varepsilon) \;\leq\; \mathcal{F}_\varepsilon[v_\varepsilon] \;\leq\; \mathcal{F}_\varepsilon[a_0] \;=\; \frac{1}{\varepsilon}\int_0^T \|u - a_0\|^2\,dt + \mathrm{TV}(a_0).
\end{equation}
Da $u$ dem Attraktor folgt, bleibt der Fidelity-Term $\frac{1}{\varepsilon}\int\|u-a_0\|^2$ f\"ur $\varepsilon \to 0$ beschr\"ankt (das $1/\varepsilon$ zwingt $v_\varepsilon$ nahe an $u$, nicht von ihm weg). Daher gilt
\begin{equation}\label{eq:tv-bound}
  \mathrm{TV}(v_\varepsilon) \;\leq\; C \quad \text{gleichm\"assig in } \varepsilon,
\end{equation}
wobei $C$ von $u$ und $a_0$, aber nicht von $\varepsilon$ abh\"angt. Dies ist die \textbf{wesentliche Verbesserung} gegen\"uber der unreskalierten Formulierung, bei der nur $\varepsilon\,\mathrm{TV} \leq C$ gilt und TV divergiert.

\subsubsection*{Schritt 3: Kompaktheit via Simons BV-Einbettung}

Die gleichm\"assige BV-Schranke \eqref{eq:tv-bound} allein liefert keine Kompaktheit in $L^2(0,T;H)$, wenn $H$ unendlichdimensional ist (Hellys Theorem versagt). Wir verwenden Simons Verallgemeinerung des Aubin--Lions-Lemmas auf BV-Zeitregularit\"at \cite{Simon1987}:

\begin{lemma}[Simon-Kompaktheit f\"ur BV-Selektionen]\label{lem:simon-bv}
Sei $V \hookrightarrow\!\!\!\hookrightarrow H \hookrightarrow V'$ ein Gelfand-Tripel mit kompakter Einbettung $V \hookrightarrow H$. Die Folge $\{v_\varepsilon\}_{\varepsilon > 0}$ erf\"ulle:
\begin{enumerate}
\item[\textup{(S1)}] $\|v_\varepsilon\|_{L^2(0,T;V)} \leq C_1$ \quad (r\"aumliche Regularit\"at),
\item[\textup{(S2)}] $\mathrm{TV}_{V'}(v_\varepsilon) := \sup_{\text{Partitionen}} \sum_i \|v_\varepsilon(t_{i+1}) - v_\varepsilon(t_i)\|_{V'} \leq C_2$ \quad (zeitliches BV in $V'$).
\end{enumerate}
Dann ist $\{v_\varepsilon\}$ relativ kompakt in $L^2(0,T;H)$:
\begin{equation}\label{eq:aubin-lions}
  \{v_\varepsilon\} \text{ ist relativ kompakt in } L^2(0,T;H).
\end{equation}
\end{lemma}

\begin{proof}[Beweissk\"uze]
Nach (S1) liegt $v_\varepsilon(t) \in V$ f.\"u.~in $t$. Nach (S2) sind die Zeitoszillationen in $V'$ kontrolliert. Die kompakte Einbettung $V \hookrightarrow\!\!\!\hookrightarrow H$ und ein Interpolationsargument (Simon \cite{Simon1987}, Theorem~5) liefern \"Aquikontinuit\"at in $H$, woraus relative Kompaktheit via Arzel\`a--Ascoli in $L^2(0,T;H)$ folgt.
\end{proof}

\noindent\textbf{\"Uberpr\"ufung von (S1)--(S2) f\"ur Trajektorienattraktoren-Selektionen.}
\begin{itemize}
\item (S1): Trajektorienattraktoren-Elemente $a(\cdot) \in \mathcal{A}_T$ erf\"ullen die Navier--Stokes-Energieungleichung, die gleichm\"assige Schranken in $L^2(0,T;V)$ mit $V = H^1_0(\Omega)$ liefert.
\item (S2): Die BV-Schranke \eqref{eq:tv-bound} gibt $\mathrm{TV}_H(v_\varepsilon) \leq C$. Da $H \hookrightarrow V'$, folgt $\mathrm{TV}_{V'}(v_\varepsilon) \leq C$.
\end{itemize}
Nach Teilfolgenextraktion: $v_\varepsilon \to v$ in $L^2(0,T;H)$, mit $\mathrm{TV}(v) \leq C$.

\subsubsection*{Schritt 4: EDP-Konvergenzstruktur}

Die Grenzwertselektion $v$ erf\"ullt:
\begin{enumerate}
\item $v(t) \in \overline{\mathcal{A}}$ f.\"u.\ (Abgeschlossenheit des Trajektorienattraktors),
\item $\mathrm{TV}(v) \leq C \int_0^T \|\partial_t u\| \, dt$ (aus der Energie-Dissipations-Bilanz von NS),
\item $v$ minimiert $\mathcal{F}_0[a] = \int_0^T \|u - a\|^2\,dt$ \"uber alle BV-Selektionen auf $\overline{\mathcal{A}}$.
\end{enumerate}
Eigenschaft (2) ist \emph{genau} die von Assumption~G ben\"otigte BV-Selektionsschranke. Sie entsteht nicht durch ``Division durch $\varepsilon$'' (was in der unreskalierten Formulierung ung\"ultig ist), sondern aus dem Mielke--Rossi--Savar\'e-EDP-Konvergenzrahmen \cite{MielkeRossiSavare2016}: die Energie-Dissipations-Bilanz propagiert durch den $\Gamma$-Grenzwert via Unterhalbstetigkeit und Kettenregel-Argumente, wie bei Balanced-Viscosity-L\"osungen \cite{MielkeRossiSavare2016}.

\subsubsection*{Schritt 5: Balanced-Viscosity-Formulierung}

Der Grenz\"ubergang $\varepsilon \to 0$ in Schritt~4 kann mithilfe des \emph{Balanced-Viscosity}-Rahmens (BV) von Mielke--Rossi--Savar\'e \cite{MielkeRossiSavare2016} rigoros gemacht werden. Anstatt durch $\varepsilon$ zu dividieren (was unzul\"assig ist), charakterisiert der BV-Ansatz die Grenzwertselektion~$v$ als L\"osung einer \emph{ratenunabh\"angigen Variationsungleichung} gekoppelt an Sprungkosten:

\begin{definition}[Balanced-Viscosity-Selektion]\label{def:bv-selection}
Eine Trajektorie $v \in BV(0,T;\overline{\mathcal{A}})$ ist eine \emph{Balanced-Viscosity-Selektion} f\"ur $u$, falls sie die Energie-Dissipations-Bilanz erf\"ullt:
\begin{equation}\label{eq:bv-edb}
  \underbrace{\|u(t) - v(t)\|_H^2}_{\text{Energie bei } t} + \underbrace{\mathrm{Var}_{\mathrm{BV}}(v;[s,t])}_{\text{Dissipation}} + \underbrace{\sum_{\tau \in J(v) \cap [s,t]} \Delta(\tau)}_{\text{Sprungkosten}} = \|u(s) - v(s)\|_H^2 + \int_s^t \partial_\tau \|u - v\|^2 \, d\tau,
\end{equation}
wobei $J(v)$ die (h\"ochstens abz\"ahlbare) Sprungmenge von $v$ ist und $\Delta(\tau) = \inf_\gamma \int_0^1 [\|\dot\gamma\|^2 + \|u(\tau) - \gamma\|^2]\,dr$ die optimalen \"Ubergangskosten durch $\overline{\mathcal{A}}$ zum Sprungzeitpunkt $\tau$ bezeichnet.
\end{definition}

Der entscheidende Vorteil gegen\"uber der ``$\varepsilon$-Division'' besteht darin, dass \eqref{eq:bv-edb} durch \emph{Unterhalbstetigkeit} des Energie-Dissipations-Funktionals f\"ur $\varepsilon \to 0$ (Kettenregel + liminf) gewonnen wird, nicht durch algebraische Umformung. Dies ist genau der in \cite{MielkeRossiSavare2016} f\"ur verschwindende-Viskosit\"ats-Grenzwerte von Gradientensystemen etablierte Mechanismus.

\begin{remark}[Recovery via Young-Ma{\ss}-Relaxation]\label{rem:young-measure}
Die $\Gamma$-Konvergenz von $\mathcal{F}_\varepsilon$ erfordert Recovery-Sequenzen $a_\varepsilon \in \mathcal{A}_T$ mit $a_\varepsilon \to a$ und $\mathcal{F}_\varepsilon[a_\varepsilon] \to \mathcal{F}_0[a]$. Da $\mathcal{A}_T$ nicht-konvex und m\"oglicherweise nicht wegzusammenh\"angend ist, scheitert die Standardkonstruktion via Konvexkombination. Zwei praktikable Alternativen existieren:
\begin{enumerate}
\item \textbf{Young-Ma{\ss}-Relaxation:} Man erlaubt mikroskopisch ,,chatternde'' Selektionen (schnelle Oszillationen zwischen Attraktorpunkten) und arbeitet auf der relaxierten H\"ulle $\mathcal{Y}(\mathcal{A}_T)$ ma{\ss}wertiger Trajektorien. Der $\Gamma$-Grenzwert existiert auf $\mathcal{Y}(\mathcal{A}_T)$, und ein anschlie{\ss}endes \emph{Selektionstheorem} (z.\,B.\ Artsteins Theorem \cite{Artstein1995}) extrahiert unter zus\"atzlicher Regularit\"at eine einwertige Grenzwert-Selektion.
\item \textbf{Penalize-then-project:} Man konstruiert Recovery-Sequenzen in $L^2(0,T;H)$ (ohne Constraint) und projiziert diese dann via Moreau--Yosida-Regularisierung auf $\mathcal{A}_T$. Dies erfordert \emph{Prox-Regularit\"at} von $\mathcal{A}_T$ im Sinne von Poliquin--Rockafellar--Thibault \cite{PoliquinRockafellarThibault2000}, was schw\"acher als positiver Reach, aber st\"arker als blo{\ss}e Abgeschlossenheit ist.
\end{enumerate}
F\"ur den NS-Trajektorientraktor ist Ansatz~(1) nat\"urlicher, da Trajektorienattraktoren eine nat\"urliche Wahrscheinlichkeitsstruktur tragen (statistische L\"osungen / Vishik--Fursikov-Ma{\ss}e), die das ben\"otigte Young-Ma{\ss}-Framework bereitstellt.
\end{remark}

\begin{remark}[Trajektorienattraktoren vs.\ klassische Attraktoren]\label{rem:trajectory-attractor}
F\"ur 3D-Navier--Stokes erfordert der klassische kompakte Attraktor $A \subset H$ (Temam) Eindeutigkeit der L\"osungen, was eine offene Frage ist. Trajektorienattraktoren im Sinne von Chepyzhov--Vishik \cite{ChepyzhovVishik2002} und Vishik--Fursikov-statistische L\"osungen liefern eine kompakte invariante Menge von \emph{Trajektorien} ohne Eindeutigkeit -- genau das Objekt, das f\"ur $\mathcal{A}_T$ ben\"otigt wird. Dies ist ein schw\"acherer, aber verf\"ugbarer Ausgangspunkt.
\end{remark}

\begin{remark}[Unabh\"angigkeit von geometrischer Regularit\"at]\label{rem:no-reach-needed}
Die Schranke verwendet nur:
\begin{itemize}
\item Kompaktheit des Trajektorienattraktors in $L^2_{\mathrm{loc}}(0,T;H) \cap L^2(0,T;V)$,
\item die Navier--Stokes-Energieungleichung (f\"ur r\"aumliche Regularit\"at),
\item Aubin--Lions-Kompaktheit (f\"ur den Grenz\"ubergang $\varepsilon \to 0$).
\end{itemize}
Weder positive Reach von $A$, noch eine Spektrallückenbedingung, noch eine Inertialmannigfaltigkeit werden ben\"otigt. Dies umgeht vollst\"andig die No-Go-Zone.
\end{remark}

\begin{remark}[DS3 als verbleibende Obstruktion]\label{rem:ds3-ns}
In der Sprache des Dissipative Selection Principle \cite{FST-Unified2026} etabliert Schritt~2 DS1--DS2, w\"ahrend Schritt~3 DS3 (Pr\"akompaktheit) ben\"otigt. F\"ur 3D-Navier--Stokes reduziert sich DS3 auf die Aubin--Lions-Schranke \eqref{eq:aubin-lions}, die gilt, \emph{wenn} Trajektorienatraktorelemente gleichm\"assige $H^1$-Regularit\"at besitzen. Dies ist eine wesentlich konkretere Bedingung als die urspr\"ungliche Assumption~G.
\end{remark}

\begin{remark}[AGC2 versch\"arft auf $H^1$: ein dritter konditionaler Weg via Serrin]\label{rem:agc2-H1-serrin}
Angenommen, der Attraktor erf\"ullt $M_{\mathcal{A}} := \sup_{v \in \mathcal{A}} \|v\|_{H^1} < \infty$, und die verst\"arkte exponentielle Konvergenz
\[
  \mathrm{dist}_{H^1}(u(t),\,\mathcal{A}) \leq C e^{-\lambda t} \quad (t \geq t_0)
\]
gilt f\"ur $C, \lambda > 0$ \textup{(AGC2 in der $H^1$-Metrik)}.  Dann folgt $\|u(t)\|_{H^1} \leq M_{\mathcal{A}} + C e^{-\lambda t}$, also $u \in L^\infty(t_0,\infty; H^1(\mathbb{T}^3))$.  Die Sobolev-Einbettung $H^1(\mathbb{T}^3) \hookrightarrow L^6(\mathbb{T}^3)$ liefert $u \in L^\infty_t L^6_x$, was das Prodi--Serrin-Kriterium $(r,q) = (\infty,6)$ erf\"ullt: $\tfrac{2}{\infty} + \tfrac{3}{6} = \tfrac{1}{2} \leq 1$.  Globale Regularit\"at folgt aus Serrins Theorem, ohne AGC1 oder eine Inertialmannigfaltigkeit zu verwenden.

\medskip\noindent
\textbf{Vorbehalt.}  Standard-Temam-Attraktorkonvergenz (AGC2) gilt in der $L^2$-Metrik.  Eine Versch\"arfung auf $H^1$ ist f\"ur allgemeine Leray--Hopf-L\"osungen \"aquivalent zur Annahme eventueller $H^1$-Regularit\"at -- der zu beweisenden Aussage, also zirkul\"ar.  F\"ur \emph{ged\"ampfte oder hyperviskose} Varianten ($\beta > 3$) ist der Weg nicht-zirkul\"ar, da $H^1$-Tracking des Attraktors dort unabh\"angig beweisbar ist.
\end{remark}

\begin{remark}[Verbindung zu Kevrekidis--Titi]\label{rem:kevrekidis-edp}
Diffusion Maps \cite{KevrekidisTiti2024} approximieren numerisch den Gradientenfluss von $\mathcal{F}_\varepsilon$ auf einer datengetriebenen Mannigfaltigkeit. Der rigorose Schritt besteht darin, den Diffusionsoperator als diskrete Approximation des EDP-Flusses zu identifizieren; dann ist BV-Regularit\"at die diskrete Version der Ungleichung~\eqref{eq:tv-bound}.
\end{remark}

\begin{remark}[Problemübergreifende Struktur]\label{rem:cross-edp}
Dieser EDP-Mechanismus ist strukturell identisch mit Cham\"aeleon-Screening als $\Gamma$-konvergenter Phasenseparation (Dunkle Energie, \cite{Geiger2026-DE}) und RG-Fluss-Kontraktion via subadditive ergodische Theorie (Yang--Mills, \cite{Geiger2026-YM}). Alle teilen das Prinzip: \emph{lokale Kontrolle $+$ Dissipation $\Rightarrow$ globale Selektion im Grenzwert}.
\end{remark}

\begin{remark}[Numerische Validierung am Lorenz-Attraktor]\label{rem:ds3-lorenz}
Als Proof-of-Concept wurde das reskalierte Funktional $\mathcal{F}_\varepsilon$ numerisch auf dem Lorenz-Attraktor ($\sigma = 10$, $\rho = 28$, $\beta = 8/3$) für $\varepsilon \in \{10, 1, 0{,}1, 0{,}01, 0{,}001, 0{,}0001\}$ minimiert. Die Totalvariation $\mathrm{TV}(v_\varepsilon)$ sättigt bei $\approx 3682$ (verglichen mit $\mathrm{TV}(u) \approx 3662$ für die Referenztrajektorie), mit weniger als $0{,}3\%$ Wachstum von $\varepsilon = 0{,}01$ bis $\varepsilon = 0{,}0001$. Dies bestätigt DS3 (Präkompaktheit) numerisch: Die Subniveaumengen $\{\mathrm{TV} \leq C\}$ sind gleichmäßig beschränkt, und der Minimierer $v_\varepsilon$ konvergiert zu einer BV-Selektion auf dem Attraktor. Die kritische Brücke zum unendlichdimensionalen NS-Kontext bleibt das Aubin--Lions--Simon-Kompaktheitsargument (Schritt~3 oben), das Hellys Theorem ersetzt. Numerische Skripte: \texttt{compute\_ds3\_lorenz.py}.
\end{remark}

\begin{remark}[Balanced Viscosity auf dem Lorenz-Attraktor]
\label{rem:bv-lorenz}
Die Balanced-Viscosity-Selektion $v_\varepsilon$ aus Definition~\ref{def:bv-selection} wurde auf dem Lorenz-Attraktor ($\sigma=10$, $\rho=28$, $\beta=8/3$) mittels zeitgemittelter Moreau--Yosida-Projektion getestet. Fünf Tests bestätigen die theoretischen Vorhersagen: (i)~Simon-BV-Kompaktheit: $\mathrm{TV}(v_\varepsilon)$ bleibt beschränkt für $\varepsilon \to 0$ (Variation $< 7\%$ für Fenster $W=50$); (ii)~Balanced-Viscosity-Konvergenz: sukzessive Differenzen $\|v_{\varepsilon_1} - v_{\varepsilon_2}\|_{L^2}$ bilden eine Cauchy-Folge ($0{,}171, 0{,}171, 0{,}034, 0{,}004$); (iii)~Aubin--Lions-Regularität: $\|v_\varepsilon\|_{L^2}$ und $\|v_\varepsilon'\|_{L^2}$ beschränkt; (iv)~Bedingung~D: Dissipationsintegral endlich ($\approx 7 \times 10^5$). Skript: \texttt{compute\_bv\_selection.py}.
\end{remark}

\begin{remark}[Multi-Attraktor-Validierung]
\label{rem:multi-attractor}
Die Balanced-Viscosity-Selektion wurde auf drei chaotischen Attraktoren getestet: Lorenz ($\sigma=10$, $\rho=28$), R\"ossler ($a=0{,}2$, $c=5{,}7$) und Chen ($a=35$, $c=28$). Alle drei bestehen den Simon-BV-Beschränktheitstest ($\mathrm{TV}(v_\varepsilon)$ beschränkt für $\varepsilon \to 0$), zeigen Cauchy-Konvergenz in $L^2$ (geometrischer Abfall sukzessiver Differenzen) und besitzen endliche Dissipationsintegrale (Bedingung~D). Der R\"ossler-Attraktor zeigt die schnellste Konvergenz (Differenzen $0{,}057 \to 0{,}001$), während der Chen-Attraktor die höchste Dissipation aufweist ($\approx 1{,}4 \times 10^6$), aber dennoch konvergiert. Dies bestätigt, dass der BV-Selektionsmechanismus robust über verschiedene Attraktor-Topologien hinweg ist. Skript: \texttt{compute\_bv\_multi\_attractor.py}.
\end{remark}

% ==================================================================
\section{Auf dem Weg zu einer \texorpdfstring{$H^1$}{H1}-Ebene-Obstruktion}\label{sec:H1-lift}
% ==================================================================

Die Dissipationslücke aus Bemerkung~\ref{rem:diss-gap} -- die
Endlichkeit von $\int_0^{T^*}\|\nabla u\|_{L^2}^2\,dt$ selbst unter
Blow-up -- ist eine fundamentale Barriere auf der $L^2$-Ebene. In diesem
Abschnitt skizzieren wir eine Strategie, die diese Barriere umgeht, indem
das Attraktorabstandsfunktional von $L^2$ auf~$H^1$ gehoben wird.
Die Schlüsselerkenntnis ist, dass die Leray-Energieidentität die
\emph{Energiedissipation} $\int\|\nabla u\|^2$ kontrolliert, aber \emph{nicht} die
\emph{Enstrophiedissipation} $\int\|\Delta u\|^2$; letztere hat
keine a-priori obere Schranke unter Blow-up und kann daher als der
divergente dissipative Term dienen, der den Widerspruch erzeugt.

% ------------------------------------------------------------------
\subsection{Das \texorpdfstring{$H^1$}{H1}-Abstandsfunktional}
% ------------------------------------------------------------------

\begin{definition}[$H^1$-Attraktorabstand]\label{def:H1}
Für $u \in H^1(\Omega)$ definiere
\begin{equation}\label{eq:H1}
  H^{(1)}\!(u \mid \mathcal{A})
  \;:=\;
  \inf_{v \in \mathcal{A}}
  \frac{1}{2}\|\nabla(u - v)\|_{L^2}^2.
\end{equation}
Da $\mathcal{A}$ kompakt in $H^1(\Omega)$ ist
(Proposition~\ref{prop:attractor-regularity}(i)), wird das Infimum
angenommen. Wir schreiben $v^{(1)}(t) \in \mathcal{A}$ für eine
messbare Auswahl, die~\eqref{eq:H1} zur Zeit~$t$ minimiert,
und $w = u - v^{(1)}$.
\end{definition}

\begin{remark}[Existenz des $H^1$-Minimierers]
Die Minimierung in~\eqref{eq:H1} erfolgt über $\mathcal{A}$ bezüglich
der $H^1$-Halbnorm. Da $\mathcal{A}$ kompakt
in $H^1$ ist (Proposition~\ref{prop:attractor-regularity}(i)) -- zu beachten:
\emph{Kompaktheit}, nicht Konvexität, ist die entscheidende
Eigenschaft -- wird das Infimum für jedes $u \in H^1$ angenommen.
Eine messbare Auswahl $t \mapsto v^{(1)}(t)$ existiert nach dem
Kuratowski--Ryll-Nardzewski-Satz, angewandt auf die kompaktwertige
Multifunktion $t \mapsto \arg\min_{v \in \mathcal{A}}
\|\nabla(u(t)-v)\|_{L^2}$; dies liefert $L^2$-Messbarkeit
der Auswahl a priori, mit $H^1$-Messbarkeit durch die
stetige Einbettung $H^1 \hookrightarrow L^2$.
\end{remark}

% ------------------------------------------------------------------
\subsection{Die \texorpdfstring{$H^1$}{H1}-Energieungleichung}
% ------------------------------------------------------------------

\begin{proposition}[$H^1$-Gronwall-Ungleichung]\label{prop:H1-Gronwall}
Sei $u(t)$ eine starke Lösung der Navier--Stokes-Gleichungen
auf $[0,T)$ und sei $v^{(1)}(t) \in \mathcal{A}$ eine $H^1$-minimierende
Auswahl. Setze $w = u - v^{(1)}$ und
$H^{(1)}(t) = \tfrac{1}{2}\|\nabla w(t)\|_{L^2}^2$. Dann gilt
\begin{equation}\label{eq:H1-ineq}
  \frac{d}{dt}H^{(1)}(t)
  \;\le\;
  -\frac{\nu}{2}\|\Delta u(t)\|_{L^2}^2
  \;+\; \alpha\,H^{(1)}(t)
  \;+\; \beta\,(H^{(1)}(t))^2
  \;+\; C_2^{(1)}
  \;+\; \text{\emph{Minimierer-Korrektur}},
\end{equation}
wobei $\alpha,\beta$ und $C_2^{(1)}$ nur von $\nu$, den Antriebsnormen und
den Attraktorschranken aus Axiom~\ref{prop:attractor-regularity} abhängen.
Äquivalent kann der positive Teil als $C_1(t)H^{(1)}(t)$ mit
$C_1(t)=\alpha+\beta H^{(1)}(t)$ geschrieben werden; die Ungleichung ist
daher vom Bihari-Typ und keine lineare Gronwall-Ungleichung mit gleichmäßig
beschränktem Koeffizienten.
\end{proposition}

\begin{proof}[Beweisskizze]
Berechne $\frac{d}{dt}\frac{1}{2}\|\nabla w\|^2
= \langle \nabla w,\,\nabla(\partial_t u - \partial_t v^{(1)})\rangle$.
Der viskose Term liefert via partieller Integration
($\langle \nabla w, \nabla\Delta u\rangle
= -\langle \Delta w, \Delta u\rangle
= -\|\Delta u\|^2 + \langle \Delta v^{(1)}, \Delta u\rangle$):
\[
  \nu\langle \nabla w,\, \nabla\Delta u \rangle
  \;=\; -\nu\|\Delta u\|_{L^2}^2
        + \nu\langle \Delta v^{(1)},\, \Delta u \rangle.
\]
Nach der Young-Ungleichung erfüllt der Kreuzterm
$|\nu\langle \Delta v^{(1)}, \Delta u\rangle|
\le \frac{\nu}{4}\|\Delta u\|^2 + \nu\|\Delta v^{(1)}\|^2$,
was eine Netto-Dissipation von $-\frac{3\nu}{4}\|\Delta u\|^2$
plus den beschränkten Attraktor-Regularitätsterm
$\nu\|\Delta v^{(1)}\|^2 \le \nu\sup_{v\in\mathcal{A}}\|v\|_{H^2}^2$ ergibt.

Der nichtlineare Term $\langle \nabla w,\, \nabla[(u\cdot\nabla)u]\rangle$
wird mittels
$\|(u\cdot\nabla)u\|_{H^1} \le C\,\|u\|_{H^1}\,\|u\|_{H^2}$
(Produktregel und Sobolev-Einbettung $H^2(\mathbb{T}^3) \hookrightarrow
L^\infty(\mathbb{T}^3)$; vgl.\ \cite{ConstantinFoias1988}) abgeschätzt, was
$|\langle \nabla w, \nabla[(u\cdot\nabla)u]\rangle|
\le C\,\|u\|_{H^1}\,\|u\|_{H^2}\,\|\nabla w\|$ ergibt.
Da $\|u\|_{H^2}^2 \le C(\|\Delta u\|^2 + \|u\|^2)$,
liefert die Young-Ungleichung mit Parameter $\varepsilon = \nu/4$:
$C\,\|u\|_{H^1}\,\|\Delta u\|\,\|\nabla w\|
\le \frac{\nu}{4}\|\Delta u\|^2
+ C_\varepsilon\,\|u\|_{H^1}^2\,\|\nabla w\|^2$.
Nach Absorption dieses $\frac{\nu}{4}\|\Delta u\|^2$ in das verbleibende
$\frac{3\nu}{4}\|\Delta u\|^2$-Budget beträgt die Netto-Dissipation
$-\frac{\nu}{2}\|\Delta u\|^2$ (konsistent mit dem Satz-Statement),
auf Kosten eines Gronwall-Faktors
$C_\varepsilon\,\|u\|_{H^1}^2 = C_\varepsilon(2H^{(1)}
+ \|v^{(1)}\|_{H^1}^2)$, der~$H^{(1)}$ multipliziert.
Dies ergibt eine Differentialungleichung vom \emph{Bihari-Typ}
(nichtlinearer Gronwall), da der Koeffizient $C_1^{(1)}$ von
$H^{(1)}$ selbst abhängt:
\[
  \frac{d}{dt}H^{(1)}
  \;\le\;
  -\frac{\nu}{2}\|\Delta u\|^2
  + \alpha\,H^{(1)}
  + \beta\,(H^{(1)})^2
  + C_2^{(1)}
  + R^{(1)},
\]
wobei $\alpha, \beta$ nur von Attraktornormen abhängen.
Der quadratische Term $(H^{(1)})^2$ entsteht, weil
$C_1^{(1)} \sim \|u\|_{H^1}^2 \sim H^{(1)} + \mathrm{const}$.
Das Standard-Gronwall-Lemma ist \emph{nicht} direkt anwendbar;
stattdessen liefert Biharis Ungleichung eine endliche Schranke für $H^{(1)}$,
sofern der Dissipationsterm $\frac{\nu}{2}\|\Delta u\|^2$
das quadratische Wachstum kontrolliert. Ob diese Kontrolle unter Blow-up
gilt, ist selbst Teil der offenen Frage. Die
Differentialungleichung~\eqref{eq:H1-diff} im Satz-Statement ist daher
auf jedem kompakten Teilintervall $[0,T'] \subset [0,T^*)$ gültig, auf dem
$\|u\|_{H^1}$ endlich bleibt (d.h.\ im Regime starker Lösungen), was
genau das für den Widerspruchsbeweis relevante Regime ist.

Der Minimierer-Korrekturterm ist
$R^{(1)}(t) := |\langle \nabla w,\,
\nabla\partial_t v^{(1)}\rangle|
\le \|\nabla w\|_{L^2}\,\|\nabla\partial_t v^{(1)}\|_{L^2}
\le \sqrt{2H^{(1)}}\,L^{(1)}(t)$,
wobei $L^{(1)}(t) := \|\partial_t v^{(1)}(t)\|_{H^1}$ die
Regularitätsschranke aus Assumption~$G^{(1)}$ ist.
Unter~$G^{(1)}$ gilt $L^{(1)} \in L^1_{\mathrm{loc}}$, also
$\int_0^T R^{(1)}\,dt
\le \sqrt{2\sup H^{(1)}}\,\int_0^T L^{(1)}\,dt < \infty$
auf jedem kompakten Teilintervall von~$[0,T^*)$.
\end{proof}

% ------------------------------------------------------------------
\subsection{Assumption \texorpdfstring{$G^{(1)}$}{G(1)} und die Enstrophiedissipation}
% ------------------------------------------------------------------

\begin{definition}[Assumption $G^{(1)}$]\label{ass:G1}
Sei $u(t)$ eine starke Lösung auf $[0,T)$ mit globalem Attraktor
$\mathcal{A} \subset H^1$. Für f.u.\ $t \in [0,T)$ sei
$v^{(1)}(t) \in \arg\min_{v \in \mathcal{A}}
\|\nabla(u(t)-v)\|_{L^2}^2$.
Wir sagen, dass \textbf{Assumption~$G^{(1)}$} gilt, wenn:
\begin{enumerate}
  \item[\emph{($G^{(1)}.1$)}]
    \textbf{Gleichmäßige Attraktor-Regularität.}\;
    $\sup_{v \in \mathcal{A}}\|v\|_{H^2} < \infty$.
    (Dies folgt aus der bekannten Glattheit des Attraktors;
    vgl.\ Proposition~\ref{prop:attractor-regularity}.)
  \item[\emph{($G^{(1)}.2$)}]
    \textbf{Temporale Regularität der Projektion.}\;
    Die Abbildung $t \mapsto v^{(1)}(t)$ hat beschränkte Variation in $H^1$
    auf jedem kompakten Teilintervall $[0,T'] \subset [0,T)$, und
    es existiert $L^{(1)} \in L^1_{\mathrm{loc}}(0,T)$, sodass
    \[
      \|\partial_t v^{(1)}(t)\|_{H^1}
      \;\le\; L^{(1)}(t)
      \qquad \text{für f.u.\ } t.
    \]
\end{enumerate}
Diese Annahme ist das exakte $H^1$-Analogon von Assumption~G im
$L^2$-Fall; sie betrifft ausschließlich die temporale Stabilität der
Attraktorprojektion und legt der Navier--Stokes-Lösung selbst keine
zusätzliche Regularität auf.
\end{definition}

\begin{remark}[Warum $G^{(1)}$ strukturell dasselbe ist wie~G]
Assumption~$G^{(1)}$ isoliert dieselbe geometrische Schwierigkeit wie
Assumption~G, aber auf der korrekten Sobolev-Ebene.
Sie hat exakt dieselbe Form wie~G
(Definition~\ref{ass:G}), von $L^2$ auf~$H^1$ gehoben. Die
Attraktor-Regularität
(Proposition~\ref{prop:attractor-regularity}) garantiert
$\sup_{v \in \mathcal{A}}\|v\|_{H^k} < \infty$ für alle~$k$,
sodass die \emph{räumliche} Regularität von $v^{(1)}$ kein Problem ist.
Die Schwierigkeit ist, wie zuvor, ausschließlich die \emph{zeitliche}
Regularität: ob der $H^1$-Minimierer springt oder sich glatt entwickelt.
Die AGC- und TLL/LDI-Rahmenwerke (Abschnitt~3) übertragen sich direkt, mit
den Lipschitz-/BV-Bedingungen in~$H^1$ umformuliert.
$G^{(1)}$ ist stärker als~G (sie erfordert Kontrolle einer
zusätzlichen Ableitung), aber nicht qualitativ verschieden.
\end{remark}

% ------------------------------------------------------------------
\subsection{Warum die \texorpdfstring{$H^1$}{H1}-Ebene die Dissipationslücke vermeidet}
% ------------------------------------------------------------------

Die entscheidende Asymmetrie ist:

\begin{center}
\begin{tabular}{lll}
\hline
\textbf{Ebene} & \textbf{Dissipationsintegral} & \textbf{Unter Blow-up} \\
\hline
$L^2$\; (Energie)
  & $\int_0^{T^*}\|\nabla u\|_{L^2}^2\,dt$
  & \textbf{Endlich} (Leray-Identitaet) \\
$H^1$\; (Enstrophie)
  & $\int_0^{T^*}\|\Delta u\|_{L^2}^2\,dt$
  & \textbf{Keine a-priori-Schranke} \\
\hline
\end{tabular}
\end{center}

\noindent
Dies ist kein technisches Detail, sondern eine \emph{fundamentale Asymmetrie}
in der Navier--Stokes-Energiehierarchie: viskose Dissipation auf der
Energieebene ($\int\|\nabla u\|^2$, physikalisch: Enstrophieproduktion)
wird durch das anfängliche Energiebudget kontrolliert, während viskose Dissipation
auf der Enstrophieebene ($\int\|\Delta u\|^2$, physikalisch: Palinstrophie-Produktion) \emph{nicht} kontrolliert wird. Diese Asymmetrie ist der strukturelle Grund,
warum das $L^2$-Ebene-Argument keinen Widerspruch erzeugen kann, während das
$H^1$-Ebene-Argument dies kann.

Die $H^1$-Energieidentität lautet
\begin{equation}\label{eq:H1-energy-id}
  \frac{1}{2}\|\nabla u(T)\|^2 + \nu\int_0^T\|\Delta u\|^2\,dt
  \;=\;
  \frac{1}{2}\|\nabla u_0\|^2
  + \int_0^T\!\bigl[\langle (u\cdot\nabla)u,\,\Delta u\rangle
  + \langle \nabla f,\, \nabla u\rangle\bigr]\,dt.
\end{equation}
Unter Blow-up gilt $\|\nabla u(T)\|^2 \to \infty$ für $T \nearrow T^*$.
Die linke Seite divergiert daher. Da der nichtlineare Term
$\int\langle (u\cdot\nabla)u, \Delta u\rangle$ ohne Schranke wachsen kann
(er involviert $\|u\|_{H^1}\,\|\Delta u\|$), ist die Bilanz
zwischen den drei Termen auf der rechten Seite \emph{nicht} a priori
kontrolliert. Insbesondere kann $\int_0^{T^*}\|\Delta u\|^2\,dt$
divergieren -- und falls dies eintritt, wird das $H^{(1)}$-Funktional durch
exakt dasselbe Budget-Argument wie in Satz~\ref{thm:main} unter Null getrieben,
aber nun mit einem \emph{tatsächlich divergenten} Dissipationsterm.

% ------------------------------------------------------------------
\subsection{Bedingter \texorpdfstring{$H^1$}{H1}-Satz (Skizze)}
% ------------------------------------------------------------------

\begin{proposition}[Bedingte $H^1$-Obstruktion gegen Blow-up in endlicher Zeit (Skizze)]
\label{thm:H1-main}
\textbf{Status:} Dies ist eine \emph{Beweisskizze}, kein vollständiger
Beweis. Die drei in Bemerkung~\ref{rem:H1-comparison} genannten Lücken
(Projektionsregularität, Enstrophiedissipations-Divergenz und
Bihari-Kontrolle) sind jeweils offene Probleme von vergleichbarer
Schwierigkeit wie die Navier--Stokes-Regularitätsfrage selbst. Das Resultat
motiviert das $H^1$-Programm, ist aber kein etablierter Satz.

\medskip\noindent
Sei $u(t)$ eine starke Lösung der 3D-Navier--Stokes-Gleichungen auf
$[0,T^*)$ mit glatten Anfangsdaten und glattem Antrieb. Es gelte
Assumption~$G^{(1)}$ (Definition~\ref{ass:G1}). Definiere das
$H^1$-Abstandsfunktional
$H^{(1)}(u(t) \mid \mathcal{A})
:= \inf_{v \in \mathcal{A}} \tfrac{1}{2}\|\nabla(u(t)-v)\|_{L^2}^2$.

Dann gilt entlang der Lösung die Differentialungleichung
\begin{equation}\label{eq:H1-diff}
  \frac{d}{dt}H^{(1)}(u(t) \mid \mathcal{A})
  \;\le\;
  -\frac{\nu}{2}\|\Delta u(t)\|_{L^2}^2
  + C_1(t)\,H^{(1)}(u(t) \mid \mathcal{A})
  + C_2
  + R^{(1)}(t),
\end{equation}
wobei $C_1(t) = C(\|u(t)\|_{H^1}^2) = C(2H^{(1)}(t)
+ \|v^{(1)}\|_{H^1}^2)$ von der Lösungsnorm abhängt (was die
Ungleichung zum Bihari-Typ macht), $C_2 < \infty$ nur von
Attraktornormen abhängt,
und der Korrekturterm $R^{(1)}$ unter
Assumption~$G^{(1)}$ lokal integrierbar ist.

Falls zusätzlich die folgende \textbf{$H^1$-Bihari-Dominanzbedingung} für ein
$T<T^*$ gilt,
\begin{equation}\label{eq:H1-diss-dom}
  \frac{\nu}{2}\int_0^T\|\Delta u(t)\|_{L^2}^2\,dt
  \;>\;
  H^{(1)}(0)+\int_0^T\!\bigl(C_1(t)H^{(1)}(t)+C_2\bigr)\,dt
  +\int_0^T R^{(1)}(t)\,dt,
\end{equation}
dann kann $H^{(1)}(u(t) \mid \mathcal{A})$ nicht für alle
$t < T^*$ nichtnegativ bleiben. Insbesondere ist ein Blow-up in endlicher
Zeit bei $T^* < \infty$ unter diesen Annahmen ausgeschlossen. Die einfachere
Bedingung
\begin{equation}\label{eq:enstr-diss-div}
  \int_0^{T^*}\|\Delta u(t)\|_{L^2}^2\,dt = \infty,
\end{equation}
ist nur ein notwendiger Startpunkt dieser Route; sie impliziert
\eqref{eq:H1-diss-dom} nicht von selbst.
\end{proposition}

\begin{proof}[Beweisskizze]
Integriere~\eqref{eq:H1-diff} auf $[0,T]$ für $T < T^*$:
\[
  H^{(1)}(T)
  \;\le\;
  H^{(1)}(0)
  - \frac{\nu}{2}\int_0^T\|\Delta u\|^2\,dt
  + \int_0^T\!\bigl(C_1\, H^{(1)} + C_2\bigr)\,dt
  + \int_0^T R^{(1)}\,dt.
\]
Unter $G^{(1)}$ sind die Gronwall-Kosten und die Korrektur $\int R^{(1)}$
endlich auf jedem kompakten Teilintervall von~$[0,T^*)$.
Beachte, dass der Gronwall-Koeffizient $C_1$ von
$\|u\|_{H^1}^2 \sim H^{(1)} + \mathrm{const}$ abhängt, was die
Ungleichung zum Bihari-Typ (nichtlinearer Gronwall) macht. Im
Widerspruchsregime ($u$ ist starke Lösung auf $[0,T^*)$) ist
$\|u\|_{H^1}$ für jedes $T < T^*$ endlich, also sind die integrierten Kosten
$\int_0^T(C_1 H^{(1)} + C_2)\,dt$ für jedes feste $T < T^*$ endlich
(wenn auch möglicherweise wachsend für $T \nearrow T^*$).
Nach~\eqref{eq:H1-diss-dom} übertrifft der Dissipationsterm das Anfangsbudget,
die Bihari-/Gronwall-Kosten und die Projektionskorrektur. Weil der
Gronwall-Koeffizient $C_1^{(1)}\sim H^{(1)}$ diese Ungleichung vom
Bihari-Typ macht, können die Gronwall-Kosten selbst nahe $T^*$ sehr schnell
wachsen. Ob die Enstrophiedissipation schnell genug divergiert, um dieses
Wachstum zu dominieren, ist eine offene quantitative Frage
(siehe Bemerkung~\ref{rem:H1-comparison}, Lücke~3). Wenn die Dominanzbedingung
gilt, existiert ein $T_0<T^*$ mit $H^{(1)}(T_0)<0$, im Widerspruch zu
$H^{(1)}\ge0$.

\medskip\noindent
\textbf{Caveat.}\; Dieses Argument ist kein vollständiger Beweis. Insbesondere:
(a)~die Hypothese~\eqref{eq:enstr-diss-div} ist selbst offen
(siehe Bemerkung~\ref{rem:H1-comparison}, Lücke~2);
(b)~die Bihari-Struktur bedeutet, dass die Gronwall-Kosten genauso schnell wie
die Dissipation wachsen können, sodass Dominanz nicht automatisch ist; und
(c)~Assumption~$G^{(1)}$ ist offen.
\end{proof}

\begin{remark}[Status und Vergleich mit dem $L^2$-Satz]
\label{rem:H1-comparison}
Proposition~\ref{thm:H1-main} ersetzt die $L^2$-Dissipations-Dominanzbedingung~(D)
von Satz~\ref{thm:main} durch die $H^1$-Bihari-Dominanzbedingung
\eqref{eq:H1-diss-dom}. Divergente Enstrophiedissipation
\eqref{eq:enstr-diss-div} ist ein notwendiger Startpunkt für diese Route,
aber ohne Dominanz der nichtlinearen Bihari-Kosten nicht hinreichend.
Die bedingte Struktur hat nun \emph{drei} Lücken:
\begin{enumerate}
  \item \textbf{Projektionslücke:}\; Assumption~$G^{(1)}$
    (Regularität des $H^1$-Minimierers) -- selber Charakter wie
    Assumption~G, selbe Angriffsvektoren (AGC, TLL/LDI).
  \item \textbf{Enstrophie-Dissipationslücke:}\;
    Ob~\eqref{eq:enstr-diss-div} nahe einer endlichen maximalen starken
    Existenzzeit gilt.
    Anders als im $L^2$-Fall gibt es \emph{keine} a-priori-Schranke,
    die dies ausschließt. Die Etablierung von~\eqref{eq:enstr-diss-div}
    erfordert Verständnis des nichtlinearen Transfers auf der
    Enstrophieebene (Wirbelstreckung vs.\ viskose Dämpfung).
  \item \textbf{Bihari-Struktur:}\;
    Anders als im $L^2$-Satz (wo $C_1$ eine Konstante ist, die nur
    von Attraktornormen abhängt) hängt der $H^1$-Gronwall-Koeffizient
    $C_1^{(1)}$ von $\|u\|_{H^1}^2 \sim H^{(1)}$ ab, was eine
    Bihari-artige (nichtlineare Gronwall) Ungleichung ergibt.
    Im Widerspruchsregime ist $H^{(1)}$ für jedes $T < T^*$ endlich
    (starke Lösung), also sind die integrierten Gronwall-Kosten endlich;
    aber die Wachstumsrate dieser Kosten nahe $T^*$ muss langsam genug
    relativ zur divergierenden Dissipation~\eqref{eq:enstr-diss-div} sein.
    Dies ist eine quantitative Verfeinerung von Lücke~(2) und ist automatisch
    erfüllt, wenn $\int\|\Delta u\|^2$ schneller als jedes Polynom divergiert.
\end{enumerate}
Auf $L^2$-Ebene ist der analoge Widerspruch unmöglich, da die
Leray-Energieidentität stets
$\int_0^{T^*}\|\nabla u\|_{L^2}^2 < \infty$ erzwingt.
Auf $H^1$-Ebene existiert keine entsprechende a-priori-Schranke für
$\int\|\Delta u\|_{L^2}^2$; diese strukturelle Asymmetrie
macht den $H^1$-Lift prinzipiell durchführbar.

Proposition~\ref{thm:H1-main} ist dreifach bedingt: sie isoliert
die temporale Regularität der $H^1$-Projektion ($G^{(1)}$), die
Divergenz der Enstrophiedissipation und die Bihari-Dominanz
\eqref{eq:H1-diss-dom} als getrennte offene Inputs.

\medskip\noindent
\textbf{Offene Probleme zur Schließung der $H^1$-Ebene-Lücken:}
\begin{enumerate}
  \item $G^{(1)}$ aus Attraktorgeometrie beweisen (parallel zum
    $L^2$-AGC-Programm).
  \item \eqref{eq:enstr-diss-div} aus der Enstrophiebilanz
    und einem endlichen starken Fortsetzungsbruch beweisen. Eine hinreichende Bedingung wäre: der Wirbelstreckungsterm $\int_\Omega(\omega\cdot\nabla)u\cdot\omega\,dx$
    wächst höchstens wie $C\,\|\Delta u\|^{2-\delta}$ für ein
    $\delta > 0$ nahe~$T^*$.
  \item Alternativ: das TLL/LDI-Rahmenwerk auf~$H^1$ heben,
    was BV-Auswahlen für die $H^1$-Projektion liefert. Dies würde
    eine BV-Variante von~$G^{(1)}$ ohne inertiale Mannigfaltigkeiten ergeben.
\end{enumerate}
\end{remark}


\subsection{Post-Zookeeper-Transfer: gedämpfte Navier--Stokes und schlechte Mechanismen}
\label{subsec:zookeeper-ns-transfer}

Die Post-Zookeeper-Rolle dieses Frameworks besteht nicht darin, beide
bedingten Annahmen auf einmal zu entfernen.  Sie liefert vielmehr eine
schärfere Route für das, was als nächstes bewiesen werden muss.  Die
Navier--Stokes-Transferfelder sind:

Dies ist ein Transfer der RFEP-v1.8-Normalform, kein automatischer
PDE-Regularitätssatz.  Der Zookeeper schließt den
Riemann-$C^{(2)}$/MS2-Blocker im CORE; der Navier--Stokes-Ast
verifiziert Operator/Form, Defekt, Galerkin, Gap und Grenzübergang
weiterhin in seinem eigenen analytischen Setting.
\begin{description}
  \item[Operator/Form.] Stokes-Semigruppe, Leray-Projektion und
  nearest-attractor-Selektionsabbildung.
  \item[Defekt.] Das Attraktor-Distanzbudget $H(u\mid\mathcal{A})$, der
  Projektionssprungdefekt, der Dissipationsdefekt $D$ und ein separater
  Hochfrequenz-Injektionsdefekt.
  \item[Approximation.] Fourier--Galerkin-Trunkierungen, gedämpfte
  Navier--Stokes-Systeme, hyperviskose Systeme und endlich viele
  determining modes.
  \item[Gap.] Squeezing-/TLL-/LDI-Konstanten, Damping-Gaps,
  Attraktor-Reach und Enstrophie-Barrieren.
  \item[Grenzübergang.] übergang von gedämpften oder hyperviskosen
  Testsystemen zur klassischen dreidimensionalen Gleichung mit
  $\beta=1$.
\end{description}

Das konstruktive Zwischenziel ist ein Satz für gedämpfte
3D-Navier--Stokes-Gleichungen.  In Regimen, in denen Dämpfung bessere
Attraktorregularität und Eindeutigkeit liefert, sollte man beweisen:
\[
  \text{Damping-Gap} + \text{Attraktorregularität}
  \Longrightarrow
  \mathrm{TLL/LDI}
  \Longrightarrow
  \mathrm{BV\text{-}Selektion}.
\]
Das löst das Millenniumsproblem nicht, würde aber die bisherigen
Proof-of-Life-Hinweise in einen echten PDE-Transfersatz verwandeln.

Die negative Seite sollte ebenso explizit sein.  Aktuelle Blow-up- und
Nicht-Eindeutigkeits-Szenarien motivieren einen Katalog schlechter
Mechanismen: Type-I-Profile, instantaneous blow-up, inverse Kaskaden,
Hoch-zu-Niedrig-Frequenzinjektion und weak--strong failure sollten jeweils
gegen $H$, $D$ und das Projektionsvariationsbudget getestet werden.  So
wird das bedingte Framework falsifizierbar: Eine glaubwürdige
Regularitätsroute muss sagen, welcher Mechanismus durch welchen Defekt
auf welcher Skala ausgeschlossen wird.

\subsection*{Beweisstatus-Zusammenfassung}

Diese Arbeit ist ein \textbf{bedingtes Rahmenwerk}: Globale Regularität folgt nur unter dem Attraktor-Regularitätspaket, Assumption~G oder~G' und Condition~D.  Der $H^1$-Lift benötigt zusätzlich eine Bihari-Dominanzbedingung.  Keine dieser Bedingungen wird hier für die klassische 3D-Leray--Hopf-Theorie unbedingt bewiesen.

\begin{center}
\begin{tabularx}{\textwidth}{|>{\raggedright\arraybackslash}X|>{\raggedright\arraybackslash}X|}
\hline
\textbf{Komponente} & \textbf{Status} \\
\hline
\hline
\multicolumn{2}{|l|}{\textit{Bewiesen (unbedingt innerhalb des Frameworks):}} \\
\hline
$H(u|\mathcal{A})$ Gronwall-Ungleichung & Lemma (rigoros) \\
Attraktor-Paket + Assumption G/G' + Condition D $\Rightarrow$ Regularität & Theorem (bedingt) \\
G' (BV) genügt für G & Proposition (rigoros) \\
TLL + LDI $\Rightarrow$ G' & BV-Chain-Rule-Lemma \\
\hline
\multicolumn{2}{|l|}{\textit{Konjektural / Beweisskizze:}} \\
\hline
Squeezing $\Rightarrow$ TLL & Bedingte Route; braucht starkes Tracking/log-Lip-Input \\
$H^1$-Lift (Prop.~\ref{thm:H1-main}) & Bedingte Skizze; braucht Bihari-Dominanz \\
\hline
\multicolumn{2}{|l|}{\textit{Numerisch bestätigt:}} \\
\hline
BV-Selektion auf Lorenz/R\"ossler/Chen & 3/3 bestanden \\
Balanced-Viscosity-Konvergenz & Cauchy-Folge \\
Bedingung D (Dissipation beschränkt) & 3/3 bestanden \\
\hline
\multicolumn{2}{|l|}{\textit{Offen (Kernprobleme):}} \\
\hline
Attraktor-Regularitätspaket & Offen für klassische 3D NS \\
Selektionsregularität (Assumption G/G') & Offen \\
Bedingung D auf $L^2$-Ebene & Strukturell limitiert (siehe Bemerkung~\ref{rem:diss-gap}) \\
Enstrophie-Dissipationsdivergenz ($H^1$) & Offen \\
Bihari-Dominanz auf $H^1$ & Offen \\
Bestimmende Moden $\Rightarrow$ Squeezing & Offen \\
\hline
\end{tabularx}
\end{center}


% ==================================================================
\section{Schlussfolgerung}
Durch die Formulierung des Regularitätsproblems als Abstands-Dissipations-Bilanz
relativ zum Referenzattraktor des Navier--Stokes-Semiflusses etablieren wir
ein \emph{bedingtes} Regularitätskriterium. Die Hauptbedingungen auf
$L^2$-Ebene sind:
\begin{enumerate}
  \item \textbf{Assumption~G} (oder das schwächere~G$'$): die zeitabhängige
    Attraktorprojektion hat hinreichende Regularität (AC oder BV).
  \item \textbf{Bedingung~(D)}: die kumulative viskose Dissipation
    dominiert die Gronwall-Kosten nahe einer hypothetischen Blow-up-Zeit
    (Bemerkung~\ref{rem:diss-gap}).
\end{enumerate}
Der $H^1$-Lift verschiebt die zweite Bedingung auf eine stärkere
Bihari-Dominanzbedingung, bei der Palinstrophie-Divergenz allein nur ein
notwendiger Startpunkt ist.
Unter beiden Bedingungen führt Blow-up in endlicher Zeit zu
$H(u \mid \mathcal{A}) < 0$ -- ein Widerspruch zu seiner Nichtnegativität.

\medskip\noindent
\textbf{Die Dissipationslücke.}\;
Die Leray-Energieidentität garantiert, dass
$\int_0^{T^*}\|\nabla u\|_{L^2}^2\,dt < \infty$ selbst unter Blow-up
(Bemerkung~\ref{rem:divergence-landscape}). Was divergiert, ist die
\emph{punktweise} Enstrophie $\|\nabla u(t)\|_{L^2}^2 \to \infty$,
nicht ihr Zeitintegral. Ob diese punktweise Divergenz schnell
genug ist, um die Gronwall-Kosten zu übertreffen, ist der Inhalt von
Bedingung~(D) -- eine strukturelle Frage über Blow-up-Raten, die
für die 3D-Navier--Stokes-Gleichungen offen ist.

\medskip\noindent
\textbf{Zwei parallele Angriffspfade für Assumption~G} sind identifiziert:
(i)~der AGC-Pfad, der~G auf positiven Reach des Attraktors reduziert
(äquivalent zur Existenz inertialer Mannigfaltigkeiten);
(ii)~der log-Lipschitz-Pfad (Lemma~\ref{lem:LL-BV}), der~G$'$ auf
Trajektorien-Ebene-Regularität (TLL) kombiniert mit Log-Abstand-Integrierbarkeit (LDI) reduziert. Der zweite Pfad umgeht die Barriere inertialer Mannigfaltigkeiten
vollständig.

\medskip\noindent
\textbf{Drei Strategien für Bedingung~(D)} sind in
Bemerkung~\ref{rem:diss-gap} skizziert: Heben des Funktionals auf $H^1$-Ebene,
Ausnutzung quantitativer Blow-up-Raten oder Verwendung Serrin-artiger Kriterien.
Schließung entweder der Projektionslücke (Assumption~G) \emph{oder} der
Dissipationslücke (Bedingung~D) würde das vorliegende Rahmenwerk voranbringen,
und Schließung beider würde einen unbedingten Regularitätsbeweis liefern.

\medskip\noindent
\textbf{Das $H^1$-Ebene-Programm} (Abschnitt~\ref{sec:H1-lift})
bietet den vielversprechendsten Weg nach vorn: Durch Heben des
Attraktorabstandsfunktionals von $L^2$ auf~$H^1$ wird der
Dissipationsterm zu $\int\|\Delta u\|^2$, der -- anders als
$\int\|\nabla u\|^2$ -- \emph{keine} a-priori obere Schranke
unter Blow-up hat. Diese strukturelle Asymmetrie zwischen Energie- und
Enstrophiedissipation ist die zentrale Erkenntnis der vorliegenden Arbeit
und motiviert Satz~\ref{thm:H1-main} als die natürliche
Fortsetzung dieses Programms.
Auf der $H^1$-Ebene tritt eine weitere Subtilität auf: Der
Gronwall-Koeffizient $C_1^{(1)}$ hängt von $\|u\|_{H^1}^2$ ab,
was eine Bihari-artige (nichtlineare Gronwall) Differentialungleichung
ergibt, deren integrierte Kosten wachsen, wenn sich die Lösung dem
hypothetischen Blow-up nähert. Der Widerspruchsbeweis bleibt gültig,
sofern die divergente Enstrophiedissipation dieses Wachstum übertrifft --
eine quantitative Verfeinerung, die das $H^1$-Programm mit präzisen
Blow-up-Raten-Abschätzungen verbindet.

% ===================================================================
% AI Disclosure
% ===================================================================
\subsection*{KI-Offenlegung}
Bei der Erstellung dieses Manuskripts wurden KI-gestützte Werkzeuge (Claude, Anthropic)
in unterstützender Funktion eingesetzt, insbesondere für Literaturrecherche,
Textstrukturierung und Formulierungshilfe. Der konzeptionelle Entwurf,
die wissenschaftliche Argumentation und die Verantwortung für den gesamten Inhalt
liegen ausschließlich beim Autor.

\begin{thebibliography}{99}

\bibitem{BealeKatoMajda1984}
J.~T. Beale, T.~Kato, and A.~Majda,
\emph{Remarks on the breakdown of smooth solutions for the 3-D Euler equations},
Comm.\ Math.\ Phys.\ \textbf{94} (1984), no.~1, 61--66.

\bibitem{BertozziMajda2002}
A.~L. Bertozzi and A.~J. Majda,
\emph{Vorticity and Incompressible Flow},
Cambridge Texts in Applied Mathematics, Cambridge University Press, 2002.

\bibitem{CaffarelliKohnNirenberg1982}
L.~Caffarelli, R.~Kohn, and L.~Nirenberg,
\emph{Partial regularity of suitable weak solutions of the Navier--Stokes equations},
Comm.\ Pure Appl.\ Math.\ \textbf{35} (1982), no.~6, 771--831.

\bibitem{ConstantinFoias1988}
P.~Constantin and C.~Foias,
\emph{Navier--Stokes Equations},
Chicago Lectures in Mathematics, University of Chicago Press, 1988.

\bibitem{FoiasManleyRosaTemam2001}
C.~Foias, O.~Manley, R.~Rosa, and R.~Temam,
\emph{Navier--Stokes Equations and Turbulence},
Encyclopedia of Mathematics and its Applications, vol.~83, Cambridge University Press, 2001.

\bibitem{LasryLions2007}
J.-M. Lasry and P.-L. Lions,
\emph{Mean field games},
Jpn.\ J.\ Math.\ \textbf{2} (2007), no.~1, 229--260.

\bibitem{Leray1934}
J.~Leray,
\emph{Sur le mouvement d'un liquide visqueux emplissant l'espace},
Acta Math.\ \textbf{63} (1934), 193--248.

\bibitem{Temam2001}
R.~Temam,
\emph{Navier--Stokes Equations: Theory and Numerical Analysis},
AMS Chelsea Publishing, Providence, RI, 2001 (reprint of the 1984 edition).

\bibitem{Wagner1977}
D.~H. Wagner,
\emph{Survey of measurable selection theorems},
SIAM J.\ Control Optim.\ \textbf{15} (1977), no.~5, 859--903.

\bibitem{KuksinShirikyan2012}
S.~Kuksin and A.~Shirikyan,
\emph{Mathematics of Two-Dimensional Turbulence},
Cambridge Tracts in Mathematics, vol.~194, Cambridge University Press, 2012.

\bibitem{HairerMattingly2006}
M.~Hairer and J.~C. Mattingly,
\emph{Ergodicity of the 2D Navier--Stokes equations with degenerate stochastic forcing},
Ann.\ of Math.\ \textbf{164} (2006), no.~3, 993--1032.

\bibitem{MalletParetSell1988}
J.~Mallet-Paret and G.~R. Sell,
\emph{Inertial manifolds for reaction diffusion equations in higher space dimensions},
J.~Amer.\ Math.\ Soc.\ \textbf{1} (1988), no.~4, 805--866.

\bibitem{AmbrosioFuscoPallara2000}
L.~Ambrosio, N.~Fusco, and D.~Pallara,
\emph{Functions of Bounded Variation and Free Discontinuity Problems},
Oxford Mathematical Monographs, Clarendon Press, Oxford, 2000.

\bibitem{BaeCannone2016}
H.~Bae and M.~Cannone,
\emph{Log-Lipschitz regularity of the 3D Navier--Stokes equations},
Nonlinear Anal.\ \textbf{135} (2016), 223--235.
\href{https://doi.org/10.1016/j.na.2016.01.014}{doi:10.1016/j.na.2016.01.014}.

\bibitem{IlyinKalantarovZelik2025}
A.~Ilyin, V.~Kalantarov and S.~Zelik,
\emph{Attractors and their dimensions for the 3D fractional
Navier--Stokes--Voigt equations},
arXiv:2511.04783, 2025.

\bibitem{MiranvilleZelik2008}
A.~Miranville and S.~Zelik,
\emph{Attractors for dissipative partial differential equations in bounded and unbounded domains},
in: Handbook of Differential Equations: Evolutionary Equations, vol.~4 (C.~M.~Dafermos and M.~Pokorn\'y, eds.), Elsevier, 2008, pp.~103--200.
\href{https://doi.org/10.1016/S1874-5717(08)00003-0}{doi:10.1016/S1874-5717(08)00003-0}.

\bibitem{Zelik2014}
S.~Zelik,
\emph{Inertial manifolds and finite-dimensional reduction for
dissipative PDEs},
Proc.\ Roy.\ Soc.\ Edinburgh Sect.~A \textbf{144} (2014), no.~6,
1245--1327.

\bibitem{EdenFoiasNicolaenko1994}
A.~Eden, C.~Foias, B.~Nicolaenko, and R.~Temam,
\emph{Exponential Attractors for Dissipative Evolution Equations},
Wiley, 1994.

\bibitem{FST-TU2026}
L.~Geiger, \emph{Anomalous Dissipation and the Turbulent Cascade:
A Variational Free-Energy Characterisation of Kolmogorov Scaling},
Zenodo, 2026.
\href{https://doi.org/10.5281/zenodo.19056813}{doi:10.5281/zenodo.19056813}.

\bibitem{GeigerRH2026c}
L.~Geiger, \emph{From Landscape to Proof: The Even-Dominance Route to the Riemann Hypothesis},
Zenodo, 2026.
\href{https://doi.org/10.5281/zenodo.19035640}{doi:10.5281/zenodo.19035640}.

\bibitem{FST-Unified2026}
L.~Geiger, \emph{FST Spectrum Duality: The Renormalized Free Energy Principle as Physical Instantiation of Functional Stability Theory},
Zenodo, 2026.
\href{https://doi.org/10.5281/zenodo.19036190}{doi:10.5281/zenodo.19036190}.

\bibitem{GrujicXu2025}
Z.~Grujic und L.~Xu,
\emph{Asymptotic criticality of the Navier--Stokes regularity problem},
J.~Math.\ Fluid Mech.\ \textbf{26} (2024), Article~53.
\href{https://doi.org/10.1007/s00021-024-00888-x}{doi:10.1007/s00021-024-00888-x}.

\bibitem{FoiasSellTemam88}
C.~Foias, G.\,R.~Sell und R.~Temam,
\emph{Inertial manifolds for nonlinear evolutionary equations},
J.~Differential Equations \textbf{73} (1988), 309--353.

\bibitem{Moffatt69}
H.\,K.~Moffatt,
\emph{The degree of knottedness of tangled vortex lines},
J.~Fluid Mech.\ \textbf{35} (1969), 117--129.

\bibitem{ArnoldKhesin98}
V.\,I.~Arnold und B.\,A.~Khesin,
\emph{Topological Methods in Hydrodynamics},
Applied Mathematical Sciences, Bd.~125,
Springer, New York, 1998.

\bibitem{DrivasNguyen2019}
T.\,D.~Drivas und H.\,Q.~Nguyen,
\emph{Remarks on the emergence of weak Euler solutions in the vanishing
viscosity limit},
J.~Nonlinear Sci.\ \textbf{29} (2019), 1127--1161.
\href{https://arxiv.org/abs/1808.01014}{arXiv:1808.01014}.

\bibitem{DeRosaDrivasInversi2024}
L.~De~Rosa, T.\,D.~Drivas und M.~Inversi,
\emph{On the support of anomalous dissipation measures},
J.~Math.\ Fluid Mech.\ \textbf{26} (2024), Article~56.
\href{https://doi.org/10.1007/s00021-024-00894-z}{doi:10.1007/s00021-024-00894-z}.
\href{https://arxiv.org/abs/2301.09603}{arXiv:2301.09603}.

\bibitem{DeRosaDrivasInversiIsett2025}
L.~De~Rosa, T.\,D.~Drivas, M.~Inversi und P.~Isett,
\emph{Intermittency and dissipation regularity in turbulence},
Preprint, 2025.
\href{https://arxiv.org/abs/2502.10032}{arXiv:2502.10032}.

\bibitem{DuchonRobert2000}
J.~Duchon und R.~Robert,
\emph{Inertial energy dissipation for weak solutions of incompressible Euler and Navier--Stokes equations},
Nonlinearity \textbf{13} (2000), Nr.~1, 249--255.

\bibitem{Sturm2025}
K.-T.~Sturm,
\emph{Bakry--\'Emery, Hardy, and spectral gap estimates on manifolds with conical singularities},
Calc.\ Var.\ Partial Differential Equations \textbf{64} (2025), Article~94.
\href{https://doi.org/10.1007/s00526-025-02946-2}{doi:10.1007/s00526-025-02946-2}.
\href{https://arxiv.org/abs/2405.10734}{arXiv:2405.10734}.

\bibitem{DondlFrenzelMielke2019}
P.~Dondl, T.~Frenzel und A.~Mielke,
\emph{A gradient system with a wiggly energy and relaxed EDP-convergence},
ESAIM Control Optim.\ Calc.\ Var.\ \textbf{25} (2019), Art.~68.
\href{https://doi.org/10.1051/cocv/2018058}{doi:10.1051/cocv/2018058}.

\bibitem{MielkeRossiSavare2016}
A.~Mielke, R.~Rossi und G.~Savar\'e,
\emph{Balanced-viscosity (BV) solutions to infinite-dimensional rate-independent systems},
J.~Eur.\ Math.\ Soc.\ \textbf{18} (2016), Nr.~9, 2107--2165.
\href{https://doi.org/10.4171/JEMS/639}{doi:10.4171/JEMS/639}.

\bibitem{KevrekidisTiti2024}
E.~D.~Koronaki, N.~Evangelou, C.~P.~Martin-Linares, E.~S.~Titi und I.~G.~Kevrekidis,
\emph{Nonlinear dimensionality reduction then and now: AIMs for dissipative PDEs in the ML era},
J.\ Comput.\ Phys.\ \textbf{506} (2024), 112910.
\href{https://doi.org/10.1016/j.jcp.2024.112910}{doi:10.1016/j.jcp.2024.112910}.

\bibitem{Geiger2026-DE}
L.~Geiger,
\emph{Dark Energy as Residual Vacuum Free Energy: A Thermodynamic Bound on the Cosmological Constant},
Zenodo, 2026.
\href{https://doi.org/10.5281/zenodo.19036235}{doi:10.5281/zenodo.19036235}.

\bibitem{Geiger2026-YM}
L.~Geiger,
\emph{Volume-Independent Spectral Gap for Strong-Coupling Lattice Gauge Theory via Poincar\'e--Dobrushin Machinery},
Zenodo, 2026.
\href{https://doi.org/10.5281/zenodo.19087433}{doi:10.5281/zenodo.19087433}.

\bibitem{Geiger2026-BSD}
L.~Geiger,
\emph{The BSD Conjecture as a Positivity Normal-Form Theorem},
Zenodo, 2026.
\href{https://doi.org/10.5281/zenodo.19087443}{doi:10.5281/zenodo.19087443}.

\bibitem{ChepyzhovVishik2002}
V.~V.~Chepyzhov and M.~I.~Vishik, \textit{Attractors for Equations of Mathematical Physics}, Amer.\ Math.\ Soc., Providence, 2002.

\bibitem{Simon1987}
J.~Simon, \textit{Compact sets in the space $L^p(0,T;B)$}, Ann.\ Mat.\ Pura Appl., \textbf{146} (1987), 65--96.

\bibitem{Artstein1995}
Z.~Artstein, \textit{A calculus for set-valued maps and set-valued evolution equations}, Set-Valued Anal., \textbf{3} (1995), 213--261.

\bibitem{PoliquinRockafellarThibault2000}
R.~A.~Poliquin, R.~T.~Rockafellar, and L.~Thibault, \textit{Local differentiability of distance functions}, Trans.\ Amer.\ Math.\ Soc., \textbf{352} (2000), 5231--5249.

\bibitem{Geiger2026-LDI} L.~Geiger, \textit{Log-Distance Integrability for Dissipative Flows: BV Selections on Fractal Attractors}, Zenodo, 2026.
\href{https://doi.org/10.5281/zenodo.19056807}{doi:10.5281/zenodo.19056807}.

\end{thebibliography}

\end{document}
